% Momentum as a Vector — Pure Mathematics with Mechanics Upper Sixth — Cours signé OnBuch+
% Official MINESEC syllabus. Compiler avec : tectonic PMM18-momentum-as-a-vector.tex
\newcommand{\DOCMATIERE}{Pure Mathematics with Mechanics}
\newcommand{\DOCNIVEAU}{Upper Sixth}
\newcommand{\DOCMODULE}{Upper Sixth — Pure mathematics with mechanics}
\newcommand{\DOCLECON}{18}
\newcommand{\DOCTITRE}{Momentum as a Vector}
% OnBuch+ — préambule « cours signés » ENRICHI (compilé avec tectonic / XeLaTeX)
% Système visuel « Pop » d'OnBuch+ : fond crème (jamais blanc), bordures encre
% épaisses, ombres « dures » décalées, titres Archivo Black en capitales.
% Version MATHÉMATIQUES : graphes (pgfplots/tikz), boîtes pédagogiques
% (définition, propriété, méthode, attention, le savais-tu, exercice résolu,
% exercice, TP, bilan), page de garde et signature OnBuch+ ancrée en bas.
%
% Macros attendues dans main.tex AVANT \input{preamble} :
%   \DOCMATIERE  \DOCNIVEAU  \DOCMODULE  \DOCLECON (numéro)  \DOCTITRE
\documentclass[11pt]{article}

\usepackage{fontspec}
\usepackage[english]{babel}
\usepackage[a4paper,margin=2cm,top=3.4cm,bottom=2.8cm,headheight=1.4cm,footskip=1.3cm]{geometry}
\usepackage{amsmath,amssymb,amsfonts}
\usepackage{mathtools} % \xmapsto, \coloneqq... souvent utilisés par les modèles
\usepackage{cancel} % simplifications barrées (bilans de forces)
% \abs{z} (module, valeur absolue) et \norm{u} (norme), utilisés par les modèles
\providecommand{\abs}{}\let\abs\relax\DeclarePairedDelimiter{\abs}{\lvert}{\rvert}
\providecommand{\norm}{}\let\norm\relax\DeclarePairedDelimiter{\norm}{\lVert}{\rVert}
\usepackage[table]{xcolor} % [table] : fournit aussi \rowcolors (alternance de lignes)
\usepackage{pagecolor}
\usepackage{tikz}
\usetikzlibrary{arrows.meta,positioning,calc,shapes.geometric,shapes.misc,shapes.arrows,decorations.pathmorphing,decorations.pathreplacing,decorations.markings,patterns,shadows,mindmap,trees,fit,backgrounds,matrix,intersections,angles,quotes}
\usepackage{pgfplots}
\usepackage[version=4]{mhchem} % équations nucléaires \ce{^{235}_{92}U}
\usepackage[european,siunitx]{circuitikz} % schémas électriques
\pgfplotsset{compat=1.18}
\usepgfplotslibrary{fillbetween,groupplots} % aires entre courbes (calcul intégral)
\usepackage{enumitem}
\usepackage{fancyhdr}
% [most] charge toutes les bibliothèques tcolorbox (listings, minted, raster,
% documentation...) et gonfle inutilement la mémoire TeX sur un document long
% (~300 boîtes) : on ne charge que ce qui est réellement utilisé.
\usepackage[skins,breakable]{tcolorbox}
\usepackage{graphicx}
\usepackage{colortbl}
\usepackage{booktabs}
\usepackage{tabularx}
\usepackage{diagbox} % case d'en-tête diagonale des tableaux à double entrée
% Colonnes extensibles centrée / gauche / droite (C, L, R), souvent utilisées par les modèles
\newcolumntype{C}{>{\centering\arraybackslash}X}
\newcolumntype{L}{>{\raggedright\arraybackslash}X}
\newcolumntype{R}{>{\raggedleft\arraybackslash}X}
\usepackage{array}
\usepackage{multicol}
\usepackage{siunitx}
% parse-numbers=false : accepte aussi \SI{2,5 \times 10^{-3}}{...}
\sisetup{locale=UK,output-decimal-marker={.},group-minimum-digits=4,parse-numbers=false}
\DeclareSIUnit\ohm{\ensuremath{\symup{\Omega}}} % U+2126 absent de la police
\DeclareSIUnit\division{div} % réglages d'oscilloscope (ms/div, V/div)
\DeclareSIUnit\tr{tr} % tours (vitesse de rotation, ex. \SI{1500}{\tr\per\minute})
\usepackage{qrcode}
% \degree : commande fréquemment utilisée par les modèles pour un angle ou une
% température en dehors de \SI{}{} (non fournie par les paquets chargés).
\providecommand{\degree}{^{\circ}}
% \qed (fin de démonstration), utilisé par les modèles sans amsthm
\providecommand{\qed}{\hfill$\square$}
% \barre : double barre des valeurs interdites dans un tableau de variations (tkz-tab)
\providecommand{\barre}{\ensuremath{\|}}
% environnement proof (amsthm non chargé) utilisé par les modèles
\ifdefined\proof\else\newenvironment{proof}[1][Proof]{\par\noindent\textit{#1.}\ }{\qed\par}\fi
\usepackage{lastpage}
\usepackage{hyperref}
\hypersetup{hidelinks}

% ============================================================
% Palette « Pop » OnBuch+ — mêmes hex que la classe P (plus_ui.dart)
% ============================================================
\definecolor{poporange}{HTML}{F59321}
\definecolor{popdark}{HTML}{E07A0C}
\definecolor{popcream}{HTML}{FFF6E8}
\definecolor{popink}{HTML}{1C1714}
\definecolor{poppurple}{HTML}{7A5AE0}
\definecolor{popgreen}{HTML}{2E9E6A}
\definecolor{poppink}{HTML}{E0457B}
\definecolor{popblue}{HTML}{2F7DE1}
\definecolor{popgold}{HTML}{C98A12}
\definecolor{popmuted}{HTML}{8A7F76}
\definecolor{popline}{HTML}{EADFD2}
\definecolor{popgray}{HTML}{8A7F76}
\colorlet{popgrey}{popgray}
% Alias tolérants (le modèle nomme parfois les couleurs différemment)
\colorlet{popwhite}{white}
\colorlet{popblack}{popink}
\colorlet{popred}{poppink}
\colorlet{popyellow}{popgold}
% Teintes claires pour fonds de tableaux / zones
\colorlet{popblueL}{popblue!10!white}
\colorlet{poporangeL}{poporange!14!white}
\colorlet{popgreenL}{popgreen!12!white}
\colorlet{poppurpleL}{poppurple!12!white}
\colorlet{poppinkL}{poppink!12!white}
\colorlet{popgoldL}{popgold!14!white}

\pagecolor{popcream}

% ============================================================
% Typographie
% ============================================================
\defaultfontfeatures{Ligatures=TeX}
\setmainfont{PlusJakartaSans-Medium.ttf}[Path=../../../fonts/,BoldFont=PlusJakartaSans-Bold.ttf]
% Maths en sans-serif (Fira Math) avec chiffres et lettres droites de Plus
% Jakarta, pour que les formules se fondent dans le texte.
\usepackage[math-style=ISO,bold-style=ISO]{unicode-math}
\setmathfont{FiraMath-Regular.otf}[Path=../../../fonts/]
\setmathfont{PlusJakartaSans-Medium.ttf}[Path=../../../fonts/,range={up/num,up/{Latin,latin},\mathup}]
% (icomma retiré : décimales avec point en anglais)
% Caractères mathématiques Unicode absents des polices de texte (Plus Jakarta,
% Archivo Black) : rendus par leur équivalent mathématique, en texte comme en
% formule — sinon ils s'affichent en carré vide (ex. « Arithmétique dans ℤ »).
\usepackage{newunicodechar}
\newunicodechar{ℝ}{\ensuremath{\mathbb{R}}}
\newunicodechar{ℤ}{\ensuremath{\mathbb{Z}}}
\newunicodechar{ℂ}{\ensuremath{\mathbb{C}}}
\newunicodechar{ℕ}{\ensuremath{\mathbb{N}}}
\newunicodechar{ℚ}{\ensuremath{\mathbb{Q}}}
\newunicodechar{𝔻}{\ensuremath{\mathbb{D}}}
\newunicodechar{α}{\ensuremath{\alpha}}
\newunicodechar{β}{\ensuremath{\beta}}
\newunicodechar{γ}{\ensuremath{\gamma}}
\newunicodechar{δ}{\ensuremath{\delta}}
\newunicodechar{ε}{\ensuremath{\varepsilon}}
\newunicodechar{θ}{\ensuremath{\theta}}
\newunicodechar{λ}{\ensuremath{\lambda}}
\newunicodechar{μ}{\ensuremath{\mu}}
\newunicodechar{σ}{\ensuremath{\sigma}}
\newunicodechar{τ}{\ensuremath{\tau}}
\newunicodechar{φ}{\ensuremath{\varphi}}
\newunicodechar{ψ}{\ensuremath{\psi}}
\newunicodechar{ω}{\ensuremath{\omega}}
\newunicodechar{Σ}{\ensuremath{\Sigma}}
% majuscules produites par \MakeUppercase dans les titres (α-aminés -> Α)
\newunicodechar{Α}{\ensuremath{\alpha}}
\newunicodechar{Β}{\ensuremath{\beta}}
\newunicodechar{Γ}{\ensuremath{\Gamma}}
\newunicodechar{Δ}{\ensuremath{\Delta}}
\newunicodechar{Ε}{\ensuremath{\varepsilon}}
\newunicodechar{Θ}{\ensuremath{\Theta}}
\newunicodechar{Λ}{\ensuremath{\Lambda}}
\newunicodechar{Μ}{\ensuremath{\mu}}
\newunicodechar{Τ}{\ensuremath{\tau}}
\newunicodechar{Φ}{\ensuremath{\Phi}}
\newunicodechar{Ψ}{\ensuremath{\Psi}}
\newunicodechar{Ω}{\ensuremath{\Omega}}
\newunicodechar{↦}{\ensuremath{\mapsto}}
\newunicodechar{∈}{\ensuremath{\in}}
\newunicodechar{∉}{\ensuremath{\notin}}
\newunicodechar{∧}{\ensuremath{\wedge}}
\newunicodechar{⇔}{\ensuremath{\Leftrightarrow}}
\newunicodechar{⟺}{\ensuremath{\Longleftrightarrow}}
\newunicodechar{⇒}{\ensuremath{\Rightarrow}}
\newunicodechar{⟹}{\ensuremath{\Longrightarrow}}
\newunicodechar{∘}{\ensuremath{\circ}}
\newunicodechar{∪}{\ensuremath{\cup}}
\newunicodechar{∩}{\ensuremath{\cap}}
\newunicodechar{≡}{\ensuremath{\equiv}}
\newunicodechar{⊂}{\ensuremath{\subset}}
\newunicodechar{⊥}{\ensuremath{\perp}}
\newunicodechar{∃}{\ensuremath{\exists}}
\newunicodechar{∀}{\ensuremath{\forall}}
\newunicodechar{∛}{\ensuremath{\sqrt[3]{\,}}}
\newunicodechar{∜}{\ensuremath{\sqrt[4]{\,}}}
\newunicodechar{⃗}{\ensuremath{{}^{\to}}}
\newunicodechar{ⁿ}{\ifmmode^{n}\else\textsuperscript{\ensuremath{n}}\fi}
\newunicodechar{ˣ}{\ifmmode^{x}\else\textsuperscript{\ensuremath{x}}\fi}
\newunicodechar{ᵇ}{\ifmmode^{b}\else\textsuperscript{\ensuremath{b}}\fi}
\newunicodechar{ʳ}{\ifmmode^{r}\else\textsuperscript{\ensuremath{r}}\fi}
\newunicodechar{ᵘ}{\ifmmode^{u}\else\textsuperscript{\ensuremath{u}}\fi}
\newunicodechar{ʸ}{\ifmmode^{y}\else\textsuperscript{\ensuremath{y}}\fi}
\newunicodechar{ᵃ}{\ifmmode^{a}\else\textsuperscript{\ensuremath{a}}\fi}
\newunicodechar{ᵉ}{\ifmmode^{e}\else\textsuperscript{\ensuremath{e}}\fi}
\newunicodechar{ᵏ}{\ifmmode^{k}\else\textsuperscript{\ensuremath{k}}\fi}
\newunicodechar{ᵖ}{\ifmmode^{p}\else\textsuperscript{\ensuremath{p}}\fi}
\newunicodechar{ᴺ}{\ifmmode^{N}\else\textsuperscript{\ensuremath{N}}\fi}
\newunicodechar{ᶜ}{\ifmmode^{c}\else\textsuperscript{\ensuremath{c}}\fi}
\newunicodechar{ʰ}{\ifmmode^{h}\else\textsuperscript{\ensuremath{h}}\fi}
\newunicodechar{ᵠ}{\ifmmode^{q}\else\textsuperscript{\ensuremath{q}}\fi}
\newunicodechar{ₙ}{\ifmmode_{n}\else\textsubscript{\ensuremath{n}}\fi}
\newunicodechar{ₐ}{\ifmmode_{a}\else\textsubscript{\ensuremath{a}}\fi}
\newunicodechar{ᵢ}{\ifmmode_{i}\else\textsubscript{\ensuremath{i}}\fi}
\newunicodechar{ᵣ}{\ifmmode_{r}\else\textsubscript{\ensuremath{r}}\fi}
\newunicodechar{ₖ}{\ifmmode_{k}\else\textsubscript{\ensuremath{k}}\fi}
\newunicodechar{ᵦ}{\ifmmode_{\beta}\else\textsubscript{\ensuremath{\beta}}\fi}
\newunicodechar{ₓ}{\ifmmode_{x}\else\textsubscript{\ensuremath{x}}\fi}
\newfontfamily\popdisplay{ArchivoBlack-Regular.ttf}[Path=../../../fonts/]
\newfontfamily\poplogo{SpaceGrotesk-Bold.ttf}[Path=../../../fonts/]
\newfontfamily\popsemi{PlusJakartaSans-SemiBold.ttf}[Path=../../../fonts/]
\newfontfamily\popxbold{PlusJakartaSans-ExtraBold.ttf}[Path=../../../fonts/]
\newfontfamily\popxboldls{PlusJakartaSans-ExtraBold.ttf}[Path=../../../fonts/,LetterSpace=3.0]

\setlist[enumerate]{leftmargin=1.7em,itemsep=5pt,topsep=4pt}
\setlist[itemize]{leftmargin=1.7em,itemsep=4pt,topsep=4pt,label={\color{poporange}\textbullet}}
\setlength{\parindent}{0pt}
\setlength{\parskip}{5pt}
\renewcommand{\arraystretch}{1.25}

% pgfplots : style Pop par défaut
\pgfplotsset{
  every axis/.append style={axis line style={popink,line width=1pt},tick style={popink},
    grid style={popline,line width=0.5pt},label style={font=\small},tick label style={font=\footnotesize}},
  popcurve/.style={poporange,line width=1.8pt,no markers,smooth},
}
% Même style disponible dans un \draw TikZ simple (hors pgfplots)
\tikzset{popcurve/.style={poporange,line width=1.8pt}}
\tikzset{popgrid/.style={popline,very thin}}
\colorlet{popcurve}{poporange} % le nom du style est parfois utilisé comme couleur

% ============================================================
% Pill / squiggle
% ============================================================
\newcommand{\poppill}[3]{%
  \tikz[baseline=-0.55ex]\node[draw=popink,line width=1.2pt,rounded corners=2.6pt,%
    fill=#2,inner xsep=6.5pt,inner ysep=3pt]{%
    \popxboldls\fontsize{7}{7}\selectfont\color{#3}\MakeUppercase{#1}};%
}
\newcommand{\popsquiggle}[1]{%
  \tikz[baseline=-0.4ex]{
    \draw[#1,line width=2.6pt,line cap=round]
      plot [smooth] coordinates {(0,0.09) (0.34,0.01) (0.68,0.21) (1.02,0.01) (1.36,0.10)};
  }%
}
% Mot-clé surligné (marqueur orange sous le texte)
\newcommand{\cle}[1]{{\popxbold\color{popdark}#1}}

% ============================================================
% En-tête / pied
% ============================================================
\pagestyle{fancy}
\fancyhf{}
\renewcommand{\headrulewidth}{0pt}
\renewcommand{\footrulewidth}{0pt}
\lhead{\raisebox{-0.15ex}{\poplogo\fontsize{13}{13}\selectfont\color{popink}OnBuch\color{poporange}+}}
\rhead{\poppill{\DOCMATIERE\ \textbullet\ \DOCNIVEAU\ \textbullet\ Chapter \DOCLECON}{white}{popink}}
\lfoot{\popxbold\fontsize{7.6}{7.6}\selectfont\color{popmuted}\MakeUppercase{Signed course OnBuch+}\ \textbullet\ {\popsemi\DOCTITRE}}
\rfoot{\tikz[baseline=-0.6ex]\node[rounded corners=8.5pt,fill=popink,minimum height=17pt,minimum width=17pt,inner xsep=5pt,inner sep=0]{\popxbold\fontsize{8}{8}\selectfont\color{popcream}\thepage\,/\,\pageref*{LastPage}};}

% ============================================================
% Titres
% ============================================================
\newcommand{\chaptitre}[2]{%
  \begin{tcolorbox}[enhanced,colback=poporange,colframe=popink,boxrule=2.6pt,arc=11pt,
      drop shadow=popink,left=15pt,right=15pt,top=12pt,bottom=11pt,before skip=3pt,after skip=18pt]
    \poppill{#2}{popink}{popcream}\par\vspace{9pt}
    {\raggedright\hyphenpenalty=10000\exhyphenpenalty=10000\popdisplay\fontsize{22}{24}\selectfont\color{popcream}\MakeUppercase{#1}\par}
  \end{tcolorbox}
}
% Section numérotée : pastille orange avec le numéro + titre Archivo
\newcounter{popsec}
\newcommand{\coursec}[1]{%
  \stepcounter{popsec}\setcounter{popsub}{0}%
  \par\vspace{16pt}\noindent
  \begin{minipage}{\linewidth}
  \tikz[baseline=-0.7ex]\node[draw=popink,line width=1.6pt,fill=poporange,rounded corners=4pt,minimum size=22pt,inner sep=0]{\popdisplay\fontsize{12}{12}\selectfont\color{popcream}\thepopsec};%
  \hspace{8pt}{\popdisplay\fontsize{15}{17}\selectfont\color{popink}\MakeUppercase{#1}}\par
  \vspace{4pt}\noindent{\color{poporange}\rule{36pt}{3.6pt}}
  \end{minipage}\par\nopagebreak\vspace{8pt}
}
\newcounter{popsub}
\newcommand{\courssub}[1]{%
  \stepcounter{popsub}%
  \par\vspace{9pt}\noindent{\popxbold\fontsize{11.5}{13}\selectfont\color{popink}{\color{poporange}\thepopsec.\thepopsub}\hspace{6pt}#1}\par\nopagebreak\vspace{4pt}
}

% ============================================================
% Boîtes pédagogiques — même recette : fond blanc, bordure épaisse colorée,
% ombre dure encre, badge plein. Titre optionnel [#1].
% ============================================================
\tcbset{popbase/.style={breakable,enhanced,colback=white,boxrule=2.3pt,arc=9pt,
  shadow={3pt}{-3pt}{0pt}{popink},left=14pt,right=14pt,top=11pt,bottom=11pt,before skip=11pt,after skip=15pt,
  fontupper=\popsemi\fontsize{10.6}{14.5}\selectfont\color{popink},
  fontlower=\popsemi\fontsize{10.6}{14.5}\selectfont\color{popink}}}
% Coche dessinée (le glyphe ✓ n'existe pas dans les polices du document)
\newcommand{\popcheck}[1]{\tikz[baseline=-0.5ex]\draw[#1,line width=1.6pt,line cap=round,line join=round](-0.13,0.0)--(-0.04,-0.09)--(0.13,0.1);}
\providecommand{\coche}{\popcheck{popgreen}}
\providecommand{\resultat}[1]{\par\smallskip\noindent\fcolorbox{popgreen}{popgreenL}{\parbox{\dimexpr\linewidth-2\fboxsep-2\fboxrule\relax}{\textbf{Result.} #1}}\par\smallskip}
\newcommand{\popboxhead}[3]{\poppill{#1}{#2}{white}\ifx\relax#3\relax\else\hspace{6pt}{\popxbold\fontsize{10.6}{12}\selectfont\color{popink}#3}\fi\par\vspace{7pt}}

\newtcolorbox{aretenir}[1][]{popbase,colframe=popgreen,before upper={\popboxhead{Key points}{popgreen}{#1}}}
\newtcolorbox{exemplebox}[1][]{popbase,colframe=poporange,before upper={\popboxhead{Exemple}{poporange}{#1}}}
\newtcolorbox{definition}[1][]{popbase,colframe=popblue,before upper={\popboxhead{Definition}{popblue}{#1}}}
\newtcolorbox{propriete}[1][]{popbase,colframe=poppurple,before upper={\popboxhead{Property}{poppurple}{#1}}}
\newtcolorbox{methode}[1][]{popbase,colframe=popgold,before upper={\popboxhead{Method}{popgold}{#1}}}
\newtcolorbox{attention}[1][]{popbase,colframe=poppink,before upper={\popboxhead{Warning}{poppink}{#1}}}
\newtcolorbox{experience}[1][]{popbase,colframe=popblue,colback=popblueL,before upper={\popboxhead{Experiment}{popblue}{#1}}}
% « Le savais-tu ? » : carte violette pleine (contexte, culture, Cameroun)
\newtcolorbox{savaistu}[1][]{popbase,colback=poppurple,colframe=popink,
  fontupper=\popsemi\fontsize{10.4}{14.2}\selectfont\color{white},
  before upper={\poppill{Did you know?}{white}{poppurple}\ifx\relax#1\relax\else\hspace{6pt}{\popxbold\color{white}#1}\fi\par\vspace{7pt}}}
% Exercice résolu : énoncé en haut, \tcblower puis la solution sur fond crème
\newcounter{popexo}\newcounter{popexr}
\newtcolorbox{exoresolu}[1][]{popbase,colframe=poporange,colbacklower=popcream,
  segmentation style={popink,line width=1.2pt,dashed},
  before upper={\stepcounter{popexr}\popboxhead{Worked example \thepopexr}{poporange}{#1}},
  before lower={\poppill{Solution}{popink}{popcream}\par\vspace{6pt}}}
% Exercice (non résolu en place) — la correction est dans la partie Corrigés
\newtcolorbox{exercice}[1][]{popbase,colframe=popink,
  before upper={\stepcounter{popexo}\popboxhead{Exercise \thepopexo}{popink}{#1}}}
% Corrigé
\newtcolorbox{corrige}[1][]{popbase,colframe=popgreen,colback=popgreenL,
  before upper={\popboxhead{Solution}{popgreen}{#1}}}
% Objectifs / prérequis / bilan
\newtcolorbox{objectifs}{popbase,colframe=popink,colback=poporangeL,
  before upper={\popboxhead{Chapter objectives}{popink}{}}}
\newtcolorbox{prerequis}{popbase,colframe=popink,colback=popblueL,
  before upper={\popboxhead{Prerequisites}{popblue}{}}}
\newtcolorbox{bilan}{popbase,colframe=popink,colback=white,boxrule=2.8pt,
  before upper={\popboxhead{Fiche bilan}{poporange}{L'essentiel à connaître pour l'examen}}}
% Figure encadrée avec légende
\newtcolorbox{popfigure}[1][]{enhanced,breakable=false,colback=white,colframe=popline,boxrule=1.4pt,arc=8pt,
  left=10pt,right=10pt,top=10pt,bottom=8pt,before skip=10pt,after skip=12pt,halign=center,
  lower separated=false,halign lower=center,
  fontlower=\popsemi\fontsize{9}{11.5}\selectfont\color{popmuted},
  #1}
\newcommand{\legende}[1]{\par\vspace{6pt}{\popsemi\fontsize{9}{11.5}\selectfont\color{popmuted}#1}}

% ============================================================
% Signature OnBuch+
% ============================================================
\newcommand{\poplogoinline}[1]{{\poplogo\fontsize{#1}{#1}\selectfont\color{popink}OnBuch\color{poporange}+}}
\newcommand{\signatureonbuch}{%
  \begin{tcolorbox}[enhanced,colback=popink,colframe=popink,boxrule=2.2pt,arc=10pt,
      drop shadow=poporange,left=14pt,right=14pt,top=10pt,bottom=10pt,before skip=0pt,after skip=0pt]
    \tikz[baseline=-0.7ex]\node[circle,fill=popgreen,draw=popcream,line width=1.3pt,minimum size=22pt,inner sep=0]{\popcheck{white}};%
    \hspace{9pt}\begin{minipage}[c]{0.62\linewidth}
      {\popxbold\fontsize{12}{13}\selectfont\color{popcream}Signed\ \textbullet\ {\poplogo OnBuch\color{poporange}+}}\par\vspace{2pt}
      {\raggedright\popsemi\fontsize{8}{10.5}\selectfont\color{popline}\MakeUppercase{OnBuch+ teaching team}\\\MakeUppercase{Follows the official MINESEC syllabus}\par}
    \end{minipage}\hfill
    {\poplogo\fontsize{22}{22}\selectfont\color{popcream}OnBuch\color{poporange}+}
  \end{tcolorbox}%
}

% ============================================================
% Page de garde
% #1 titre · #2 matière · #3 niveau · #4 module · #5 n° leçon · #6 durée ·
% #7 description / objectifs (texte ou itemize)
% ============================================================
\newcommand{\pagegarde}[7]{%
  \thispagestyle{empty}
  \newgeometry{a4paper,margin=1.7cm,top=1.6cm,bottom=1.4cm}%
  \begin{tikzpicture}[remember picture,overlay]
    \fill[poporange!18] ([xshift=-3.2cm,yshift=-3.6cm]current page.north east) circle (4.6cm);
    \fill[poppurple!14] ([xshift=2.4cm,yshift=7.6cm]current page.south west) circle (3.8cm);
  \end{tikzpicture}%
  {\poplogo\fontsize{30}{30}\selectfont\color{popink}OnBuch\color{poporange}+}\hfill
  \raisebox{0.6ex}{\poppill{Signed course \textbullet\ Premium}{popink}{popcream}}\par
  \vspace{1pt}\popsquiggle{poppurple}\par
  \vspace{8pt}
  \begin{tcolorbox}[enhanced,colback=poporange,colframe=popink,boxrule=2.8pt,arc=13pt,
      drop shadow=popink,left=18pt,right=18pt,top=12pt,bottom=13pt,after skip=10pt]
    \poppill{#2 \textbullet\ #3}{popink}{popcream}\hspace{5pt}\poppill{#4}{white}{popink}\par\vspace{8pt}
    {\popdisplay\fontsize{14}{14}\selectfont\color{popink}CHAPTER #5}\par\vspace{4pt}
    {\raggedright\hyphenpenalty=10000\exhyphenpenalty=10000\popdisplay\fontsize{25}{27}\selectfont\color{popcream}\MakeUppercase{#1}\par}
    \vspace{8pt}
    {\popxbold\fontsize{9.5}{9.5}\selectfont\color{popink}\MakeUppercase{Suggested time: #6}}
  \end{tcolorbox}
  \begin{tcolorbox}[enhanced,colback=white,colframe=popink,boxrule=2.2pt,arc=10pt,
      drop shadow=popink,left=16pt,right=16pt,top=10pt,bottom=10pt,after skip=8pt]
    \poppill{In this chapter}{poppurple}{white}\par\vspace{5pt}
    \popsemi\fontsize{9.3}{12.6}\selectfont\color{popink}#7
  \end{tcolorbox}
  \noindent\begin{minipage}[t]{0.487\linewidth}
    \begin{tcolorbox}[enhanced,colback=popgreen,colframe=popink,boxrule=2.2pt,arc=10pt,
        drop shadow=popink,left=11pt,right=11pt,top=10pt,bottom=10pt,height=3.1cm,valign=center]
      \tcbox[colback=white,colframe=popink,boxrule=1.3pt,arc=5pt,boxsep=0pt,nobeforeafter,
        left=3pt,right=3pt,top=3pt,bottom=3pt,valign=center]{\qrcode[height=1.9cm]{https://play.google.com/store/apps/details?id=com.luvvix.onbuch&hl=en}}%
      \hspace{8pt}\begin{minipage}[c]{0.5\linewidth}
        {\popdisplay\fontsize{11}{12.5}\selectfont\color{white}\MakeUppercase{Download OnBuch}}\par\vspace{4pt}
        {\raggedright\popsemi\fontsize{8.4}{11}\selectfont\color{white}Free \textbullet\ Google Play\\AI tutor, past papers, results\par}
      \end{minipage}
    \end{tcolorbox}
  \end{minipage}\hfill
  \begin{minipage}[t]{0.487\linewidth}
    \begin{tcolorbox}[enhanced,colback=poppurple,colframe=popink,boxrule=2.2pt,arc=10pt,
        drop shadow=popink,left=11pt,right=11pt,top=10pt,bottom=10pt,height=3.1cm,valign=center]
      \tikz[baseline=-0.7ex]\node[circle,fill=white,draw=popink,line width=1.3pt,minimum size=1.6cm,inner sep=0]{\popdisplay\fontsize{24}{24}\selectfont\color{poppurple}+};%
      \hspace{8pt}\begin{minipage}[c]{0.6\linewidth}
        {\popdisplay\fontsize{11}{12.5}\selectfont\color{white}\MakeUppercase{Go OnBuch+}}\par\vspace{4pt}
        {\raggedright\popsemi\fontsize{8.4}{11}\selectfont\color{white}All signed courses, unlimited AI tutor, PDF sheets — OnBuch+ tab in the app\par}
      \end{minipage}
    \end{tcolorbox}
  \end{minipage}
  \vspace{8pt}
  \popcontenu
  \vfill
  \signatureonbuch
  \restoregeometry
  \newpage
}

% Rangée « Ce cours contient » (page de garde)
\newcommand{\poptile}[3]{% #1 couleur #2 gros texte #3 légende
  \begin{tcolorbox}[enhanced,colback=#1,colframe=popink,boxrule=2pt,arc=9pt,shadow={2.5pt}{-2.5pt}{0pt}{popink},
      left=6pt,right=6pt,top=9pt,bottom=9pt,halign=center,valign=center,height=1.75cm,width=\linewidth,nobeforeafter]
    {\popdisplay\fontsize{12.5}{13}\selectfont\color{white}\MakeUppercase{#2}}\par\vspace{3pt}
    {\centering\popsemi\fontsize{7.8}{9.5}\selectfont\color{white}#3\par}
  \end{tcolorbox}}
\newcommand{\popcontenu}{%
  \par\noindent{\popxboldls\fontsize{7.5}{7.5}\selectfont\color{popmuted}THIS SIGNED COURSE CONTAINS}\par\vspace{7pt}
  \noindent\begin{minipage}[t]{0.235\linewidth}\poptile{popblue}{Course}{complete, illustrated, step by step}\end{minipage}\hfill
  \begin{minipage}[t]{0.235\linewidth}\poptile{poporange}{Methods}{and worked examples}\end{minipage}\hfill
  \begin{minipage}[t]{0.235\linewidth}\poptile{poppink}{Exercises}{exam style + solutions}\end{minipage}\hfill
  \begin{minipage}[t]{0.235\linewidth}\poptile{popgreen}{Summary sheet}{the essentials to remember}\end{minipage}\par}

% Sommaire
\newtcolorbox{sommaire}{enhanced,colback=white,colframe=popink,boxrule=2.2pt,arc=10pt,
  drop shadow=popink,left=18pt,right=18pt,top=13pt,bottom=13pt,before skip=6pt,after skip=20pt,
  before upper={\poppill{Contents}{poppurple}{white}\par\vspace{9pt}\popsemi\fontsize{10.6}{14.5}\selectfont\color{popink}}}

% Signature de fin de cours — poussée tout en bas de la dernière page
\newcommand{\findecours}{%
  \par\vfill\nopagebreak
  \begin{center}\popsquiggle{poporange}\end{center}\vspace{4pt}
  \signatureonbuch
}
\AtBeginDocument{\def\llbracket{\mathopen{[\mkern-3.5mu[}}\def\rrbracket{\mathclose{]\mkern-3.5mu]}}} % intervalles d'entiers (glyphe ⟦ absent de la police)
% Glyphes absents de Fira Math : redéfinis à partir de glyphes disponibles
\AtBeginDocument{%
  \def\setminus{\mathbin{\backslash}}\let\smallsetminus\setminus
  \def\vdots{\mathord{\vbox{\baselineskip3.2pt\lineskiplimit0pt\kern4pt\hbox{.}\hbox{.}\hbox{.}}}}%
  \def\ddots{\mathinner{\mkern1mu\raise6.5pt\hbox{.}\mkern2mu\raise3.6pt\hbox{.}\mkern2mu\raise0.7pt\hbox{.}\mkern1mu}}%
  \def\triangle{\mathord{\scalebox{0.9}{$\symup{\Delta}$}}}%
  \def\nabla{\mathord{\raisebox{\depth}{\scalebox{1}[-1]{$\symup{\Delta}$}}}}%
  \def\checkmark{\popcheck{popdark}}%
  \def\star{\mathbin{\ast}}\def\bigstar{\mathord{\ast}}%
  \def\diamond{\mathbin{\scalebox{0.6}{$\Diamond$}}}%
  \def\bigcup{\mathop{\mathchoice{\raisebox{-0.3em}{\scalebox{1.7}{$\cup$}}}{\scalebox{1.25}{$\cup$}}{\cup}{\cup}}\displaylimits}%
  \def\bigcap{\mathop{\mathchoice{\raisebox{-0.3em}{\scalebox{1.7}{$\cap$}}}{\scalebox{1.25}{$\cap$}}{\cap}{\cap}}\displaylimits}%
  \def\lesseqqgtr{\mathrel{\lessgtr}}\def\bigoplus{\mathop{\oplus}\displaylimits}%
}
\providecommand{\ssi}{if and only if} % abréviation fréquente
\AtBeginDocument{\providecommand{\wideparen}{\overparen}} % arc de cercle

%% FIGFIX-BEGIN v3 — correctifs globaux des figures (docs/onbuch-generation/tools/figfix)
\usepackage{adjustbox}\usepackage{etoolbox}
% toute figure TikZ/pgfplots plus large que la ligne est réduite à la largeur disponible
\BeforeBeginEnvironment{tikzpicture}{\begin{adjustbox}{max width=\linewidth}}
\AfterEndEnvironment{tikzpicture}{\end{adjustbox}}
% légendes pgfplots lisibles par-dessus les courbes
\pgfplotsset{every axis/.append style={legend style={fill=white,fill opacity=0.92,text opacity=1,draw=black!15,inner sep=3pt}}}
% étiquettes d'arêtes (syntaxe edge["..."]) avec fond blanc
\tikzset{every edge quotes/.append style={fill=white,inner sep=1.2pt}}
% bulles de cartes mentales : texte à la ligne dans la bulle, bulle un peu plus grande
\tikzset{every concept/.append style={text width=2.0cm,align=center,minimum size=2.3cm,inner sep=2pt}}
%% FIGFIX-END
\begin{document}

\pagegarde{Momentum as a Vector}{Pure Mathematics with Mechanics}{Upper Sixth}{Upper Sixth — Pure mathematics with mechanics}{18}{12 periods}{This chapter introduces linear momentum as a vector quantity and uses it to give a deeper meaning to Newton's laws of motion. Through the principle of conservation of momentum, learners analyse collisions, explosions, impulse, rocket motion and the motion of the centre of mass in one and two dimensions, with direct applications to road safety and everyday life in Cameroon.
\vspace{4pt}
\begin{multicols}{2}\small
\begin{itemize}
\item By the end of this chapter I can define linear momentum, state its vector nature and its SI unit, and compute the momentum of one or several bodies.
\item By the end of this chapter I can relate force to the rate of change of momentum and use this as the general form of Newton's second law.
\item By the end of this chapter I can state and apply the principle of conservation of linear momentum to isolated systems of interacting bodies.
\item By the end of this chapter I can distinguish elastic from inelastic collisions and solve one-dimensional collision problems using conservation of momentum and, where appropriate, conservation of kinetic energy.
\item By the end of this chapter I can analyse two-dimensional collisions by resolving momentum into perpendicular components.
\item By the end of this chapter I can define impulse, use the impulse-momentum theorem, and interpret force-time graphs, including applications to safety devices.
\end{itemize}\end{multicols}}

\begin{sommaire}
\begin{multicols}{2}
\begin{enumerate}
\item Momentum, Force and Impulse
\item Conservation of Linear Momentum
\item Collisions in One Dimension: Elastic and Inelastic
\item Collisions and Momentum in Two Dimensions
\item Explosions, Rockets and Centre of Mass
\item Real-Life Applications and Synthesis of Momentum as a Vector
\item Integration activity
\item Methods and worked examples
\item Exercises
\item Solutions to the exercises
\item Summary sheet
\end{enumerate}
\end{multicols}
\end{sommaire}

\begin{objectifs}
\begin{itemize}
\item By the end of this chapter I can define linear momentum, state its vector nature and its SI unit, and compute the momentum of one or several bodies.
\item By the end of this chapter I can relate force to the rate of change of momentum and use this as the general form of Newton's second law.
\item By the end of this chapter I can state and apply the principle of conservation of linear momentum to isolated systems of interacting bodies.
\item By the end of this chapter I can distinguish elastic from inelastic collisions and solve one-dimensional collision problems using conservation of momentum and, where appropriate, conservation of kinetic energy.
\item By the end of this chapter I can analyse two-dimensional collisions by resolving momentum into perpendicular components.
\item By the end of this chapter I can define impulse, use the impulse-momentum theorem, and interpret force-time graphs, including applications to safety devices.
\item By the end of this chapter I can explain rocket propulsion and explosions using conservation of momentum.
\item By the end of this chapter I can locate the centre of mass of a system of particles and relate its motion to the total momentum of the system.
\item By the end of this chapter I can model real-life situations (vehicle crashes, recoil of a gun, sports impacts) quantitatively using momentum as a vector.
\end{itemize}
\end{objectifs}

\begin{prerequis}
\begin{itemize}
\item Newton's three laws of motion and their use with F = ma (Lower Sixth Mechanics)
\item Kinematics of motion in a straight line: velocity, acceleration, suvat equations
\item Kinetic energy and the work-energy principle
\item Vector algebra: components, resolution, addition of vectors, use of unit vectors i and j
\item Basic trigonometry and solving simultaneous linear equations
\end{itemize}
\end{prerequis}

\newpage

\coursec{Momentum, Force and Impulse}

During a busy morning at the Mile 4 motor park in Limbe, a fully loaded 70-seater bus and an empty taxi collide at low speed: the taxi is thrown violently backwards while the bus barely slows down. Yet a few kilometres away, a goalkeeper at the Limbe Omnisport stadium stops a fast football with his bare hands and feels only a sting, whereas the same ball hitting a concrete wall rebounds sharply. Why does the ``quantity of motion'' of a body depend on both its mass \emph{and} its velocity? Why do some collisions bounce and others stick? And how do engineers use these ideas to design safer vehicles and launch rockets? This chapter answers these questions by treating \cle{momentum} as a \emph{vector} quantity. We begin with the fundamental definitions: momentum, force and impulse.

\courssub{Linear Momentum: Definition and Vector Nature}

\textbf{An observation to build intuition.}\par
Imagine two objects moving towards you at $10\ \mathrm{m\,s^{-1}}$: a football and a lorry. You could stop the football with your hands; the lorry would be impossible to stop. Now imagine the same lorry moving at $0.5\ \mathrm{m\,s^{-1}}$: a few strong people could stop it. So the ``quantity of motion'' of a body grows with \emph{both} its mass and its velocity. Moreover, a lorry moving \emph{towards} you is very different from one moving \emph{away} from you: direction matters. This leads us to define momentum as a vector.

\begin{definition}[Linear momentum]
The \cle{linear momentum} $\vec{p}$ of a body of mass $m$ moving with velocity $\vec{v}$ is
\[
\vec{p} = m\vec{v}.
\]
Momentum is a \textbf{vector quantity}: it has the \emph{same direction as the velocity} of the body. Its SI unit is the kilogram-metre per second, $\mathrm{kg\,m\,s^{-1}}$, which is equivalent to the newton-second ($\mathrm{N\,s}$), since $1\ \mathrm{N} = 1\ \mathrm{kg\,m\,s^{-2}}$.
\end{definition}

In one-dimensional problems we work with the \emph{component} of momentum along a chosen positive direction: $p = mv$, where $v$ carries a sign ($+$ or $-$). This sign is the mathematical translation of the direction of motion.

\begin{exemplebox}[Momentum of the bus and the taxi]
A loaded bus of mass $8000\ \mathrm{kg}$ moves east at $10\ \mathrm{m\,s^{-1}}$; a taxi of mass $1000\ \mathrm{kg}$ moves west at $15\ \mathrm{m\,s^{-1}}$. Taking east as positive:
\[
p_{\text{bus}} = 8000 \times 10 = 80\,000\ \mathrm{kg\,m\,s^{-1}}, \qquad
p_{\text{taxi}} = 1000 \times (-15) = -15\,000\ \mathrm{kg\,m\,s^{-1}}.
\]
The bus has more than five times the magnitude of momentum of the taxi, in the opposite direction. This is why, in a collision, the bus barely notices the impact while the taxi is thrown backwards.
\end{exemplebox}

\begin{popfigure}
\begin{center}
\begin{tikzpicture}[scale=1.0]
% road
\draw[thick, popline] (-6.5,0) -- (6.5,0);
\foreach \x in {-6,-4,...,6} \draw[popmuted, dashed] (\x,0) -- (\x+0.9,0);
% bus
\draw[fill=popblueL, draw=popblue, thick] (-5.6,0.25) rectangle (-2.4,1.35);
\draw[fill=white, draw=popblue] (-5.2,0.85) rectangle (-4.6,1.2);
\draw[fill=white, draw=popblue] (-4.3,0.85) rectangle (-3.7,1.2);
\draw[fill=white, draw=popblue] (-3.4,0.85) rectangle (-2.8,1.2);
\draw[fill=popdark] (-5.0,0.25) circle (0.22);
\draw[fill=popdark] (-3.0,0.25) circle (0.22);
\node[popblue] at (-4.0,1.6) {\small bus $8000$ kg};
% taxi
\draw[fill=popgold, draw=poporange, thick] (2.6,0.25) rectangle (4.4,0.95);
\draw[fill=popgold, draw=poporange, thick] (3.0,0.95) -- (3.2,1.25) -- (3.9,1.25) -- (4.1,0.95);
\draw[fill=popdark] (3.0,0.25) circle (0.18);
\draw[fill=popdark] (4.0,0.25) circle (0.18);
\node[poporange] at (3.5,1.55) {\small taxi $1000$ kg};
% momentum vectors (scaled: 80000 -> 4cm, 15000 -> 0.75cm)
\draw[-{Stealth}, ultra thick, popblue] (-4.0,2.1) -- (0.0,2.1) node[midway, above, popblue] {\small $\vec{p}_{\text{bus}} = +80\,000$};
\draw[-{Stealth}, ultra thick, poporange] (3.5,2.1) -- (2.75,2.1) node[midway, above, poporange] {\small $\vec{p}_{\text{taxi}} = -15\,000$};
% velocity arrows
\draw[-{Stealth}, thick, popgreen] (-2.4,0.8) -- (-1.2,0.8) node[right, popgreen] {\small $+10\ \mathrm{m\,s^{-1}}$};
\draw[-{Stealth}, thick, popgreen] (2.6,0.6) -- (1.6,0.6) node[left, popgreen] {\small $-15\ \mathrm{m\,s^{-1}}$};
% positive direction
\draw[-{Stealth}, thick, popdark] (5.0,-0.6) -- (6.2,-0.6) node[right, popdark] {\small $+x$ (east)};
\end{tikzpicture}
\end{center}
\legende{Momentum vectors of a bus and a taxi moving in opposite directions. Arrow lengths are to scale: the bus carries far more momentum than the taxi, in the opposite direction.}
\end{popfigure}

\courssub{Momentum of a System of Particles}

\begin{definition}[Total momentum of a system]
The total momentum of a system of particles is the \textbf{vector sum} of the individual momenta:
\[
\vec{P} = \vec{p}_1 + \vec{p}_2 + \cdots + \vec{p}_n = \sum_{i=1}^{n} m_i \vec{v}_i.
\]
\end{definition}

In one dimension this vector sum becomes an \emph{algebraic} sum once a positive direction is chosen.

\begin{exoresolu}[Total momentum of two vehicles]
A lorry of mass $3000\ \mathrm{kg}$ travels east at $12\ \mathrm{m\,s^{-1}}$. A car of mass $900\ \mathrm{kg}$ travels west at $20\ \mathrm{m\,s^{-1}}$. Find the total momentum of the system.
\tcblower
\textbf{Solution.} Take east as positive.
\begin{align*}
P &= m_1 v_1 + m_2 v_2 = 3000(12) + 900(-20)\\
&= 36\,000 - 18\,000 = 18\,000\ \mathrm{kg\,m\,s^{-1}}.
\end{align*}
The total momentum is $18\,000\ \mathrm{kg\,m\,s^{-1}}$ directed \textbf{east} (the positive sign tells us the direction).
\end{exoresolu}

\begin{attention}[Momentum is a vector -- mind the signs!]
The single most common error in GCE momentum questions is forgetting that a body moving in the negative direction has \emph{negative} velocity and therefore \emph{negative} momentum. When a ball rebounds, its velocity \emph{changes sign}: if $u = +5\ \mathrm{m\,s^{-1}}$ and it rebounds at $3\ \mathrm{m\,s^{-1}}$, then $v = -3\ \mathrm{m\,s^{-1}}$, \textbf{not} $+3$. Always state your positive direction \emph{before} substituting numbers.
\end{attention}

\courssub{Newton's Second Law in Momentum Form}

Newton originally stated his second law not as $F = ma$ but in terms of the change of ``motion'' (momentum). This more general form is the one required at Advanced Level.

\begin{propriete}[Newton's second law -- general form (proof examinable)]
The resultant force acting on a body equals the \textbf{rate of change of its momentum}:
\[
\vec{F} = \frac{d\vec{p}}{dt} = \frac{d}{dt}(m\vec{v}).
\]
\textbf{Derivation of $\vec{F} = m\vec{a}$ for constant mass.} If the mass $m$ is constant, it can be taken out of the derivative:
\[
\vec{F} = \frac{d}{dt}(m\vec{v}) = m\frac{d\vec{v}}{dt} = m\vec{a}.
\]
Thus $F = ma$ is a \emph{special case} of $F = \dfrac{dp}{dt}$, valid only when the mass does not change.
\end{propriete}

\textbf{When is $F = \dfrac{dp}{dt}$ more general than $F = ma$?}\par
If the mass varies with time, the product rule gives
\[
\vec{F} = \frac{d}{dt}(m\vec{v}) = m\frac{d\vec{v}}{dt} + \vec{v}\frac{dm}{dt},
\]
which is \emph{not} simply $m\vec{a}$. Variable-mass systems include:
\begin{itemize}
\item a \cle{rocket} burning fuel and ejecting exhaust gases (its mass decreases);
\item a conveyor belt onto which sand falls at a steady rate (mass increases);
\item a raindrop growing as it falls through a cloud.
\end{itemize}
For such systems, only the momentum form of Newton's second law is correct. We shall exploit this in the study of rocket motion later in the chapter.

\courssub{Impulse and the Impulse--Momentum Theorem}

\textbf{Intuition.}\par
To change the momentum of a body you must apply a force \emph{for some time}. A small force applied for a long time can produce the same change of momentum as a large force applied briefly. The quantity that captures ``force $\times$ duration'' is called impulse.

\begin{definition}[Impulse]
The \cle{impulse} $\vec{J}$ of a constant force $\vec{F}$ acting for a time interval $\Delta t$ is
\[
\vec{J} = \vec{F}\,\Delta t.
\]
For a variable force, the impulse over $[t_1, t_2]$ is
\[
\vec{J} = \int_{t_1}^{t_2} \vec{F}\, dt.
\]
Impulse is a \textbf{vector} in the direction of the force; its SI unit is the newton-second ($\mathrm{N\,s}$), equivalent to $\mathrm{kg\,m\,s^{-1}}$.
\end{definition}

\begin{propriete}[Impulse--momentum theorem (proof examinable)]
The impulse of the resultant force on a body equals the \textbf{change of momentum} of the body:
\[
\vec{J} = \Delta \vec{p} = m\vec{v} - m\vec{u},
\]
where $\vec{u}$ and $\vec{v}$ are the velocities before and after the force acts.
\textbf{Proof.} From Newton's second law, $\vec{F} = \dfrac{d\vec{p}}{dt}$. Integrating between $t_1$ and $t_2$:
\[
\vec{J} = \int_{t_1}^{t_2} \vec{F}\, dt = \int_{t_1}^{t_2} \frac{d\vec{p}}{dt}\, dt = \vec{p}(t_2) - \vec{p}(t_1) = m\vec{v} - m\vec{u}. \qquad \blacksquare
\]
\end{propriete}

\begin{methode}[Solving impulse problems]
\begin{enumerate}
\item \textbf{Choose a positive direction} and state it clearly (e.g.\ ``take the initial direction of motion as positive'').
\item \textbf{Write all velocities with their signs}: $u$ before, $v$ after. A reversed velocity is negative.
\item \textbf{Compute the change of momentum}: $\Delta p = m(v - u)$.
\item \textbf{Apply the theorem}: $J = \Delta p$, and if the contact time $\Delta t$ is known, the average force is
\[
F_{\text{avg}} = \frac{\Delta p}{\Delta t} = \frac{m(v-u)}{\Delta t}.
\]
\item \textbf{Interpret the sign} of $J$ or $F$: it gives the direction of the impulse/force relative to your chosen positive direction.
\end{enumerate}
\end{methode}

\begin{exoresolu}[Ball bouncing off a wall]
A ball of mass $0.5\ \mathrm{kg}$ moving horizontally at $8\ \mathrm{m\,s^{-1}}$ strikes a wall perpendicularly and rebounds at $6\ \mathrm{m\,s^{-1}}$. The contact with the wall lasts $0.05\ \mathrm{s}$. Find (a) the impulse on the ball; (b) the average force exerted by the wall.
\tcblower
\textbf{Solution.} Take the initial direction of motion (towards the wall) as positive. Then $u = +8\ \mathrm{m\,s^{-1}}$ and, since the ball rebounds, $v = -6\ \mathrm{m\,s^{-1}}$.

(a) Impulse:
\[
J = m(v - u) = 0.5(-6 - 8) = 0.5 \times (-14) = -7\ \mathrm{N\,s}.
\]
The impulse has magnitude $7\ \mathrm{N\,s}$, directed \emph{away from the wall}.

(b) Average force:
\[
F_{\text{avg}} = \frac{J}{\Delta t} = \frac{-7}{0.05} = -140\ \mathrm{N}.
\]
The wall pushes the ball with an average force of $140\ \mathrm{N}$ directed away from the wall. Note that had we written $v = +6$, we would have obtained the absurd answer $J = -1\ \mathrm{N\,s}$: the sign convention is essential.
\end{exoresolu}

\courssub{Force--Time Graphs}

In real impacts (a ball on a bat, a car in a crash test) the force is \emph{not} constant: it rises from zero, peaks, and falls back to zero. Since $J = \displaystyle\int F\, dt$, the impulse has a beautiful graphical interpretation.

\begin{propriete}[Impulse as an area]
The impulse of a force over a time interval equals the \textbf{area under the force--time graph} between the curve and the time axis:
\[
J = \int_{t_1}^{t_2} F\, dt = \text{area under the } F\text{-}t \text{ curve}.
\]
The \textbf{average force} is then the constant force that would give the same area over the same duration:
\[
F_{\text{avg}} = \frac{J}{\Delta t} = \frac{\text{area under the curve}}{\text{duration}}.
\]
\end{propriete}

\begin{popfigure}
\begin{center}
\begin{tikzpicture}
\begin{axis}[
    width=11cm, height=6.2cm,
    axis lines=left,
    xlabel={$t$ (s)}, ylabel={$F$ (N)},
    xmin=0, xmax=0.09, ymin=0, ymax=180,
    xtick={0,0.02,0.04,0.06,0.08},
    ytick={0,50,100,150},
    tick label style={font=\small},
    label style={font=\small},
]
\addplot[draw=none, name path=xaxis] coordinates {(0,0) (0.08,0)};
\addplot[very thick, popblue, smooth, name path=curve] coordinates {
    (0,0) (0.01,40) (0.02,110) (0.03,160) (0.04,150)
    (0.05,95) (0.06,45) (0.07,10) (0.08,0)
};
\addplot[popblueL] fill between[of=curve and xaxis];
\node[popblue, font=\small] at (axis cs:0.04,70) {$J = \text{ area } = \Delta p$};
\node[popmuted, font=\small] at (axis cs:0.062,140) {ball--wall impact};
\end{axis}
\end{tikzpicture}
\end{center}
\legende{Force--time graph for a ball bouncing off a wall. The shaded area under the curve is the impulse, equal to the change of momentum of the ball.}
\end{popfigure}

\begin{exemplebox}[Average force from a graph]
Suppose the area under the curve above is estimated (by counting squares or by trapezia) as $7.0\ \mathrm{N\,s}$ over a contact time of $0.08\ \mathrm{s}$. Then
\[
F_{\text{avg}} = \frac{7.0}{0.08} = 87.5\ \mathrm{N},
\]
even though the \emph{peak} force reaches about $160\ \mathrm{N}$. The average force is what the impulse--momentum theorem delivers directly.
\end{exemplebox}

\begin{savaistu}[A frequent GCE question]
``Show that the area under a force--time graph represents the impulse'' and ``estimate the impulse from the given graph by counting squares'' appear regularly in GCE Advanced Level papers. Practise both: the theoretical link $J = \int F\,dt = \Delta p$, and the numerical skill of estimating an irregular area (count full squares, group partial squares, or use the trapezium rule).
\end{savaistu}

\courssub{Application to Safety: Increasing Impact Time}

Rewrite the impulse--momentum theorem as
\[
F_{\text{avg}} = \frac{\Delta p}{\Delta t}.
\]
In any given impact, the change of momentum $\Delta p$ is fixed (a car of given mass must be brought from its speed to rest). The only way to \textbf{reduce the average force} on the occupants is therefore to \textbf{increase the impact time} $\Delta t$. This single idea underlies most modern safety engineering:

\begin{itemize}
\item \textbf{Airbags}: the driver's head sinks into the inflating bag, so it stops over perhaps $0.1\ \mathrm{s}$ instead of $0.01\ \mathrm{s}$ against the steering wheel -- the average force is divided by about ten.
\item \textbf{Crumple zones}: the front and rear of a vehicle are designed to deform progressively, lengthening the collision time while the passenger cell stays rigid.
\item \textbf{Seatbelts}: they stretch slightly and spread the force over the strong parts of the body (chest, pelvis), extending the stopping time.
\item \textbf{Bending the knees on landing}: a paratrooper or basketball player who lands with bent legs extends the stopping time of his body's momentum, sparing his joints.
\item \textbf{Padded goalposts and crash mats}: the padding compresses, increasing $\Delta t$ and reducing $F_{\text{avg}}$.
\end{itemize}

\begin{popfigure}
\begin{center}
\begin{tikzpicture}
\begin{axis}[
    width=11.5cm, height=6.5cm,
    axis lines=left,
    xlabel={$t$ (s)}, ylabel={$F$ (N)},
    xmin=0, xmax=0.35, ymin=0, ymax=5200,
    xtick={0,0.05,0.1,0.15,0.2,0.25,0.3},
    ytick={0,1000,2000,3000,4000,5000},
    tick label style={font=\small},
    label style={font=\small},
    legend style={font=\small, at={(0.98,0.97)}, anchor=north east},
]
% hard impact: narrow tall triangle, area = 0.5*0.04*5000 = 100 N s
\addplot[very thick, poppink, smooth, name path=hard] coordinates {(0.02,0) (0.04,5000) (0.06,0)};
\addplot[draw=none, name path=axh] coordinates {(0.02,0) (0.06,0)};
\addplot[poppink, opacity=0.25] fill between[of=hard and axh];
% airbag: wide low triangle, area = 0.5*0.2*1000 = 100 N s
\addplot[very thick, popgreen, smooth, name path=soft] coordinates {(0.05,0) (0.15,1000) (0.25,0)};
\addplot[draw=none, name path=axs] coordinates {(0.05,0) (0.25,0)};
\addplot[popgreenL] fill between[of=soft and axs];
\legend{hard surface, with airbag}
\node[font=\small, popdark] at (axis cs:0.24,3300) {equal areas};
\node[font=\small, popdark] at (axis cs:0.24,2900) {same $\Delta p$};
\end{axis}
\end{tikzpicture}
\end{center}
\legende{The airbag principle: both force--time curves enclose the same area (same change of momentum), but the airbag spreads the impact over a longer time, so the peak and average forces are much smaller.}
\end{popfigure}

\begin{exoresolu}[Why the goalkeeper pulls his hands back]
A $0.43\ \mathrm{kg}$ football arrives at a goalkeeper's hands at $25\ \mathrm{m\,s^{-1}}$ and is stopped. (a) Find the change of momentum. (b) Find the average force on the hands if the ball is stopped in $0.02\ \mathrm{s}$ (rigid hands) and in $0.15\ \mathrm{s}$ (hands drawn back).
\tcblower
\textbf{Solution.} Take the ball's initial direction as positive: $u = 25\ \mathrm{m\,s^{-1}}$, $v = 0$.

(a) $\Delta p = m(v-u) = 0.43(0 - 25) = -10.75\ \mathrm{kg\,m\,s^{-1}}$; magnitude $10.75\ \mathrm{kg\,m\,s^{-1}}$.

(b) Rigid hands: $F_{\text{avg}} = \dfrac{10.75}{0.02} \approx 538\ \mathrm{N}$ -- painful!

Hands drawn back: $F_{\text{avg}} = \dfrac{10.75}{0.15} \approx 72\ \mathrm{N}$ -- a mere sting.

Multiplying the impact time by $7.5$ divides the force by $7.5$: this is exactly why good goalkeepers ``give'' with the ball.
\end{exoresolu}

\begin{attention}[Impulse vs.\ force: do not confuse them]
Impulse ($\mathrm{N\,s}$) measures the \emph{total effect} of a force over time and equals the change of momentum; force ($\mathrm{N}$) is the \emph{rate} of change of momentum. Writing $J = F\Delta t$ is only valid for a \emph{constant} force (or when $F$ is the average force). For a varying force you must use the area under the $F$--$t$ graph or the integral $\int F\,dt$.
\end{attention}

\begin{aretenir}[Key points of this section]
\begin{itemize}
\item Linear momentum: $\vec{p} = m\vec{v}$, a vector in the direction of motion; unit $\mathrm{kg\,m\,s^{-1}} \equiv \mathrm{N\,s}$.
\item Total momentum of a system = \emph{vector} sum of individual momenta; in 1-D, use signs.
\item Newton's second law, general form: $\vec{F} = \dfrac{d\vec{p}}{dt}$; it reduces to $\vec{F} = m\vec{a}$ only when mass is constant.
\item Impulse: $\vec{J} = \vec{F}\Delta t$ (constant force) or $\vec{J} = \int \vec{F}\,dt$ (variable force); equals change of momentum $\Delta \vec{p}$.
\item Impulse = area under $F$--$t$ graph.
\item To reduce impact force, increase impact time (airbags, crumple zones, seatbelts, bending knees).
\end{itemize}
\end{aretenir}

\coursec{Conservation of Linear Momentum}

In the previous section we established that the momentum $\vec{p}=m\vec{v}$ of a body changes only when a net external impulse acts on it, and that $\vec{F}_{\text{net}}=\dfrac{d\vec{p}}{dt}$. We now ask a deeper question: what happens when \emph{two or more bodies interact with each other}? The answer is one of the most powerful principles in all of physics — the \cle{principle of conservation of linear momentum}. It allows us to predict the outcome of collisions, explosions and recoils \emph{without knowing any detail of the forces involved}.

\courssub{Internal and External Forces: the Idea of a System}

\begin{experience}[Two skaters push each other]
Imagine two students on roller skates, standing face to face on a smooth floor. Student $A$ pushes student $B$. Both students move — $B$ moves forward and $A$ moves backward. If we look at each student \emph{separately}, each one's momentum has changed, because each experienced a force from the other. But if we look at the \emph{pair together}, something remarkable happens: the total momentum of the pair is still zero, exactly as it was before the push. The push was an \emph{internal} affair of the pair.
\end{experience}

To make this precise we must first say what we mean by a \cle{system}.

\begin{definition}[System, internal and external forces]
A \cle{system} is any collection of bodies that we choose to consider together. Forces acting on the bodies of the system are of two kinds:
\begin{itemize}
\item \cle{Internal forces}: forces exerted \emph{between bodies of the system} (e.g.\ the contact forces between two colliding trolleys, the explosive forces between fragments of a shell).
\item \cle{External forces}: forces exerted on bodies of the system \emph{by agents outside the system} (e.g.\ weight, normal reaction of the ground, friction from the road, air resistance).
\end{itemize}
\end{definition}

\begin{definition}[Isolated (closed) system]
An \cle{isolated system} (or closed system) is a system on which the \textbf{resultant external force is zero}:
\[
\sum \vec{F}_{\text{ext}} = \vec{0}.
\]
This happens either because no external forces act at all, or because the external forces cancel one another (e.g.\ weight balanced by the normal reaction on a smooth horizontal table).
\end{definition}

\begin{attention}[Choosing the system is your decision]
The same interaction can be isolated or not, depending on what you include in the system. A single billiard ball struck by another is \emph{not} isolated (it feels the contact force from outside). The \emph{pair} of billiard balls on a smooth table \emph{is} isolated, because the contact forces are now internal and the external forces (weight, reaction) cancel.
\end{attention}

\courssub{Statement of the Principle}

\begin{propriete}[Principle of conservation of linear momentum]
If the resultant external force on a system is zero (isolated system), then the \textbf{total linear momentum of the system remains constant}:
\[
\vec{P}_{\text{total}} = \vec{p}_1+\vec{p}_2+\dots+\vec{p}_n = \text{constant},
\]
that is,
\[
\boxed{\;\vec{P}_{\text{before}} = \vec{P}_{\text{after}}\;}
\]
for any two instants before and after an internal interaction. This is a \textbf{vector} statement: it holds component by component along each independent direction. The derivation from Newton's laws is examinable.
\end{propriete}

\courssub{Derivation from Newton's Second and Third Laws}

Consider two bodies $A$ and $B$, of masses $m_1$ and $m_2$, forming an isolated system. During their interaction, body $A$ exerts a force $\vec{F}_{A\text{ on }B}$ on $B$, and $B$ exerts $\vec{F}_{B\text{ on }A}$ on $A$.

\textbf{Step 1 — Newton's third law.} The interaction forces are equal and opposite at every instant:
\[
\vec{F}_{B\text{ on }A} = -\,\vec{F}_{A\text{ on }B}.
\]

\textbf{Step 2 — Newton's second law in momentum form.} Since the system is isolated, these internal forces are the \emph{only} forces acting on each body, so
\[
\vec{F}_{B\text{ on }A} = \frac{d\vec{p}_1}{dt}, \qquad
\vec{F}_{A\text{ on }B} = \frac{d\vec{p}_2}{dt}.
\]

\textbf{Step 3 — Add.} Adding the two equations and using Step 1:
\[
\frac{d\vec{p}_1}{dt}+\frac{d\vec{p}_2}{dt}
= \vec{F}_{B\text{ on }A}+\vec{F}_{A\text{ on }B} = \vec{0}.
\]
Hence
\[
\frac{d}{dt}\left(\vec{p}_1+\vec{p}_2\right)=\vec{0}
\quad\Longrightarrow\quad
\vec{p}_1+\vec{p}_2 = \text{constant}. \qquad \blacksquare
\]

The argument extends immediately to $n$ bodies, because internal forces always occur in Newton's-third-law pairs whose vector sum is zero. Notice what the proof reveals: internal forces \emph{redistribute} momentum between the bodies, but can never change the \emph{total}.

\begin{aretenir}[The vector nature of conservation]
Conservation of momentum is a vector law. In two or three dimensions it gives one scalar equation per axis:
\[
\sum m_i u_{ix} = \sum m_i v_{ix}, \qquad
\sum m_i u_{iy} = \sum m_i v_{iy}.
\]
Momentum along each independent direction is conserved \emph{separately}. A body can gain momentum in the $x$-direction while the total $y$-momentum stays zero.
\end{aretenir}

\courssub{A Working Routine}

\begin{methode}[Conservation-of-momentum routine]
\begin{enumerate}
\item \textbf{Define the system}: decide which bodies you will consider together.
\item \textbf{Check isolation}: confirm that the resultant external force is zero, or that external impulses are negligible during the short interaction.
\item \textbf{Draw before/after diagrams}: sketch the situation just before and just after the interaction, labelling all velocities.
\item \textbf{Choose a positive direction} (one per axis) and stick to it; velocities in the opposite direction carry a minus sign.
\item \textbf{Write the vector equation} $\vec{P}_{\text{before}}=\vec{P}_{\text{after}}$, i.e.\ $\sum m_i\vec{u}_i=\sum m_i\vec{v}_i$, and solve component by component.
\item \textbf{Interpret the sign} of your answer: a negative velocity means motion opposite to the chosen positive direction.
\end{enumerate}
\end{methode}

\courssub{Applications in One Dimension}

\textbf{Two bodies pushing apart.} Two trolleys of masses $m_1$ and $m_2$, initially at rest and touching, are pushed apart by a spring between them. The initial total momentum is zero, so
\[
0 = m_1 v_1 + m_2 v_2 \quad\Longrightarrow\quad \frac{v_1}{v_2}=-\frac{m_2}{m_1}.
\]
The lighter trolley moves faster, in the opposite direction. The momenta are equal in magnitude and opposite in direction.

\textbf{Man jumping from a canoe.} A fisherman of mass $m$ stands on a stationary canoe of mass $M$ on calm water (resistance negligible). He jumps horizontally off the stern with velocity $\vec{v}$ relative to the ground. The system \{man + canoe\} is isolated horizontally, with zero initial momentum:
\[
0 = m\vec{v} + M\vec{V} \quad\Longrightarrow\quad \vec{V}=-\frac{m}{M}\vec{v}.
\]
The canoe glides backwards — this is why stepping carelessly out of a dugout canoe on Lake Nyos can leave you in the water while the canoe drifts away!

\begin{attention}[Velocities are measured relative to the ground]
In all conservation-of-momentum equations, every velocity must be expressed in the \emph{same} reference frame, normally the ground (or the water at rest). If a problem gives the man's velocity \emph{relative to the canoe}, you must first convert it to a ground-frame velocity before substituting.
\end{attention}

\begin{popfigure}
\begin{center}
\begin{tikzpicture}[scale=0.9, every node/.style={font=\small}]
% BEFORE
\node[font=\small\bfseries] at (2.2,2.6) {Before (at rest)};
\draw[popblue, thick] (0.4,0.4) .. controls (0.7,0.9) and (3.7,0.9) .. (4.0,0.4) -- (3.6,0.1) -- (0.8,0.1) -- cycle;
\fill[popblueL] (0.8,0.12) rectangle (3.6,0.38);
\draw[fill=poporange] (2.2,1.15) circle (0.16);
\draw[thick] (2.2,0.99) -- (2.2,0.55) (2.2,0.85) -- (1.9,0.65) (2.2,0.85) -- (2.5,0.65) (2.2,0.55) -- (2.0,0.28) (2.2,0.55) -- (2.4,0.28);
\draw[popblue!60!black] (-0.2,0.05) -- (4.6,0.05);
\node[popmuted] at (2.2,-0.25) {Lake Nyos};
\node at (2.2,1.9) {$\vec{P}_{\text{total}}=\vec{0}$};
% AFTER
\begin{scope}[xshift=7.2cm]
\node[font=\small\bfseries] at (2.2,2.6) {After the jump};
\draw[popblue, thick] (0.4,0.4) .. controls (0.7,0.9) and (3.7,0.9) .. (4.0,0.4) -- (3.6,0.1) -- (0.8,0.1) -- cycle;
\fill[popblueL] (0.8,0.12) rectangle (3.6,0.38);
\draw[fill=poporange] (4.9,1.15) circle (0.16);
\draw[thick] (4.9,0.99) -- (4.9,0.55) (4.9,0.85) -- (4.6,0.65) (4.9,0.85) -- (5.2,0.65) (4.9,0.55) -- (4.7,0.28) (4.9,0.55) -- (5.1,0.28);
\draw[popblue!60!black] (-0.2,0.05) -- (5.6,0.05);
\draw[-{Stealth}, very thick, popgreen] (4.9,1.5) -- (6.1,1.5) node[midway, above]{$m\vec{v}$};
\draw[-{Stealth}, very thick, poppink] (1.6,0.7) -- (0.4,0.7) node[midway, above]{$M\vec{V}$};
\node at (2.2,1.9) {$m\vec{v}+M\vec{V}=\vec{0}$};
\end{scope}
\end{tikzpicture}
\end{center}
\legende{Man jumping from a stationary canoe: the forward momentum of the man is balanced by the backward momentum of the canoe, so the total momentum remains zero.}
\end{popfigure}

\begin{exoresolu}[Shunting railway wagons]
A railway wagon of mass $12\,000\ \mathrm{kg}$ moving at $4.0\ \mathrm{m\,s^{-1}}$ on a straight level track collides with a stationary wagon of mass $8\,000\ \mathrm{kg}$, and the two couple together. Find their common velocity just after coupling.
\tcblower
\textbf{System:} the two wagons. \textbf{Isolation:} during the brief coupling, external impulses (friction) are negligible; weight cancels the reaction.

Take the initial direction of motion as positive.
\begin{align*}
P_{\text{before}} &= 12000\times 4.0 + 8000\times 0 = 48000\ \mathrm{kg\,m\,s^{-1}},\\
P_{\text{after}} &= (12000+8000)\,v = 20000\,v.
\end{align*}
Conservation of momentum:
\[
20000\,v = 48000 \quad\Longrightarrow\quad v = 2.4\ \mathrm{m\,s^{-1}}.
\]
The coupled wagons move on at $\SI{2.4}{m.s^{-1}}$ in the original direction. (Note: kinetic energy is \emph{lost} here — the coupling is inelastic — yet momentum is exactly conserved.)
\end{exoresolu}

\courssub{Recoil of a Gun}

When a rifle fires, the expanding gases push the bullet forward and the rifle backward with equal and opposite forces (Newton's third law). The system \{rifle + bullet\} is isolated during the very short firing, and its total momentum, initially zero, stays zero:
\[
m_b\vec{v}_b + m_r\vec{v}_r = \vec{0}
\quad\Longrightarrow\quad
\vec{v}_r = -\frac{m_b}{m_r}\vec{v}_b.
\]
The \cle{recoil velocity} of the rifle is small because its mass is large compared with the bullet's — but the \emph{momenta} of bullet and rifle are equal in magnitude.

\begin{popfigure}
\begin{center}
\begin{tikzpicture}[scale=1.0, every node/.style={font=\small}]
% rifle body
\draw[thick, fill=popgold!40] (0,0.55) rectangle (3.4,0.95);
\draw[thick, fill=popgold!40] (3.4,0.68) rectangle (5.2,0.82);
\draw[thick, fill=popgold!60] (0,0.55) -- (-0.5,0.1) -- (0.1,0.1) -- (0.55,0.55) -- cycle;
\draw[thick] (1.2,0.55) -- (1.35,0.2) -- (1.6,0.2) -- (1.7,0.55);
\node at (1.7,1.25) {rifle, mass $m_r$};
% bullet
\draw[fill=poppink] (5.6,0.75) circle (0.09);
\node[above] at (5.6,0.86) {bullet, mass $m_b$};
% momentum arrows
\draw[-{Stealth}, very thick, popgreen] (5.75,1.25) -- (7.3,1.25) node[midway, above]{$m_b\vec{v}_b$};
\draw[-{Stealth}, very thick, poppurple] (0.4,1.45) -- (-1.2,1.45) node[midway, above]{$m_r\vec{v}_r$};
\node[popmuted] at (3.4,-0.35) {$m_b\vec{v}_b + m_r\vec{v}_r = \vec{0}$};
\end{tikzpicture}
\end{center}
\legende{Recoil of a rifle: after firing, the forward momentum of the bullet and the backward momentum of the rifle are equal in magnitude and opposite in direction.}
\end{popfigure}

\begin{exoresolu}[Computing the recoil velocity]
A rifle of mass $4.0\ \mathrm{kg}$ fires a bullet of mass $10\ \mathrm{g}$ with a muzzle velocity of $400\ \mathrm{m\,s^{-1}}$. The rifle is held loosely. Calculate the recoil velocity of the rifle.
\tcblower
The system \{rifle + bullet\} is initially at rest: $P_{\text{before}}=0$. Taking the bullet's direction as positive, with $m_b = 0.010\ \mathrm{kg}$:
\[
0 = m_b v_b + m_r v_r
\quad\Longrightarrow\quad
v_r = -\frac{m_b v_b}{m_r} = -\frac{0.010\times 400}{4.0} = -1.0\ \mathrm{m\,s^{-1}}.
\]
The rifle recoils at $\SI{1.0}{m.s^{-1}}$ opposite to the bullet's motion. The negative sign indicates the direction of recoil. (In practice a firm shoulder adds an external impulse and reduces the free recoil.)
\end{exoresolu}

\courssub{Condition of Validity: Negligible External Impulse}

Strictly, conservation requires $\sum\vec{F}_{\text{ext}}=\vec{0}$. In real collisions there is often friction or gravity acting. Why can we still apply the principle?

The change in total momentum equals the \emph{impulse of external forces}:
\[
\Delta\vec{P} = \vec{J}_{\text{ext}} = \sum \vec{F}_{\text{ext}}\,\Delta t.
\]
During a collision or explosion, $\Delta t$ is extremely short (milliseconds), while the internal forces are enormous. The external impulse $\vec{F}_{\text{ext}}\Delta t$ is then negligible compared with the momenta involved, so
\[
\vec{P}_{\text{just before}} \approx \vec{P}_{\text{just after}}.
\]

\begin{attention}[Apply conservation only across the short interaction]
Conservation of momentum is applied between instants \emph{immediately before} and \emph{immediately after} the collision or explosion. Do not extend it over long time intervals during which friction, gravity or other external forces have time to act significantly.
\end{attention}

\courssub{Conservation of Momentum versus Conservation of Energy}

A frequent source of confusion at Advanced Level is the difference between these two great conservation laws.

\begin{center}
\begin{tabularx}{\linewidth}{|l|X|X|}
\hline
\rowcolor{popblueL} & \textbf{Linear momentum} & \textbf{Kinetic energy} \\
\hline
Nature & Vector ($m\vec{v}$) & Scalar ($\tfrac12 mv^2$) \\
\hline
Conserved in isolated system? & \textbf{Always} (Newton's third law) & Only if the collision is \textbf{elastic} \\
\hline
In inelastic collisions / explosions & Conserved & Converted to heat, sound, deformation (or released by explosives) \\
\hline
\end{tabularx}
\end{center}

In the shunting example above, the kinetic energy before coupling was $\tfrac12(12000)(4.0)^2 = 96\,000\ \mathrm{J}$, while afterwards it is $\tfrac12(20000)(2.4)^2 = 57\,600\ \mathrm{J}$: some $38\,400\ \mathrm{J}$ was dissipated in the coupling mechanism — yet momentum was conserved exactly.

\begin{attention}[Momentum is conserved for the SYSTEM, not for each body]
Each colliding body generally \emph{changes} its momentum — that is the whole point of the interaction. It is the \emph{vector sum} of the momenta of all bodies in the isolated system that stays constant. Writing $m_1u_1 = m_1v_1$ for a single body is a serious error.
\end{attention}

\begin{attention}[Never assume kinetic energy is conserved]
Unless the collision is \emph{stated} (or shown) to be elastic, do not write $\tfrac12 m_1u_1^2+\tfrac12 m_2u_2^2 = \tfrac12 m_1v_1^2+\tfrac12 m_2v_2^2$. Momentum conservation alone, applied as a vector equation, is the safe universal tool; energy conservation is a bonus valid only for elastic interactions.
\end{attention}

\begin{savaistu}[Why conservation of momentum is so fundamental]
Newton's laws are an approximation valid at ordinary speeds, but the conservation of linear momentum survives \emph{everywhere}: in relativistic physics, in quantum mechanics, and in the collisions of galaxies. Physicists regard it as a consequence of a deep symmetry of nature — the fact that the laws of physics are the same everywhere in space (homogeneity of space). When you solve a GCE collision problem, you are using one of the deepest principles known to science!
\end{savaistu}

\begin{exercice}[Practice]
\begin{enumerate}
\item A man of mass $70\ \mathrm{kg}$ stands at rest on a stationary canoe of mass $90\ \mathrm{kg}$. He jumps off horizontally at $2.5\ \mathrm{m\,s^{-1}}$ relative to the ground. Find the velocity of the canoe just after he leaves it. (Water resistance negligible.)
\item A gun of mass $5.0\ \mathrm{kg}$, free to move, fires a $20\ \mathrm{g}$ shell at $250\ \mathrm{m\,s^{-1}}$. Find the recoil speed of the gun and the ratio of the kinetic energies of shell and gun. Comment.
\item A trolley of mass $2.0\ \mathrm{kg}$ moving at $3.0\ \mathrm{m\,s^{-1}}$ catches up with a trolley of mass $1.0\ \mathrm{kg}$ moving at $1.0\ \mathrm{m\,s^{-1}}$ in the same direction; they stick together. Find their common velocity and the kinetic energy lost.
\end{enumerate}
\end{exercice}

\begin{corrige}[Exercise 1]
System \{man + canoe\}, isolated horizontally, initial momentum zero. Taking the man's jump direction as positive:
\[
0 = 70(2.5) + 90V \quad\Longrightarrow\quad V = -\frac{175}{90} \approx -1.94\ \mathrm{m\,s^{-1}}.
\]
The canoe moves backwards at about $\SI{1.9}{m.s^{-1}}$.
\end{corrige}

\begin{corrige}[Exercise 2]
$0 = (0.020)(250) + 5.0\,v_r \Rightarrow v_r = -1.0\ \mathrm{m\,s^{-1}}$; recoil speed $\SI{1.0}{m.s^{-1}}$.
Kinetic energies: shell $=\tfrac12(0.020)(250)^2 = 625\ \mathrm{J}$; gun $=\tfrac12(5.0)(1.0)^2 = 2.5\ \mathrm{J}$. Ratio $250:1$. Although the momenta are equal, the light shell carries almost all the kinetic energy — equal momenta do \emph{not} mean equal energies.
\end{corrige}

\begin{corrige}[Exercise 3]
Momentum: $2.0(3.0)+1.0(1.0) = 3.0\,v \Rightarrow v = \dfrac{7.0}{3.0} \approx 2.33\ \mathrm{m\,s^{-1}}$.
KE before $=\tfrac12(2.0)(9)+\tfrac12(1.0)(1)=9.5\ \mathrm{J}$; KE after $=\tfrac12(3.0)\left(\tfrac{7}{3}\right)^2 \approx 8.17\ \mathrm{J}$. Energy lost $\approx 1.33\ \mathrm{J}$, dissipated as heat, sound and deformation in the coupling.
\end{corrige}

\coursec{Collisions in One Dimension: Elastic and Inelastic}

In the previous section we established the \cle{Principle of Conservation of Linear Momentum}: whenever two bodies interact and no external force acts on the system, the total momentum before the interaction equals the total momentum after it. We now put this principle to work on the most important family of interactions in mechanics: \cle{collisions}. A collision is a brief, intense interaction between two bodies during which they exert large forces on each other for a short time. Think of two trolleys meeting on a linear air track in the school laboratory, of two lorries crashing on the Douala--Yaound\'e highway, or of a bullet striking a wooden block. In every case momentum is conserved; what distinguishes one collision from another is what happens to the \cle{kinetic energy}.

\courssub{Classifying Collisions}

\begin{definition}[Elastic, inelastic and perfectly inelastic collisions]
A \cle{collision} is an interaction between two bodies in which they exert forces on each other over a short time interval.
\begin{itemize}
\item A collision is \cle{elastic} if \textbf{kinetic energy is conserved}: the total kinetic energy before impact equals the total kinetic energy after impact.
\item A collision is \cle{inelastic} if some kinetic energy is lost (converted into heat, sound and permanent deformation).
\item A collision is \cle{perfectly inelastic} (or \emph{coalescing}) if the two bodies \textbf{stick together} after impact and move with a common velocity. This is the case of maximum kinetic energy loss consistent with momentum conservation.
\end{itemize}
In \textbf{all} collisions, total linear momentum is conserved (provided no external impulse acts).
\end{definition}

\begin{attention}[Momentum yes, kinetic energy not always]
The single most common error at GCE is to write ``kinetic energy before = kinetic energy after'' for \emph{every} collision. Kinetic energy is conserved \textbf{only} in elastic collisions. If the bodies stick together, or if the question mentions heat, sound or deformation, kinetic energy is \emph{not} conserved. Momentum, however, is conserved in \emph{every} collision.
\end{attention}

\courssub{Elastic Collisions in One Dimension}

Consider two particles of masses $m_1$ and $m_2$ moving along the same straight line with velocities $u_1$ and $u_2$ ($u_1 > u_2$, so that they collide). After the elastic collision their velocities are $v_1$ and $v_2$. Two conservation laws apply:

\begin{align*}
\text{Momentum:} \quad & m_1 u_1 + m_2 u_2 = m_1 v_1 + m_2 v_2 \\
\text{Kinetic energy:} \quad & \tfrac{1}{2} m_1 u_1^2 + \tfrac{1}{2} m_2 u_2^2 = \tfrac{1}{2} m_1 v_1^2 + \tfrac{1}{2} m_2 v_2^2
\end{align*}

These are two simultaneous equations in the two unknowns $v_1$ and $v_2$. The energy equation contains squares, which makes direct solution heavy; fortunately there is a beautiful shortcut.

\begin{propriete}[Speed of approach equals speed of separation]
In a one-dimensional elastic collision,
\[ u_1 - u_2 = v_2 - v_1, \]
that is, the \cle{speed of approach} before impact equals the \cle{speed of separation} after impact.
\end{propriete}

\textbf{Proof (examinable).} Rearrange the momentum equation by grouping terms in each mass:
\[ m_1(u_1 - v_1) = m_2(v_2 - u_2). \tag{1} \]
Rearrange the energy equation in the same way, cancelling the factors $\tfrac12$:
\[ m_1(u_1^2 - v_1^2) = m_2(v_2^2 - u_2^2). \]
Factorising each difference of squares:
\[ m_1(u_1 - v_1)(u_1 + v_1) = m_2(v_2 - u_2)(v_2 + u_2). \tag{2} \]
Dividing (2) by (1) (both sides of (1) are non-zero during a genuine collision) gives
\[ u_1 + v_1 = v_2 + u_2 \quad\Longrightarrow\quad u_1 - u_2 = v_2 - v_1. \qquad \blacksquare \]

This is exactly \cle{Newton's experimental law of restitution} with coefficient of restitution $e = 1$.

\begin{methode}[Strategy for elastic collisions]
\begin{enumerate}
\item Draw a clear \textbf{before/after} diagram; fix a positive direction and write every velocity with its sign.
\item Write the \textbf{momentum equation}: $m_1u_1 + m_2u_2 = m_1v_1 + m_2v_2$.
\item Write the \textbf{restitution equation} $u_1 - u_2 = v_2 - v_1$ (much easier than the energy equation, and equivalent to it for elastic impacts).
\item Solve the two linear equations simultaneously for $v_1$ and $v_2$.
\item \textbf{Check}: verify that kinetic energy is indeed conserved and that the answer is physically sensible (the front body must end up faster than the rear body, otherwise they have passed through each other!).
\end{enumerate}
\end{methode}

\begin{exemplebox}[Special case 1: equal masses exchange velocities]
Set $m_1 = m_2 = m$. The momentum equation gives $u_1 + u_2 = v_1 + v_2$, and restitution gives $u_1 - u_2 = v_2 - v_1$. Adding and subtracting these yields
\[ v_1 = u_2, \qquad v_2 = u_1. \]
The two bodies simply \textbf{swap velocities}. If the target was at rest ($u_2 = 0$), the incoming body stops dead and the target moves off with the incoming velocity. This is exactly what you see when a moving trolley hits an identical stationary trolley on an air track.
\end{exemplebox}

\begin{exemplebox}[Special case 2: light body bouncing off a massive fixed body]
Let $m_2 \to \infty$ with $u_2 = 0$ (a wall, or the Earth). Solving the general equations (or taking limits) gives
\[ v_1 \approx -u_1, \qquad v_2 \approx 0. \]
The light body \textbf{rebounds with (nearly) the same speed, in the opposite direction}, while the massive body barely moves. A tennis ball dropped onto a hard floor bouncing back to nearly the same height is a familiar example.
\end{exemplebox}

\begin{popfigure}
\begin{tikzpicture}[scale=1.0, every node/.style={font=\small}]
% track
\draw[thick, popline] (-0.5,0) -- (12.5,0);
\node[popmuted] at (6,-0.35) {linear air track (friction negligible)};
% BEFORE
\node[font=\small\bfseries, popink] at (2.2,2.3) {Before (elastic, equal masses $m$)};
\draw[fill=popblueL, draw=popblue, thick] (0.5,0) rectangle (1.7,0.7);
\node at (1.1,0.35) {$m$};
\draw[->, very thick, popblue] (1.1,1.0) -- (2.6,1.0) node[midway, above]{$u$};
\draw[fill=popblueL, draw=popblue, thick] (3.4,0) rectangle (4.6,0.7);
\node at (4.0,0.35) {$m$};
\node at (4.0,1.05) {at rest};
% AFTER
\node[font=\small\bfseries, popink] at (9.4,2.3) {After};
\draw[fill=popgreenL, draw=popgreen, thick] (6.8,0) rectangle (8.0,0.7);
\node at (7.4,0.35) {$m$};
\node at (7.4,1.05) {stops};
\draw[fill=popgreenL, draw=popgreen, thick] (9.6,0) rectangle (10.8,0.7);
\node at (10.2,0.35) {$m$};
\draw[->, very thick, popgreen] (10.2,1.0) -- (11.7,1.0) node[midway, above]{$u$};
\end{tikzpicture}
\legende{Elastic collision of two equal-mass trolleys: the incoming trolley stops and the stationary one moves off with velocity $u$ — the trolleys exchange velocities.}
\end{popfigure}

\begin{exoresolu}[Head-on elastic collision]
A car of mass $1200\ \mathrm{kg}$ travelling at $20\ \mathrm{m\,s^{-1}}$ collides head-on elastically (an idealisation!) with a car of mass $800\ \mathrm{kg}$ travelling at $10\ \mathrm{m\,s^{-1}}$ in the opposite direction. Find their velocities after the collision.
\tcblower
\textbf{Solution.} Take the direction of the $1200\ \mathrm{kg}$ car as positive. Then
\[ m_1 = 1200,\ u_1 = 20; \qquad m_2 = 800,\ u_2 = -10. \]
\textbf{Momentum:}
\[ 1200(20) + 800(-10) = 1200v_1 + 800v_2 \;\Longrightarrow\; 16000 = 1200v_1 + 800v_2, \]
i.e. $3v_1 + 2v_2 = 40$. \quad (i)

\textbf{Restitution (elastic):} $u_1 - u_2 = v_2 - v_1$:
\[ 20 - (-10) = v_2 - v_1 \;\Longrightarrow\; v_2 - v_1 = 30. \quad (ii) \]
From (ii), $v_2 = v_1 + 30$; substituting in (i): $3v_1 + 2(v_1 + 30) = 40$, so $5v_1 = -20$, giving
\[ v_1 = -4\ \mathrm{m\,s^{-1}}, \qquad v_2 = 26\ \mathrm{m\,s^{-1}}. \]
The heavy car rebounds at $4\ \mathrm{m\,s^{-1}}$; the light car bounces back at $26\ \mathrm{m\,s^{-1}}$.

\textbf{Check (energy):} $\mathrm{KE}_{\text{before}} = \tfrac12(1200)(400) + \tfrac12(800)(100) = 280000\ \mathrm{J}$; $\mathrm{KE}_{\text{after}} = \tfrac12(1200)(16) + \tfrac12(800)(676) = 9600 + 270400 = 280000\ \mathrm{J}$. \checkmark
\end{exoresolu}

\courssub{Perfectly Inelastic Collisions: Bodies Stick Together}

When two vehicles crumple and lock together in a crash, or when a bullet embeds itself in a block of wood, the bodies move afterwards with a \textbf{common velocity} $v$. Only \emph{one} unknown remains, so momentum conservation \emph{alone} suffices:
\[ m_1 u_1 + m_2 u_2 = (m_1 + m_2)\,v \quad\Longrightarrow\quad v = \frac{m_1 u_1 + m_2 u_2}{m_1 + m_2}. \]
Kinetic energy is \emph{not} conserved: part of it is converted into heat, sound and the work of permanently deforming the bodies.

\begin{methode}[Computing the kinetic energy lost]
\begin{enumerate}
\item Use momentum conservation to find the final velocity (or velocities).
\item Compute $\mathrm{KE}_{\text{before}} = \tfrac12 m_1 u_1^2 + \tfrac12 m_2 u_2^2$ and $\mathrm{KE}_{\text{after}}$.
\item Energy lost: $\Delta \mathrm{KE} = \mathrm{KE}_{\text{before}} - \mathrm{KE}_{\text{after}}$.
\item Express the loss as a fraction or percentage of the initial KE: $\dfrac{\Delta \mathrm{KE}}{\mathrm{KE}_{\text{before}}} \times 100\%$.
\item State where the energy goes: heat, sound, permanent deformation.
\end{enumerate}
\end{methode}

\begin{popfigure}
\begin{tikzpicture}[scale=1.0, every node/.style={font=\small}]
\draw[thick, popline] (-0.5,0) -- (12.5,0);
% BEFORE
\node[font=\small\bfseries, popink] at (2.6,2.3) {Before};
\draw[fill=popblueL, draw=popblue, thick] (0.3,0) rectangle (1.8,0.7);
\node at (1.05,0.35) {$1200$ kg};
\draw[->, very thick, popblue] (1.05,1.0) -- (2.55,1.0) node[midway, above]{$15$};
\draw[fill=popblueL, draw=popblue, thick] (3.6,0) rectangle (4.8,0.7);
\node at (4.2,0.35) {$800$ kg};
\node at (4.2,1.05) {at rest};
% AFTER
\node[font=\small\bfseries, popink] at (9.3,2.3) {After (coalesced)};
\draw[fill=popgreenL, draw=popgreen, thick] (7.6,0) rectangle (9.1,0.7);
\draw[fill=popgreenL, draw=popgreen, thick] (9.1,0) rectangle (10.3,0.7);
\node at (8.35,0.35) {$1200$ kg};
\node at (9.7,0.35) {$800$ kg};
\draw[->, very thick, popgreen] (8.95,1.0) -- (10.75,1.0) node[midway, above]{$v = 9$};
\end{tikzpicture}
\legende{Perfectly inelastic crash: a $1200\ \mathrm{kg}$ vehicle at $15\ \mathrm{m\,s^{-1}}$ strikes a stationary $800\ \mathrm{kg}$ vehicle; they lock together. Momentum gives $v = \dfrac{1200 \times 15}{2000} = 9\ \mathrm{m\,s^{-1}}$.}
\end{popfigure}

\begin{exoresolu}[Kinetic energy lost in a coalescing crash]
Using the collision of the figure above: a $1200\ \mathrm{kg}$ car moving at $15\ \mathrm{m\,s^{-1}}$ hits a stationary $800\ \mathrm{kg}$ car and the two lock together. Find (a) the common velocity, (b) the kinetic energy lost, (c) the percentage of the initial KE lost.
\tcblower
\textbf{(a)} Momentum: $1200(15) + 800(0) = 2000\,v \Rightarrow v = 9\ \mathrm{m\,s^{-1}}$.

\textbf{(b)} $\mathrm{KE}_{\text{before}} = \tfrac12(1200)(15^2) = 135000\ \mathrm{J}$.
$\mathrm{KE}_{\text{after}} = \tfrac12(2000)(9^2) = 81000\ \mathrm{J}$.
$\Delta \mathrm{KE} = 135000 - 81000 = 54000\ \mathrm{J}$, converted into heat, sound and deformation of the bodywork.

\textbf{(c)} Percentage lost $= \dfrac{54000}{135000}\times 100\% = 40\%$.
\end{exoresolu}

\begin{exoresolu}[The ballistic pendulum idea]
A bullet of mass $20\ \mathrm{g}$ travelling horizontally at $350\ \mathrm{m\,s^{-1}}$ embeds itself in a wooden block of mass $1.98\ \mathrm{kg}$ resting on a smooth horizontal table. Find (a) the velocity of block and bullet just after impact, (b) the fraction of the bullet's kinetic energy that is lost.
\tcblower
\textbf{(a)} Momentum (bullet + block, perfectly inelastic):
\[ 0.020(350) = (0.020 + 1.98)\,v \;\Longrightarrow\; v = \frac{7.0}{2.00} = 3.5\ \mathrm{m\,s^{-1}}. \]
\textbf{(b)} $\mathrm{KE}_{\text{before}} = \tfrac12(0.020)(350^2) = 1225\ \mathrm{J}$; $\mathrm{KE}_{\text{after}} = \tfrac12(2.00)(3.5^2) = 12.25\ \mathrm{J}$.
Fraction lost $= \dfrac{1225 - 12.25}{1225} = 0.99$: \textbf{99\%} of the kinetic energy is dissipated as heat, sound and the work of tearing the wood. (In a real ballistic pendulum the block then swings upward; the common velocity found here is the starting point for that energy calculation.)
\end{exoresolu}

\courssub{Comparing Elastic and Perfectly Inelastic Collisions}

The bar chart below compares the kinetic energy before and after impact for the two archetypal collisions studied above, each starting with the same $1200\ \mathrm{kg}$ vehicle at $15\ \mathrm{m\,s^{-1}}$ striking a stationary $800\ \mathrm{kg}$ body.

\begin{popfigure}
\begin{center}
\begin{tikzpicture}
\begin{axis}[
    ybar, bar width=22pt,
    width=11cm, height=7cm,
    ymin=0, ymax=150,
    ylabel={Kinetic energy (kJ)},
    symbolic x coords={Before, After},
    xtick=data,
    legend style={at={(0.5,-0.18)}, anchor=north, legend columns=2},
    nodes near coords, every node near coord/.append style={font=\small},
    ymajorgrids, grid style={popline!40},
]
\addplot[fill=popblue, draw=popdark] coordinates {(Before,135) (After,135)};
\addplot[fill=poporange, draw=popdark] coordinates {(Before,135) (After,81)};
\legend{Elastic, Perfectly inelastic}
\end{axis}
\end{tikzpicture}
\end{center}
\legende{Kinetic energy before and after impact. Elastic collision: KE is fully conserved (135 kJ throughout). Perfectly inelastic collision: KE falls from 135 kJ to 81 kJ — 40\% is dissipated as heat, sound and deformation.}
\end{popfigure}

\courssub{Extension: The Coefficient of Restitution}

Real collisions lie \emph{between} the elastic and perfectly inelastic extremes. \cle{Newton's experimental law} states that for a direct impact, the speed of separation is a constant fraction $e$ of the speed of approach:
\[ v_2 - v_1 = e\,(u_1 - u_2), \qquad 0 \le e \le 1, \]
where $e$ is the \cle{coefficient of restitution}, a property of the materials in contact. The limiting cases are exactly the ones we have studied: $e = 1$ gives a perfectly \textbf{elastic} collision (speed of separation $=$ speed of approach), and $e = 0$ gives a \textbf{perfectly inelastic} collision ($v_2 = v_1$, the bodies move together). Together with momentum conservation, this law solves \emph{any} one-dimensional collision — a very useful tool for GCE problems involving bouncing balls.

\begin{attention}[Exam traps]
\begin{itemize}
\item Never write $\mathrm{KE}_{\text{before}} = \mathrm{KE}_{\text{after}}$ when the bodies stick together — use momentum only.
\item Velocities are \textbf{signed}: in a head-on collision one of the $u$'s must be negative.
\item In the restitution law, approach uses \emph{initial} velocities and separation uses \emph{final} velocities, in consistent order: $v_2 - v_1 = e(u_1 - u_2)$.
\item After the collision the rear body cannot be faster than the front body — if your algebra says otherwise, recheck your signs.
\end{itemize}
\end{attention}

\begin{aretenir}[Key points]
\begin{itemize}
\item Momentum is conserved in \textbf{every} collision; kinetic energy only in \textbf{elastic} ones.
\item Elastic 1-D collision: solve $m_1u_1 + m_2u_2 = m_1v_1 + m_2v_2$ together with $u_1 - u_2 = v_2 - v_1$.
\item Equal masses exchange velocities; a light body rebounds off a massive fixed body with nearly the same speed.
\item Perfectly inelastic: common velocity $v = \dfrac{m_1u_1 + m_2u_2}{m_1 + m_2}$; the lost KE becomes heat, sound and deformation.
\item General law: $v_2 - v_1 = e(u_1 - u_2)$, with $e = 1$ elastic, $e = 0$ perfectly inelastic.
\end{itemize}
\end{aretenir}

\begin{exemplebox}[Exam tip]
At the GCE, always \textbf{state explicitly which principle you are using and why}: ``By conservation of linear momentum (no external impulse during the impact)\ldots'' and, for elastic collisions, ``Since the collision is elastic, kinetic energy is conserved (equivalently, speed of approach $=$ speed of separation)\ldots''. Examiners award method marks for these statements even before any calculation.
\end{exemplebox}

\begin{exercice}[]
A trolley of mass $2.0\ \mathrm{kg}$ moving at $6.0\ \mathrm{m\,s^{-1}}$ on a smooth track collides with a stationary trolley of mass $4.0\ \mathrm{kg}$. (a) Assuming the collision is perfectly elastic, find both final velocities. (b) Assuming instead that the trolleys stick together, find their common velocity and the percentage of kinetic energy lost.
\end{exercice}

\begin{corrige}[Exercise 1]
(a) Momentum: $2(6) = 2v_1 + 4v_2 \Rightarrow v_1 + 2v_2 = 6$. Restitution: $v_2 - v_1 = 6$. Solving: $v_2 = 4\ \mathrm{m\,s^{-1}}$, $v_1 = -2\ \mathrm{m\,s^{-1}}$ (the light trolley rebounds).
(b) Coalescing: $12 = 6v \Rightarrow v = 2\ \mathrm{m\,s^{-1}}$. $\mathrm{KE}_{\text{before}} = \tfrac12(2)(36) = 36\ \mathrm{J}$; $\mathrm{KE}_{\text{after}} = \tfrac12(6)(4) = 12\ \mathrm{J}$. Loss $= \dfrac{24}{36}\times 100\% \approx 66.7\%$.
\end{corrige}

\coursec{Collisions and Momentum in Two Dimensions}

In the previous section we studied collisions along a straight line: two bodies approaching each other on the same road, or a ball striking another head-on. Real life, however, is rarely so tidy. Two cars meet at a crossroads in Buea at right angles; a snooker ball clips another and both roll off in different directions; an alpha particle is scattered sideways by a nucleus. In all these situations the velocities before and after the impact do \emph{not} lie on one line, and the one-dimensional equation $m_1u_1+m_2u_2=m_1v_1+m_2v_2$ is no longer enough. We must treat momentum as what it truly is: a \cle{vector}. This section develops the full vector machinery of two-dimensional collisions, which is one of the most frequently examined topics in the GCE Advanced Level mechanics papers.

\courssub{Momentum as a Plane Vector}

\begin{definition}[Momentum in component form]
The momentum of a body of mass $m$ moving with velocity $\vec{v}=v_x\,\vec{i}+v_y\,\vec{j}$ is the vector
\[
\vec{p}=m\vec{v}=mv_x\,\vec{i}+mv_y\,\vec{j}=\begin{pmatrix}mv_x\\ mv_y\end{pmatrix}.
\]
Its magnitude is $|\vec{p}|=m\sqrt{v_x^2+v_y^2}=mv$ and its direction is that of $\vec{v}$.
\end{definition}

The key observation is simple: since mass is a scalar, multiplying a velocity vector by $m$ does not change its direction. Everything we know about resolving vectors therefore carries over directly to momentum. If a ball of mass $0.5\ \mathrm{kg}$ moves at $10\ \mathrm{m\,s^{-1}}$ at $30^{\circ}$ above the horizontal, its momentum has components
\[
p_x=0.5\times 10\cos 30^{\circ}=2.5\sqrt{3}\ \mathrm{kg\,m\,s^{-1}},\qquad p_y=0.5\times 10\sin 30^{\circ}=2.5\ \mathrm{kg\,m\,s^{-1}}.
\]

\begin{propriete}[Conservation of momentum along each axis]
In any collision or explosion where no external impulse acts, the total momentum vector is conserved:
\[
\sum \vec{p}_{\text{before}}=\sum \vec{p}_{\text{after}}.
\]
A vector equation in the plane is equivalent to \emph{two independent scalar equations}: the total $x$-component of momentum is conserved, and the total $y$-component of momentum is conserved, separately:
\[
\sum p_x\ \text{constant},\qquad \sum p_y\ \text{constant}.
\]
\end{propriete}

This is the single most important idea of the section. Two perpendicular directions cannot ``exchange'' momentum: whatever momentum the system has along $\vec{i}$ stays along $\vec{i}$, and likewise for $\vec{j}$. In practice this gives us two equations, which is exactly what we need to find two unknowns (typically a speed and an angle).

\courssub{The Component Method for 2-D Collisions}

\begin{methode}[2-D collision routine]
\begin{enumerate}
\item Choose perpendicular axes $\vec{i}$ and $\vec{j}$ (usually along and perpendicular to the initial motion, or along the line of centres).
\item Resolve \emph{every} velocity, before and after, into components. Draw a clear ``before'' and ``after'' sketch with all angles marked.
\item Write conservation of momentum along $\vec{i}$: one scalar equation.
\item Write conservation of momentum along $\vec{j}$: a second scalar equation.
\item Solve the two simultaneous equations for the unknown components.
\item Recombine: speed $=\sqrt{v_x^2+v_y^2}$, direction $=\tan^{-1}\!\left|\dfrac{v_y}{v_x}\right|$ measured from the chosen axis. State the angle clearly.
\end{enumerate}
\end{methode}

\begin{exoresolu}[Two cars at a crossroads]
A car $A$ of mass $1200\ \mathrm{kg}$ travelling due east at $20\ \mathrm{m\,s^{-1}}$ collides at a junction with a car $B$ of mass $800\ \mathrm{kg}$ travelling due north at $15\ \mathrm{m\,s^{-1}}$. The cars lock together on impact. Find the velocity of the wreckage immediately after the collision.
\tcblower
\textbf{Solution.} Take $\vec{i}$ due east and $\vec{j}$ due north. Let the common velocity after impact be $\vec{v}=v_x\vec{i}+v_y\vec{j}$.

Along $\vec{i}$ (only $A$ has eastward momentum before):
\[
1200\times 20=(1200+800)v_x \;\Longrightarrow\; v_x=\frac{24000}{2000}=12\ \mathrm{m\,s^{-1}}.
\]
Along $\vec{j}$ (only $B$ has northward momentum before):
\[
800\times 15=2000\,v_y \;\Longrightarrow\; v_y=\frac{12000}{2000}=6\ \mathrm{m\,s^{-1}}.
\]
Speed of the wreckage:
\[
v=\sqrt{12^2+6^2}=\sqrt{180}=6\sqrt{5}\approx 13.4\ \mathrm{m\,s^{-1}}.
\]
Direction: $\tan\theta=\dfrac{v_y}{v_x}=\dfrac{6}{12}=0.5$, so $\theta\approx 26.6^{\circ}$ north of east, i.e.\ on a bearing of approximately $063.4^{\circ}$.
\end{exoresolu}

\begin{popfigure}
\begin{center}
\begin{tikzpicture}[scale=0.55,>=stealth]
% axes
\draw[->,gray] (-1,0) -- (11,0) node[right] {$\vec{i}$ (East)};
\draw[->,gray] (0,-1) -- (0,8) node[above] {$\vec{j}$ (North)};
% car A before
\draw[->,very thick,popblue] (-0.2,0.6) -- (4.5,0.6) node[midway,above,popblue] {$A:\ 20\ \mathrm{m\,s^{-1}}$};
% car B before
\draw[->,very thick,popgreen] (0.6,-0.2) -- (0.6,3.6) node[midway,right,popgreen] {$B:\ 15\ \mathrm{m\,s^{-1}}$};
% junction
\fill[popink] (0,0) circle (0.12);
% wreckage after
\draw[->,ultra thick,poporange] (5.5,1.2) -- (10.3,3.6) node[right,poporange] {wreckage $\vec{v}$};
\draw[dashed,popmuted] (5.5,1.2) -- (10.3,1.2) -- (10.3,3.6);
\node[popmuted] at (7.9,0.8) {$v_x=12$};
\node[popmuted,rotate=90] at (10.7,2.4) {$v_y=6$};
\draw[popdark] (7.3,1.2) arc (0:30:1.8);
\node[popdark] at (7.9,1.75) {$\theta$};
\end{tikzpicture}
\end{center}
\legende{Two cars collide at a junction and move off together; the velocity of the wreckage is found from its components.}
\end{popfigure}

\begin{attention}[Direction matters]
A body described as ``at rest'' or ``stopped'' after a collision may only have lost its momentum along \emph{one} axis. Always check \emph{both} components before concluding that a body is stationary. Likewise, never add momenta as scalars when the motions are not parallel: $\vec{p}_1+\vec{p}_2$ is a \emph{vector} sum.
\end{attention}

\courssub{Oblique Collisions of Smooth Spheres}

When two smooth spheres collide, the impulse between them acts along the \cle{line of centres} at the instant of impact (because the spheres are smooth, there is no tangential force). This physical fact leads to a beautiful simplification.

\begin{propriete}[Oblique impact of smooth spheres]
In a collision between two smooth spheres:
\begin{itemize}
\item the component of each sphere's velocity \emph{perpendicular} to the line of centres is unchanged by the impact;
\item along the line of centres, momentum is conserved (and Newton's law of restitution applies) exactly as in a one-dimensional collision.
\end{itemize}
\end{propriete}

A particularly clean case occurs when a sphere strikes an \emph{identical stationary} sphere obliquely. The moving sphere keeps its component of velocity perpendicular to the line of centres; along the line of centres the impact is head-on between equal masses, and for a perfectly elastic impact the spheres simply exchange their line-of-centres velocities. The moving sphere therefore stops dead along the line of centres and moves off purely perpendicular to it, while the target sphere moves off along the line of centres. The two balls separate at $90^{\circ}$ to each other — a result every snooker player knows instinctively.

\begin{popfigure}
\begin{center}
\begin{tikzpicture}[scale=1.0,>=stealth]
% target sphere
\draw[fill=popgreenL] (3,0) circle (0.5);
\node at (3,0) {$B$};
% moving sphere at impact position
\draw[fill=popblueL] (2.13,0.87) circle (0.5);
\node at (2.13,0.87) {$A$};
% line of centres
\draw[thick,popdark,dashed] (1.2,1.74) -- (4.2,-0.87);
\node[popdark] at (4.35,-1.05) {line of centres};
% perpendicular direction
\draw[thick,popmuted,dashed] (1.2,-0.3) -- (3.06,1.74);
\node[popmuted] at (1.0,-0.5) {perpendicular};
% incoming velocity
\draw[->,very thick,popblue] (0.2,2.6) -- (1.9,1.1) node[midway,above left,popblue] {$\vec{u}$};
% outgoing velocities
\draw[->,very thick,popgreen] (3.5,-0.25) -- (5.2,-1.1) node[right,popgreen] {$\vec{v}_B$};
\draw[->,very thick,popblue] (1.9,1.35) -- (1.15,2.85) node[above,popblue] {$\vec{v}_A$};
% right angle marker
\draw[popink] (3.28,0.28) -- (3.62,0.42) -- (3.48,0.76);
\end{tikzpicture}
\end{center}
\legende{Sphere $A$ strikes identical stationary sphere $B$ obliquely. The impulse acts along the line of centres; after an elastic impact the spheres move off at right angles.}
\end{popfigure}

\begin{exoresolu}[Snooker-style impact]
A ball of mass $0.16\ \mathrm{kg}$ moving at $4\ \mathrm{m\,s^{-1}}$ strikes an identical stationary ball obliquely. At impact, the line of centres makes an angle of $60^{\circ}$ with the initial direction of motion. The impact is perfectly elastic. Find the velocities of both balls after impact.
\tcblower
\textbf{Solution.} Resolve $\vec{u}$ along and perpendicular to the line of centres:
\[
u_{\parallel}=4\cos 60^{\circ}=2\ \mathrm{m\,s^{-1}},\qquad u_{\perp}=4\sin 60^{\circ}=2\sqrt{3}\ \mathrm{m\,s^{-1}}.
\]
Perpendicular component: unchanged, so the first ball keeps $2\sqrt{3}\ \mathrm{m\,s^{-1}}$ perpendicular to the line of centres.

Along the line of centres, equal masses, elastic impact, target at rest: the velocities are exchanged, so the first ball's line-of-centres component becomes $0$ and the target moves off at $2\ \mathrm{m\,s^{-1}}$ along the line of centres.

Hence the struck ball moves at $2\ \mathrm{m\,s^{-1}}$ along the line of centres (at $60^{\circ}$ to the original direction), and the cue ball moves at $2\sqrt{3}\approx 3.46\ \mathrm{m\,s^{-1}}$ perpendicular to the line of centres (at $30^{\circ}$ to the original direction, on the other side). Note $60^{\circ}+30^{\circ}=90^{\circ}$: the balls separate at right angles, as predicted.
\end{exoresolu}

\courssub{Momentum Triangles: the Geometric Method}

Because momentum is a vector, the equation
\[
\vec{p}_1+\vec{p}_2=\vec{p}_1^{\,\prime}+\vec{p}_2^{\,\prime}
\]
can be drawn as a \emph{closed polygon of vectors} — a triangle when only two momenta are involved on each side. Trigonometry (cosine rule and sine rule) then replaces simultaneous equations. This method is often faster and more elegant, especially when one body is initially at rest, so that $\vec{p}_{\text{before}}=\vec{p}_1^{\,\prime}+\vec{p}_2^{\,\prime}$ forms a single triangle.

\begin{methode}[Momentum-triangle method]
\begin{enumerate}
\item Write the vector equation $\sum\vec{p}_{\text{before}}=\sum\vec{p}_{\text{after}}$.
\item Draw the vectors head-to-tail so that the ``before'' side and the ``after'' side close the same triangle.
\item Mark all known magnitudes ($mu$ values) and angles.
\item Apply the cosine rule to find an unknown side, and the sine rule to find an unknown angle:
\[
a^2=b^2+c^2-2bc\cos A,\qquad \frac{a}{\sin A}=\frac{b}{\sin B}.
\]
\item Convert momentum magnitudes back to velocities by dividing by the relevant mass.
\end{enumerate}
\end{methode}

\begin{popfigure}
\begin{center}
\begin{tikzpicture}[scale=1.1,>=stealth]
\draw[->,ultra thick,popink] (0,0) -- (5,0) node[midway,below] {$m_1u$ (before)};
\draw[->,ultra thick,popblue] (0,0) -- (3.2,2.1) node[midway,above left,popblue] {$m_1v_1$};
\draw[->,ultra thick,poporange] (3.2,2.1) -- (5,0) node[midway,above right,poporange] {$m_2v_2$};
\draw[popdark] (1.3,0) arc (0:33.3:1.3);
\node[popdark] at (1.65,0.32) {$\theta$};
\draw[popdark] (4.2,0.62) arc (146.7:180:0.9);
\node[popdark] at (3.75,0.42) {$\phi$};
\end{tikzpicture}
\end{center}
\legende{Momentum triangle for a projectile striking a stationary target: $m_1\vec{u}=m_1\vec{v}_1+m_2\vec{v}_2$. The angles $\theta$ and $\phi$ are the deflections of the two bodies.}
\end{popfigure}

\begin{exoresolu}[Scattering of a particle by a stationary target]
A particle of mass $2\ \mathrm{kg}$ moving at $10\ \mathrm{m\,s^{-1}}$ strikes a stationary particle of mass $3\ \mathrm{kg}$. After the collision the $2\ \mathrm{kg}$ particle moves at $6\ \mathrm{m\,s^{-1}}$, deflected through $40^{\circ}$ from its original direction. Find the speed and direction of the $3\ \mathrm{kg}$ particle.
\tcblower
\textbf{Solution by components.} Take $\vec{i}$ along the initial direction. Momentum before: $(2\times 10,\ 0)=(20,0)$. Momentum of the scattered projectile after: $(2\times 6\cos 40^{\circ},\ 2\times 6\sin 40^{\circ})=(12\cos 40^{\circ},\ 12\sin 40^{\circ})$.

Let the target move with components $(V_x,V_y)$. Conservation:
\[
3V_x=20-12\cos 40^{\circ}=20-9.193=10.807 \;\Longrightarrow\; V_x=3.602,
\]
\[
3V_y=0-12\sin 40^{\circ}=-7.713 \;\Longrightarrow\; V_y=-2.571.
\]
Speed: $V=\sqrt{3.602^2+2.571^2}=\sqrt{19.58}\approx 4.43\ \mathrm{m\,s^{-1}}$.

Direction: $\tan\alpha=\dfrac{2.571}{3.602}=0.714$, so $\alpha\approx 35.5^{\circ}$ on the \emph{opposite} side of the original direction from the projectile (the negative sign of $V_y$).

\emph{Check with the triangle:} the three momenta $20$, $12$ and $3V$ close a triangle with $40^{\circ}$ between the sides $20$ and $12$:
\[
(3V)^2=20^2+12^2-2(20)(12)\cos 40^{\circ}=544-367.7=176.3,
\]
giving $3V=13.28$, $V\approx 4.43\ \mathrm{m\,s^{-1}}$ — confirming the component method.
\end{exoresolu}

\courssub{Is the Collision Elastic? Checking Kinetic Energy in 2-D}

Conservation of momentum alone cannot tell us whether a collision is elastic. The test is always the same as in one dimension: compare the total kinetic energy before and after, remembering that kinetic energy is a \emph{scalar} and uses the \emph{full speed}:
\[
E_k=\tfrac12 m|\vec{v}|^2=\tfrac12 m\left(v_x^2+v_y^2\right).
\]

For the crossroads collision above:
\[
E_k^{\text{before}}=\tfrac12(1200)(20^2)+\tfrac12(800)(15^2)=240000+90000=330000\ \mathrm{J},
\]
\[
E_k^{\text{after}}=\tfrac12(2000)(180)=180000\ \mathrm{J}.
\]
Since $180000<330000$, kinetic energy is lost: the collision is \cle{inelastic}, as expected for cars that lock together. For the snooker example, $E_k^{\text{before}}=\tfrac12(0.16)(16)=1.28\ \mathrm{J}$ and $E_k^{\text{after}}=\tfrac12(0.16)(4)+\tfrac12(0.16)(12)=0.32+0.96=1.28\ \mathrm{J}$: energy is conserved, confirming a perfectly elastic impact.

\begin{attention}[Squaring components correctly]
A frequent error is to compute kinetic energy as $\tfrac12 m v_x^2+\tfrac12 m v_y^2$ with the \emph{wrong} pairing of masses and components, or to forget to recombine components before squaring. Always find $|\vec{v}|^2=v_x^2+v_y^2$ for each body first. Also note: momentum can be conserved while kinetic energy is not — never use ``energy is conserved'' as a substitute for the momentum equations.
\end{attention}

\begin{exercice}[Target at rest, projectile deflected]
A ball of mass $0.5\ \mathrm{kg}$ moving at $8\ \mathrm{m\,s^{-1}}$ strikes a stationary ball of mass $1.5\ \mathrm{kg}$. After impact the lighter ball moves at right angles to its original direction at $4\ \mathrm{m\,s^{-1}}$. Find the velocity of the heavier ball, and determine whether the collision is elastic.
\end{exercice}

\begin{corrige}[Exercise 1]
Take $\vec{i}$ along the initial motion, $\vec{j}$ perpendicular. Before: $(0.5\times 8,\,0)=(4,0)$. After, light ball: $(0,\,0.5\times 4)=(0,2)$. Let the heavy ball have components $(V_x,V_y)$:
\[
1.5V_x=4-0 \Rightarrow V_x=\tfrac{8}{3};\qquad 1.5V_y=0-2 \Rightarrow V_y=-\tfrac{4}{3}.
\]
$V=\sqrt{\tfrac{64}{9}+\tfrac{16}{9}}=\tfrac{\sqrt{80}}{3}=\tfrac{4\sqrt{5}}{3}\approx 2.98\ \mathrm{m\,s^{-1}}$ at $\tan^{-1}\!\tfrac{1}{2}\approx 26.6^{\circ}$ below the original direction.

Energy: before $=\tfrac12(0.5)(64)=16\ \mathrm{J}$; after $=\tfrac12(0.5)(16)+\tfrac12(1.5)\!\left(\tfrac{80}{9}\right)=4+6.67=10.67\ \mathrm{J}<16\ \mathrm{J}$: inelastic.
\end{corrige}

\begin{aretenir}[Key points of the section]
\begin{itemize}
\item Momentum is a vector: $\vec{p}=m\vec{v}=(mv_x,\,mv_y)$.
\item In 2-D, momentum is conserved \emph{independently} along each of two perpendicular axes — this gives two scalar equations.
\item For smooth spheres, resolve along and perpendicular to the line of centres; the perpendicular components are unchanged.
\item The momentum triangle plus the cosine and sine rules offers a fast geometric alternative to components.
\item Always state the direction of a final velocity unambiguously, and check kinetic energy to classify the collision as elastic or inelastic.
\end{itemize}
\end{aretenir}

\begin{methode}[Exam tip: stating the angle of deflection]
GCE questions frequently ask for the \emph{angle of deflection} or the \emph{direction of motion} after impact. Never leave the answer as a bare $\tan^{-1}$ value: state precisely what the angle is measured from, e.g.\ ``$26.6^{\circ}$ north of east'', ``on a bearing of $063^{\circ}$'', or ``deflected through $40^{\circ}$ from its original path''. A correct magnitude with a vague direction often loses a mark.
\end{methode}

\begin{savaistu}[Particle scattering and the atom]
The same mathematics you have just used for snooker balls earned Rutherford's team its place in history: by measuring the angles through which alpha particles were deflected by a thin gold foil and applying conservation of momentum and energy in two dimensions, they deduced that almost all the mass of the atom is concentrated in a tiny nucleus. Vector momentum is truly a tool that reaches from the billiard table to the heart of matter.
\end{savaistu}

\coursec{Explosions, Rockets and Centre of Mass}

In the previous section we studied collisions, where two moving bodies meet and interact. In this section we turn the situation around: a single body \emph{splits apart} into fragments. Such events — explosions, nuclear decays, skaters pushing apart — are in a sense \cle{reverse collisions}, and the very same principle of conservation of linear momentum governs them. We then use momentum conservation to explain how rockets accelerate, and we finish with a powerful new concept, the \cle{centre of mass}, which lets us describe the motion of a whole system with a single point.

\courssub{Explosions as Reverse Collisions}

\textbf{Observation.}\par
A firework shell is fired vertically and, at the top of its flight, bursts into glowing fragments that fly off in all directions. Before the burst there was one momentum vector; after the burst there are many. Yet no \emph{external} force acted horizontally during the burst — the explosion is driven by \emph{internal} forces between the fragments. By Newton's third law these internal forces cancel in pairs, so the total momentum of the system cannot change.

\begin{propriete}[Momentum conservation in an explosion]
In an explosion, all the forces between the fragments are \cle{internal forces}. Hence the total linear momentum of the system is conserved:
\[
\vec{p}_{\text{before}}=\sum_{i}\vec{p}_{i,\text{after}} .
\]
In particular, if the body is initially at rest, the vector sum of the fragments' momenta after the explosion is zero:
\[
\sum_{i}m_{i}\vec{v}_{i}=\vec{0}.
\]
\end{propriete}

\begin{attention}[Kinetic energy is NOT conserved]
An explosion is the opposite of an elastic collision. Chemical (or nuclear, or stored elastic) energy is converted into kinetic energy, so the total kinetic energy \emph{increases}:
\[
E_{k,\text{after}}=E_{k,\text{before}}+E_{\text{released}}.
\]
Never apply the elastic-collision condition (conservation of kinetic energy, or relative velocity reversed) to an explosion. The only conservation laws you may use are momentum conservation and the energy balance above.
\end{attention}

\begin{methode}[Solving explosion problems]
\begin{enumerate}
\item Draw a ``before'' and an ``after'' diagram; choose axes.
\item Write the total momentum before the explosion (zero if the body was stationary).
\item Write the sum of the fragment momenta after the explosion, with unknown velocities as vectors.
\item Equate the two, component by component, and solve.
\item If the energy released is required, compute $E_{\text{released}}=E_{k,\text{after}}-E_{k,\text{before}}$.
\end{enumerate}
\end{methode}

\begin{exoresolu}[A stationary shell splits into two fragments]
A shell of mass $\SI{6}{kg}$, initially at rest, explodes into two fragments of masses $\SI{4}{kg}$ and $\SI{2}{kg}$. The $\SI{4}{kg}$ fragment moves off at $\SI{15}{m.s^{-1}}$. Find (a) the velocity of the other fragment, (b) the kinetic energy released by the explosion.
\tcblower
\textbf{(a)} Initial momentum is zero, so
\[
m_{1}\vec{v}_{1}+m_{2}\vec{v}_{2}=\vec{0}
\quad\Longrightarrow\quad
\vec{v}_{2}=-\frac{m_{1}}{m_{2}}\vec{v}_{1}=-\frac{4}{2}(15)=-\SI{30}{m.s^{-1}}.
\]
The $\SI{2}{kg}$ fragment moves at $\SI{30}{m.s^{-1}}$ in the \emph{opposite} direction.

\textbf{(b)} Kinetic energy before: $0$. After:
\[
E_{k}=\tfrac12(4)(15)^2+\tfrac12(2)(30)^2=450+900=\SI{1350}{J}.
\]
The explosion released $\SI{1350}{J}$ of chemical energy as kinetic energy.
\end{exoresolu}

\begin{exoresolu}[A shell exploding into three fragments]
A stationary shell explodes into three fragments $A$, $B$, $C$ of masses $\SI{1}{kg}$, $\SI{2}{kg}$ and $\SI{3}{kg}$. Fragment $A$ moves at $\SI{12}{m.s^{-1}}$ due east and $B$ moves at $\SI{9}{m.s^{-1}}$ due north. Find the velocity of $C$.
\tcblower
Take $\vec{i}$ east, $\vec{j}$ north. Conservation of momentum with zero initial momentum:
\[
1(12\vec{i})+2(9\vec{j})+3\vec{v}_{C}=\vec{0}
\quad\Longrightarrow\quad
\vec{v}_{C}=-4\vec{i}-6\vec{j}\ \mathrm{m\,s^{-1}}.
\]
Hence $|\vec{v}_{C}|=\sqrt{4^2+6^2}=\sqrt{52}\approx \SI{7.2}{m.s^{-1}}$, directed south-west, at an angle $\tan^{-1}(6/4)\approx 56.3^{\circ}$ south of west.
\end{exoresolu}

\begin{popfigure}
\begin{center}
\begin{tikzpicture}[scale=1.0,>=stealth]
\fill[popgold] (0,0) circle (0.28);
\node[below] at (0,-0.35) {\small shell at rest};
\draw[->,popink,very thick] (0.4,0.1) -- (2.4,0.7) node[right] {\small $m_1\vec{v}_1$};
\draw[->,poporange,very thick] (0.1,0.4) -- (-0.9,2.2) node[above] {\small $m_2\vec{v}_2$};
\draw[->,poppurple,very thick] (-0.35,-0.15) -- (-2.2,-1.3) node[below left] {\small $m_3\vec{v}_3$};
\node[align=center] at (0,-2.2) {\small $m_1\vec{v}_1+m_2\vec{v}_2+m_3\vec{v}_3=\vec{0}$};
\end{tikzpicture}
\end{center}
\legende{A stationary shell explodes into three fragments; the momentum vectors sum to zero.}
\end{popfigure}

\begin{exemplebox}[Everyday explosions]
\begin{itemize}
\item \textbf{Skaters pushing apart:} two skaters at rest on ice push each other. They glide off in opposite directions with momenta equal in magnitude: $m_1v_1=m_2v_2$. The lighter skater moves faster.
\item \textbf{Nuclear decay:} a stationary nucleus that emits an $\alpha$-particle recoils in the opposite direction so that total momentum stays zero — this is exactly how the recoil of a gun works too.
\item \textbf{Gun recoil:} when a bullet is fired forward, the gun recoils backward. The system (gun + bullet) has zero initial momentum, so $m_{\text{bullet}}v_{\text{bullet}} + m_{\text{gun}}v_{\text{gun}} = 0$.
\end{itemize}
\end{exemplebox}

\courssub{Rocket Motion}

\textbf{Observation.}\par
Release an inflated balloon and it darts forward as air rushes out backwards. Nothing pushes the balloon from outside: the expelled air carries momentum backwards, so the balloon must carry equal momentum forwards. A rocket works exactly the same way, burning fuel and ejecting hot exhaust gas at high speed.

\begin{propriete}[Rocket propulsion by momentum conservation]
A rocket ejects exhaust gas backwards at speed $v_{\text{rel}}$ relative to the rocket. The gas carries momentum backwards; by conservation of momentum the rocket gains momentum forwards. There is no need for anything external to ``push against'' — a rocket works best in the vacuum of space.
\end{propriete}

Consider a short time interval $\Delta t$ during which the rocket ejects a mass $\Delta m$ of gas at relative speed $v_{\text{rel}}$. The gas gains backward momentum $\Delta m\,v_{\text{rel}}$, so the rocket gains forward momentum $\Delta m\,v_{\text{rel}}$. The rate of change of the rocket's momentum — that is, the \cle{thrust} — is therefore:

\begin{propriete}[Thrust of a rocket]
\[
F=v_{\text{rel}}\,\frac{dm}{dt},
\]
where $\dfrac{dm}{dt}$ is the rate at which fuel mass is ejected and $v_{\text{rel}}$ is the exhaust speed relative to the rocket. The rocket is a \cle{variable-mass} system: as fuel burns, the remaining mass decreases, so the same thrust produces an ever-increasing acceleration $a=F/m(t)$.
\end{propriete}

\begin{popfigure}
\begin{center}
\begin{tikzpicture}[scale=1.0,>=stealth]
\draw[fill=popblueL,draw=popink,thick] (0,0) -- (0,1.2) -- (0.5,1.9) -- (1,1.2) -- (1,0) -- cycle;
\draw[fill=popmuted] (0.15,0) -- (0.35,-0.35) -- (0.65,-0.35) -- (0.85,0) -- cycle;
\draw[->,poporange,very thick] (0.5,-0.4) -- (0.5,-1.6) node[below] {\small exhaust, $v_{\text{rel}}$};
\draw[->,popgreen,very thick] (1.25,0.9) -- (1.25,2.3) node[above] {\small rocket momentum};
\draw[->,poppurple,thick] (-0.3,-0.9) -- (-1.5,-0.9) node[left] {\small gas momentum};
\node at (0.5,0.75) {\small $m$};
\end{tikzpicture}
\end{center}
\legende{A rocket ejects exhaust backwards; momentum conservation drives the rocket forwards.}
\end{popfigure}

\begin{exoresolu}[Computing a thrust]
A rocket engine ejects gas at $\SI{2500}{m.s^{-1}}$ relative to the rocket at a rate of $\SI{80}{kg.s^{-1}}$. Find the thrust.
\tcblower
\[
F=v_{\text{rel}}\frac{dm}{dt}=2500\times 80=\SI{2.0e5}{N}.
\]
The engine produces a thrust of $\SI{200}{kN}$.
\end{exoresolu}

\begin{exemplebox}[Multi-stage rockets]
The thrust equation shows why engineers build \cle{multi-stage rockets}. Most of a rocket's initial mass is fuel and the tanks and engines that hold it. Once a stage's fuel is spent, its empty structure is dead weight: carrying it further would waste thrust. By jettisoning the empty stage, the remaining rocket is much lighter, so the next stage's thrust produces a far greater acceleration. Real launch vehicles — such as the Ariane rockets launched from Kourou, or the SpaceX Falcon 9 — use two or three stages for exactly this reason.
\end{exemplebox}

\begin{savaistu}[Beyond the syllabus: the Tsiolkovsky rocket equation]
If you integrate the momentum balance as the rocket's mass falls from its initial value $m_0$ to a final value $m$, you obtain the famous \cle{rocket equation} of Konstantin Tsiolkovsky:
\[
v=u+v_{\text{rel}}\ln\!\left(\frac{m_0}{m}\right),
\]
where $u$ is the initial speed. The logarithm explains why reaching orbit is so hard: each extra bit of speed requires a \emph{multiplicative} increase in fuel. This derivation is beyond the GCE syllabus, but it is a classic in olympiad problems.
\end{savaistu}

\courssub{Centre of Mass of a System of Particles}

\textbf{Observation.}\par
Toss a spanner across a room. Each point of the spanner whirls in a complicated way, yet one special point glides along a smooth parabola, as if all the mass were concentrated there. That point is the \cle{centre of mass}.

\begin{definition}[Centre of mass]
The centre of mass of a system of particles of masses $m_1,m_2,\dots$ at positions $\vec{r}_1,\vec{r}_2,\dots$ is the mass-weighted average position
\[
\vec{r}_{CM}=\frac{\sum_i m_i\vec{r}_i}{\sum_i m_i}=\frac{\sum_i m_i\vec{r}_i}{M},
\qquad M=\sum_i m_i .
\]
In coordinates: $x_{CM}=\dfrac{\sum m_ix_i}{M}$, and similarly for $y_{CM}$ and $z_{CM}$.
\end{definition}

\begin{exoresolu}[Centre of mass of two particles]
A $\SI{2}{kg}$ mass sits at $x=0$ and a $\SI{6}{kg}$ mass at $x=\SI{8}{m}$ on a light rod. Locate the centre of mass.
\tcblower
\[
x_{CM}=\frac{2(0)+6(8)}{2+6}=\frac{48}{8}=\SI{6}{m}.
\]
The centre of mass is $\SI{6}{m}$ from the origin — closer to the heavier mass, as expected. Note the useful check: it divides the rod in the inverse ratio of the masses, $m_1:m_2=1:3$, so it lies $\tfrac{3}{4}$ of the way from the lighter mass.
\end{exoresolu}

\begin{exoresolu}[Centre of mass of a simple array]
Three particles are placed as follows: $\SI{1}{kg}$ at $(0,0)$, $\SI{2}{kg}$ at $(4,0)$ and $\SI{3}{kg}$ at $(0,6)$, coordinates in metres. Find the centre of mass.
\tcblower
Total mass $M=\SI{6}{kg}$. Then
\[
x_{CM}=\frac{1(0)+2(4)+3(0)}{6}=\frac{8}{6}=\frac{4}{3}\ \mathrm{m},
\qquad
y_{CM}=\frac{1(0)+2(0)+3(6)}{6}=\frac{18}{6}=\SI{3}{m}.
\]
The centre of mass is at $\left(\tfrac{4}{3},\,3\right)\ \mathrm{m}$.
\end{exoresolu}

\courssub{Motion of the Centre of Mass}

Differentiate the definition of $\vec{r}_{CM}$ with respect to time:
\[
M\vec{v}_{CM}=\sum_i m_i\vec{v}_i=\vec{P}_{\text{total}}.
\]

\begin{propriete}[Total momentum and the centre of mass]
The total momentum of a system equals the total mass times the velocity of the centre of mass:
\[
\vec{P}_{\text{total}}=M\vec{v}_{CM}.
\]
Consequently, for an \cle{isolated system} (no external force), $\vec{P}_{\text{total}}$ is constant, so the centre of mass moves with constant velocity — in a straight line at constant speed, or remains at rest if it started at rest. This holds no matter how violently the parts of the system interact internally.
\end{propriete}

This is the deep link between explosions and the centre of mass: \emph{fragments fly apart, but the centre of mass continues exactly as if nothing had happened.}

\begin{popfigure}
\begin{center}
\begin{tikzpicture}[scale=0.9,>=stealth]
\draw[->,gray] (0,0) -- (9.5,0) node[right] {\small $x$};
\draw[->,gray] (0,0) -- (0,4.2) node[above] {\small $y$};
\draw[dashed,popink,thick,domain=0.3:8.7,samples=40] plot (\x,{-0.09*(\x-4.5)*(\x-4.5)+3.2});
\fill[popgold] (1.5,1.62) circle (0.12);
\node[below] at (1.5,1.5) {\small shell};
\fill[poporange] (4.5,3.2) circle (0.10);
\node[above] at (4.5,3.3) {\small explosion};
\draw[->,poporange,thick] (4.5,3.2) -- (6.3,3.9);
\draw[->,poppurple,thick] (4.5,3.2) -- (3.1,2.6);
\draw[->,popgreen,thick] (4.5,3.2) -- (5.2,1.9);
\fill[popink] (7.5,1.62) circle (0.10);
\node[below] at (7.5,1.45) {\small CM};
\node[align=center,popmuted] at (4.5,0.55) {\small centre of mass follows\\ \small the original parabola};
\end{tikzpicture}
\end{center}
\legende{A projectile explodes in mid-flight; the fragments scatter, but the centre of mass continues along the original parabolic path (dashed).}
\end{popfigure}

\begin{exoresolu}[An exploding projectile]
A $\SI{12}{kg}$ projectile is moving horizontally at $\SI{20}{m.s^{-1}}$ when it explodes into two fragments of $\SI{8}{kg}$ and $\SI{4}{kg}$. Just after the explosion the $\SI{8}{kg}$ fragment is moving at $\SI{26}{m.s^{-1}}$ in the original direction. Find (a) the velocity of the $\SI{4}{kg}$ fragment, (b) the velocity of the centre of mass after the explosion.
\tcblower
\textbf{(a)} Momentum conservation along the line of motion:
\[
12(20)=8(26)+4v \quad\Longrightarrow\quad 240=208+4v \quad\Longrightarrow\quad v=\SI{8}{m.s^{-1}}.
\]
\textbf{(b)} The centre of mass is unaffected by the internal explosion:
\[
v_{CM}=\frac{8(26)+4(8)}{12}=\frac{240}{12}=\SI{20}{m.s^{-1}},
\]
exactly the original velocity, as the property guarantees.
\end{exoresolu}

\begin{attention}[Common mistakes with explosions and centre of mass]
\begin{itemize}
\item Applying elastic-collision rules (kinetic energy conserved) to an explosion — kinetic energy \emph{increases}.
\item Forgetting that momentum is a \textbf{vector}: fragments moving in different directions must be handled component by component.
\item Thinking the centre of mass ``jumps'' when a body explodes: internal forces cannot move the centre of mass; only an external force can.
\item In $F=v_{\text{rel}}\,\dfrac{dm}{dt}$, using the exhaust speed relative to the ground instead of relative to the rocket.
\end{itemize}
\end{attention}

\begin{exercice}[Explosion, rocket and centre of mass]
\begin{enumerate}
\item A stationary $\SI{5}{kg}$ body explodes into fragments of $\SI{2}{kg}$ and $\SI{3}{kg}$. The $\SI{3}{kg}$ fragment moves at $\SI{10}{m.s^{-1}}$. Find the speed of the other fragment and the kinetic energy released.
\item A rocket ejects $\SI{120}{kg}$ of gas per second at a relative speed of $\SI{1800}{m.s^{-1}}$. Calculate the thrust, and the initial acceleration if the rocket's initial mass is $\SI{30000}{kg}$ (ignore gravity).
\item Particles of $\SI{4}{kg}$ at $(1,2)$ and $\SI{6}{kg}$ at $(7,8)$ (metres) are joined by a light rod. Find the centre of mass.
\item A $\SI{10}{kg}$ shell moving at $\SI{30}{m.s^{-1}}$ horizontally explodes into two equal fragments. One fragment stops dead. Find the velocity of the other and the velocity of the centre of mass.
\end{enumerate}
\end{exercice}

\begin{corrige}[Exercise 1]
\begin{enumerate}
\item Momentum: $3(10)=2v \Rightarrow v=\SI{15}{m.s^{-1}}$ opposite to the $\SI{3}{kg}$ fragment. Energy released $=\tfrac12(3)(10^2)+\tfrac12(2)(15^2)=150+225=\SI{375}{J}$.
\item $F=v_{\text{rel}}\dfrac{dm}{dt}=1800\times120=\SI{2.16e5}{N}$; $a=\dfrac{F}{m}=\dfrac{2.16\times10^{5}}{3\times10^{4}}=\SI{7.2}{m.s^{-2}}$.
\item $x_{CM}=\dfrac{4(1)+6(7)}{10}=4.6$, $y_{CM}=\dfrac{4(2)+6(8)}{10}=5.6$; centre of mass at $(4.6,\,5.6)\ \mathrm{m}$.
\item Momentum: $10(30)=5(0)+5v \Rightarrow v=\SI{60}{m.s^{-1}}$ in the original direction. Centre of mass: $v_{CM}=\dfrac{5(0)+5(60)}{10}=\SI{30}{m.s^{-1}}$, unchanged by the explosion.
\end{enumerate}
\end{corrige}

\begin{aretenir}[Key points of this section]
\begin{itemize}
\item In an explosion, total momentum is conserved (zero if the body was at rest); kinetic energy \emph{increases} by the chemical or stored energy released.
\item Rocket thrust: $F=v_{\text{rel}}\,\dfrac{dm}{dt}$; rockets are variable-mass systems, and multi-staging sheds dead weight to boost acceleration.
\item Centre of mass: $\vec{r}_{CM}=\dfrac{\sum m_i\vec{r}_i}{M}$; total momentum $\vec{P}=M\vec{v}_{CM}$.
\item For an isolated system the centre of mass moves with constant velocity — fragments of an explosion scatter, but the centre of mass continues on the original path.
\end{itemize}
\end{aretenir}

\coursec{Real-Life Applications and Synthesis of Momentum as a Vector}

In the previous sections we saw how one single principle --- the conservation of linear momentum --- governs explosions, rocket propulsion and the motion of the centre of mass. We now close the chapter by returning to the everyday world: road accidents in Douala, a goalkeeper catching a shot, a fisherman pushing his pirogue on the Wouri. In every case the same two ideas do all the work:
\[
\vec{J}=\vec{F}_{\text{av}}\,\Delta t=\Delta\vec{p}
\qquad\text{and}\qquad
\sum \vec{p}_{\text{before}}=\sum \vec{p}_{\text{after}} \quad (\text{isolated system}).
\]
This section trains you to recognise \emph{which} idea applies, to use it quantitatively, and to present complete vector answers --- exactly what the GCE Advanced Level examiner rewards. It directly addresses Lessons 198 (Force-Time Graphs/Application in Safety) and 202 (Real-life Applications of Momentum as a Vector) of the official syllabus.

\courssub{Road Safety: Impact Forces from Momentum Change}

When a vehicle of mass $m$ moving at speed $u$ is brought to rest, its momentum changes by $\Delta p = mu$ (in magnitude). If the stopping takes a time $\Delta t$, the average force on the vehicle (and, by Newton's third law, on whatever stops it) is
\[
F_{\text{av}}=\frac{\Delta p}{\Delta t}=\frac{mu}{\Delta t}.
\]
Two lessons follow immediately. First, $F_{\text{av}}$ grows \emph{linearly with speed}: doubling your speed doubles the momentum to be destroyed, and since a collision with a rigid obstacle stops the car in roughly the same (very short) time, the impact force doubles too. Second, $F_{\text{av}}$ is \emph{inversely proportional to the stopping time}: anything that lengthens $\Delta t$ reduces the force.

\begin{propriete}[The core principle of all safety devices]
For a \emph{fixed} change of momentum $\Delta p$, the average force is inversely proportional to the contact time:
\[
F_{\text{av}}=\frac{\Delta p}{\Delta t}\propto \frac{1}{\Delta t}.
\]
Every safety device --- airbag, helmet, crumple zone, seatbelt, landing mat, boxing glove --- works by \textbf{increasing} $\Delta t$, never by changing $\Delta p$.
\end{propriete}

\begin{popfigure}
\begin{center}
\begin{tikzpicture}
\begin{axis}[
  width=0.85\linewidth, height=7cm,
  xlabel={Stopping time $\Delta t$ (s)},
  ylabel={Average force $F_{\text{av}}$ (kN)},
  xmin=0, xmax=0.5, ymin=0, ymax=300,
  domain=0.02:0.5, samples=100,
  grid=major, grid style={popline!40},
  axis lines=left, thick,
  clip=false]
\addplot[popblue, very thick] {15/x};
\addplot[poporange, dashed, thick] coordinates {(0.05,0) (0.05,300)};
\addplot[popgreen, dashed, thick] coordinates {(0.3,0) (0.3,300)};
\node[poporange, align=center, font=\small] at (axis cs:0.085,220) {rigid\\ impact};
\node[popgreen, align=center, font=\small] at (axis cs:0.36,220) {airbag};
\end{axis}
\end{tikzpicture}
\end{center}
\legende{Average impact force versus stopping time for a fixed momentum change $\Delta p = 15\ \mathrm{kN\,s}$: the hyperbola $F=\Delta p/\Delta t$. Stretching the contact time from $0.05$ s to $0.3$ s divides the force by six.}
\end{popfigure}

\begin{exoresolu}[Why speed limits save lives]
A car of mass $1000$ kg carrying a driver of mass $70$ kg travels at $90\ \mathrm{km\,h^{-1}}$ ($=25\ \mathrm{m\,s^{-1}}$). In a head-on collision the car stops in $0.10$ s. (a) Find the average force on the car. (b) The driver, restrained by a seatbelt and airbag, stops in $0.35$ s; find the average force on the driver. (c) Comment on the effect of doubling the speed.
\tcblower
\textbf{Solution.}
(a) $\Delta p_{\text{car}} = 1000\times 25 = 25\,000\ \mathrm{kg\,m\,s^{-1}}$, so
\[
F_{\text{car}}=\frac{25\,000}{0.10}=2.5\times10^{5}\ \mathrm{N}.
\]
(b) $\Delta p_{\text{driver}} = 70\times 25 = 1750\ \mathrm{kg\,m\,s^{-1}}$, so
\[
F_{\text{driver}}=\frac{1750}{0.35}=5.0\times10^{3}\ \mathrm{N},
\]
about seven times the driver's weight --- survivable when spread over the chest by the belt and airbag. Without the airbag the head might stop in $0.01$ s against the steering wheel, giving forces ten times larger, concentrated on the skull.
(c) At $180\ \mathrm{km\,h^{-1}}$ every $\Delta p$ doubles, so every force above doubles. Speed limits matter because force --- not just distance --- grows with speed.
\end{exoresolu}

\courssub{Safety Engineering and Sports: Impulse--Time Devices}

All the devices studied qualitatively earlier can now be treated quantitatively with $F_{\text{av}}=\Delta p/\Delta t$:

\begin{itemize}
\item \textbf{Airbags and crumple zones} stretch the stopping time of the passenger and of the car respectively.
\item \textbf{Helmets} contain a crushable foam liner: the head's momentum is destroyed over several milliseconds instead of a fraction of one.
\item \textbf{Seatbelts} stretch slightly and hold the passenger to the car so that the passenger decelerates over the car's (longer) stopping time rather than against the dashboard.
\item \textbf{Sports:} a cricketer or goalkeeper \emph{relaxes the hands} and draws them back while catching, lengthening $\Delta t$; a footballer \emph{follows through} the kick, lengthening the contact time so that a larger impulse $F\,\Delta t$ is transferred to the ball; high-jumpers land on thick mats; boxers wear padded gloves --- all impulse--time devices.
\end{itemize}

\begin{attention}[Following through versus catching]
Do not confuse the two directions of the same formula. \emph{Following through} increases $\Delta t$ to \textbf{increase the impulse delivered} ($\Delta p = F\Delta t$ grows). \emph{Catching with relaxed hands} increases $\Delta t$ to \textbf{decrease the force received} for a fixed $\Delta p$. State clearly which quantity is fixed in each situation.
\end{attention}

\courssub{Recoil and Propulsion in Daily Life}

Recoil is conservation of momentum in a system that starts at rest --- an ``explosion'' in the language of Section 5.

\begin{exemplebox}[Everyday recoil and propulsion]
\begin{itemize}
\item \textbf{Firing a gun:} bullet forward, gun backward; $m_b v_b = m_g v_g$ in magnitude, so the heavy gun recoils slowly.
\item \textbf{Jet engine and rocket:} hot gas ejected backward at high speed pushes the craft forward; thrust $F = v_e \left|\dfrac{dm}{dt}\right|$.
\item \textbf{Fire hose / water cannon:} the hose must be held firmly because ejecting water of mass flow $\dfrac{dm}{dt}$ at speed $v$ requires a backward reaction $v\dfrac{dm}{dt}$ on the hose.
\item \textbf{The pirogue trick:} a fisherman stranded on a still pirogue throws a heavy load backward; he and the pirogue drift forward, since the total momentum of the (nearly isolated) system must remain zero.
\end{itemize}
\end{exemplebox}

\courssub{Forensic Reconstruction of Collisions}

Traffic investigators use conservation of momentum \emph{backwards}: from the motion \emph{after} a crash (skid distances, final positions) they estimate speeds \emph{before} it. The impact itself is so brief that external impulses (friction during the few milliseconds of contact) are negligible, so momentum is conserved across the collision even when kinetic energy is not.

\begin{attention}[Real collisions are never perfectly isolated]
Friction, braking forces and road gradients act during every collision. Momentum conservation is still an excellent approximation \textbf{because the impact time is extremely short}: the external impulse $\vec{F}_{\text{ext}}\Delta t$ is negligible compared with the momenta involved. Always \emph{state this justification} in an exam solution --- it is worth marks.
\end{attention}

\begin{popfigure}
\begin{center}
\begin{tikzpicture}[scale=1.0, >=stealth]
% roads
\draw[fill=popblueL, draw=none] (-3.4,-0.5) rectangle (3.4,0.5);
\draw[fill=popblueL, draw=none] (-0.5,-3.0) rectangle (0.5,3.0);
\draw[dashed, popmuted] (-3.4,0) -- (-0.6,0); \draw[dashed, popmuted] (0.6,0) -- (3.4,0);
\draw[dashed, popmuted] (0,0.6) -- (0,3.0); \draw[dashed, popmuted] (0,-0.6) -- (0,-3.0);
% car A before
\draw[fill=poporange!70, draw=popdark] (-2.6,-0.28) rectangle (-1.4,0.28);
\node[font=\small] at (-2.0,0.55) {car A, $m_1$};
\draw[->, very thick, poporange] (-1.35,0) -- (-0.35,0) node[midway, above, font=\small]{$\vec{u}_1$};
% car B before
\draw[fill=popgreen!60, draw=popdark] (-0.28,-2.4) rectangle (0.28,-1.2);
\node[font=\small] at (1.35,-1.8) {car B, $m_2$};
\draw[->, very thick, popgreen] (0,-1.15) -- (0,-0.3) node[midway, right, font=\small]{$\vec{u}_2$};
% wreck after
\draw[fill=poppink!60, draw=popdark, rotate around={40:(0.9,0.9)}] (0.35,0.62) rectangle (1.45,1.18);
\node[font=\small, align=center] at (2.5,1.5) {wreck $(m_1+m_2)$};
\draw[->, very thick, poppurple] (1.0,1.0) -- (2.4,2.18) node[midway, above left, font=\small]{$\vec{V}$};
\draw[dashed, poppurple] (1.0,1.0) -- (2.4,1.0) node[midway, below, font=\small]{};
\draw[popmuted] (1.7,1.0) arc (0:40:0.7) node[midway, right, font=\small]{$\theta$};
\node[font=\small, align=left, text width=4.6cm] at (-2.2,2.3) {After impact the cars lock and skid together; measure $\theta$ and the skid length to find $\vec{V}$, then apply conservation of momentum.};
\end{tikzpicture}
\end{center}
\legende{Junction collision reconstruction: car A approaches eastward, car B northward; the locked wreck moves off at angle $\theta$. Momentum is conserved component by component.}
\end{popfigure}

\begin{exoresolu}[Estimating speeds before a junction collision]
Car A of mass $1200$ kg travelling east and car B of mass $1000$ kg travelling north collide at a junction and lock together. The wreck skids off at $\theta = 40^{\circ}$ north of east, and from the skid marks its speed just after impact is estimated at $V = 12\ \mathrm{m\,s^{-1}}$. Estimate the speed of each car before impact.
\tcblower
\textbf{Solution.} Take axes $x$ east, $y$ north. Momentum is conserved in each component (impact of very short duration).
\[
x:\quad 1200\,u_1 = (1200+1000)(12\cos 40^{\circ})
\;\Rightarrow\;
u_1=\frac{2200\times 12\times 0.766}{1200}\approx 16.9\ \mathrm{m\,s^{-1}}.
\]
\[
y:\quad 1000\,u_2 = 2200\times 12\sin 40^{\circ}
\;\Rightarrow\;
u_2=\frac{2200\times 12\times 0.643}{1000}\approx 17.0\ \mathrm{m\,s^{-1}}.
\]
Both cars were travelling at about $61\ \mathrm{km\,h^{-1}}$ --- useful evidence if the limit was $50\ \mathrm{km\,h^{-1}}$.
\end{exoresolu}

\courssub{Choosing the Right Model: A Modelling Checklist}

The hardest part of a momentum problem is rarely the algebra --- it is recognising the \emph{type} of situation.

\begin{methode}[Modelling checklist for any momentum problem]
\begin{enumerate}
\item \textbf{Define the system.} Which bodies are included? Is it (approximately) isolated during the event?
\item \textbf{Classify the interaction.} Collision (bodies meet), explosion/recoil (one body splits or pushes off), or continuous impulse (force acting over a measurable time, e.g.\ hose, rocket, airbag)?
\item \textbf{Choose the tool.} Isolated collision/explosion $\Rightarrow$ conservation of momentum (component by component in 2D). Force and time given $\Rightarrow$ impulse $\vec{J}=\vec{F}\Delta t=\Delta\vec{p}$. Perfectly elastic $\Rightarrow$ add conservation of kinetic energy (or $e=1$).
\item \textbf{List knowns and unknowns,} fix a positive direction (or axes), and assign signs \emph{before} substituting numbers.
\item \textbf{Solve, then check:} units, signs, and physical plausibility (e.g.\ a light object can rebound faster than it arrived, but kinetic energy must not increase in a collision).
\end{enumerate}
\end{methode}

\begin{popfigure}
\begin{center}
\begin{tikzpicture}[
  font=\small,
  box/.style={draw=popblue, thick, rounded corners, fill=popblueL, align=center, text width=3.6cm, minimum height=0.9cm},
  dec/.style={draw=poporange, thick, fill=poporange!15, align=center, text width=3.2cm, minimum height=0.9cm},
  arr/.style={->, thick, popdark}]
\node[dec] (start) at (0,0) {Is a force and a time interval given?};
\node[box, align=center] (imp) at (-5.2,-2.2){Impulse method:\\ $J=F\,\Delta t=\Delta p$\\ (airbag, hose, catch)};
\node[dec] (iso) at (0,-2.2) {System isolated during the event?};
\node[box, align=center] (noiso) at (5.2,-2.2){Use Newton's 2nd law / energy;\\ momentum not conserved};
\node[dec] (type) at (0,-4.4) {Bodies join or separate?};
\node[box, align=center] (coll) at (-5.2,-6.6){Collision:\\ $m_1\vec{u}_1+m_2\vec{u}_2=m_1\vec{v}_1+m_2\vec{v}_2$\\ add $e$ or KE if elastic};
\node[box, align=center] (expl) at (0,-6.6){Explosion / recoil:\\ total $\vec{p}$ constant\\ (often $\vec{0}$ initially)};
\node[box, align=center] (cont) at (5.2,-6.6){Continuous mass flow:\\ $F=v_e\,\dfrac{dm}{dt}$\\ (rocket, jet, hose)};
\draw[arr] (start) -- node[above left, font=\scriptsize]{yes} (imp);
\draw[arr] (start) -- node[left, font=\scriptsize]{no} (iso);
\draw[arr] (iso) -- node[above right, font=\scriptsize]{no} (noiso);
\draw[arr] (iso) -- node[left, font=\scriptsize]{yes} (type);
\draw[arr] (type) -- node[above left, font=\scriptsize]{meet} (coll);
\draw[arr] (type) -- node[left, font=\scriptsize]{split} (expl);
\draw[arr] (type) -- node[above right, font=\scriptsize]{steady flow} (cont);
\end{tikzpicture}
\end{center}
\legende{Decision tree for selecting the correct momentum method.}
\end{popfigure}

\courssub{Exam Strategy and Synthesis}

\begin{attention}[Presenting vector answers]
A momentum or velocity answer in two dimensions is \emph{never} complete without \textbf{both magnitude and direction}, and units. Write, for example, ``$|\vec{p}| = 26\ \mathrm{kg\,m\,s^{-1}}$ at $23^{\circ}$ north of east''. Remember: $1\ \mathrm{N\,s}=1\ \mathrm{kg\,m\,s^{-1}}$. In one dimension, a negative sign \emph{is} the direction --- state what it means (``the ball rebounds'').
\end{attention}

\begin{exercice}[Synthesis: the runaway truck]
A truck of mass $3000$ kg moving at $20\ \mathrm{m\,s^{-1}}$ strikes a stationary car of mass $1000$ kg; the two lock together. (a) Find their common speed just after impact. (b) They then skid to rest in $4.0$ s under a constant friction force; find this force using impulse. (c) Show that kinetic energy is not conserved and find the fraction lost.
\end{exercice}

\begin{corrige}[Exercise 1]
(a) Conservation of momentum: $3000\times 20 = 4000\,V \Rightarrow V = 15\ \mathrm{m\,s^{-1}}$.
(b) Impulse: $F\Delta t = \Delta p = 4000\times 15 = 60\,000\ \mathrm{kg\,m\,s^{-1}}$, so $F = \dfrac{60\,000}{4.0} = 1.5\times10^{4}\ \mathrm{N}$.
(c) KE before $=\tfrac12(3000)(20^2)=6.0\times10^{5}$ J; KE after $=\tfrac12(4000)(15^2)=4.5\times10^{5}$ J. Fraction lost $=\dfrac{1.5\times10^{5}}{6.0\times10^{5}}=0.25$, i.e.\ $25\%$, converted to heat, sound and deformation.
\end{corrige}

\begin{aretenir}[One principle unifies the whole chapter]
\begin{itemize}
\item \textbf{Impulse:} $\vec{J}=\vec{F}_{\text{av}}\Delta t=\Delta\vec{p}$ --- links force, time and momentum; the basis of all safety devices ($F_{\text{av}}\propto 1/\Delta t$).
\item \textbf{Conservation:} for an isolated system, $\sum\vec{p}$ is constant --- the \emph{same} law solves collisions (elastic and inelastic, 1D and 2D), explosions, recoil and forensic reconstructions.
\item \textbf{Continuous flow:} thrust $F=v_e\,\dfrac{dm}{dt}$ for rockets, jets and hoses.
\item \textbf{Centre of mass:} the centre of mass of an isolated system moves with constant velocity --- momentum conservation seen ``from above''.
\item \textbf{Always:} define the system, justify isolation by the short impact time, fix axes and signs, and conclude with magnitude \emph{and} direction, with units ($\mathrm{kg\,m\,s^{-1}}$ or $\mathrm{N\,s}$).
\end{itemize}
\end{aretenir}

\begin{savaistu}[Momentum around us]
Crash-test programmes used by vehicle manufacturers worldwide are essentially large-scale momentum experiments: dummies instrumented with accelerometers measure $F\,\Delta t$ during impacts, and every improvement --- collapsible steering columns, side-impact bars, helmet standards --- is an application of $F_{\text{av}}=\Delta p/\Delta t$. The same equation that a goalkeeper uses instinctively when softening a catch guides the engineers designing the vehicles on the Yaound\'e--Douala highway. One principle, countless lives saved.
\end{savaistu}

\newpage

\coursec{Integration activity}

\coursec{Investigative Activity: The Mile 4 Junction Crash}

\courssub{Aims of the Activity}

By the end of this group investigation, you should be able to:

\begin{itemize}
\item apply the \cle{principle of conservation of linear momentum} in \textbf{two dimensions}, working component by component;
\item compute the \cle{kinetic energy} before and after a collision and use the result to \textbf{classify} the collision as elastic, inelastic or perfectly inelastic;
\item use the \cle{impulse--momentum relation} $F_{\text{av}}\,\Delta t = \Delta p$ to estimate average forces and to explain, quantitatively, how \textbf{safety devices} (seat belts, airbags, crumple zones) reduce injuries;
\item communicate a complete scientific argument: set up axes, state assumptions, present calculations with units, and defend conclusions in front of the class;
\item connect classroom mechanics to a genuine \textbf{road-safety concern} in Cameroon.
\end{itemize}

\begin{attention}[Ground Rules]
Work in groups of four. Every numerical answer must carry its \textbf{unit} and be justified by a \textbf{formula written in symbols first}. Take $g = 9.81\ \mathrm{m\,s^{-2}}$ where needed. Neglect friction during the impact itself (the collision time is very short, so external impulses during the impact are negligible).
\end{attention}

\courssub{Situation}

\textbf{The Mile 4 Junction Crash (Limbe).}\par
At 7:40 on a Monday morning, at a cross junction near Mile 4, a \emph{taxi} (a saloon car) of mass $m_1 = 1200\ \mathrm{kg}$ was travelling \textbf{due East} at $u_1 = 20\ \mathrm{m\,s^{-1}}$ (about $72\ \mathrm{km\,h^{-1}}$). At the same moment, a \emph{pickup truck} of mass $m_2 = 1800\ \mathrm{kg}$ was travelling \textbf{due North} at $u_2 = 15\ \mathrm{m\,s^{-1}}$ (about $54\ \mathrm{km\,h^{-1}}$). Neither driver stopped at the junction. The vehicles collided at right angles, \textbf{locked together} at impact, and the combined wreck skidded off along the tarmac.

A $70\ \mathrm{kg}$ passenger was sitting in the front seat of the taxi, moving with the taxi at $20\ \mathrm{m\,s^{-1}}$ before the crash. Investigators consider two scenarios for this passenger during the collision:

\begin{itemize}
\item \textbf{Scenario A (no airbag, no effective restraint):} the passenger is brought to rest against the dashboard in about $\Delta t_A = 0.02\ \mathrm{s}$;
\item \textbf{Scenario B (seat belt and airbag):} the passenger is brought to rest over about $\Delta t_B = 0.15\ \mathrm{s}$.
\end{itemize}

The diagram below summarises the situation. Take the $x$-axis pointing East and the $y$-axis pointing North.

\begin{popfigure}
\begin{tikzpicture}[scale=0.9, >=stealth]
% axes
\draw[->, gray] (-0.5,0) -- (7.5,0) node[below left] {\small East ($x$)};
\draw[->, gray] (0,-0.5) -- (0,4.5) node[below left] {\small North ($y$)};
% taxi
\draw[fill=popblueL, draw=popblue] (1.2,-0.4) rectangle (2.6,0.4);
\node at (1.9,0) {\small Taxi};
\draw[->, very thick, popblue] (2.6,0) -- (4.4,0) node[midway, above] {\small $\vec{u}_1 = 20\ \mathrm{m\,s^{-1}}$};
% pickup
\draw[fill=popgreenL, draw=popgreen] (5.6,0.9) rectangle (6.4,2.3);
\node at (6,1.6) {\small Pickup};
\draw[->, very thick, popgreen] (6,2.3) -- (6,3.9) node[midway, right] {\small $\vec{u}_2 = 15\ \mathrm{m\,s^{-1}}$};
% impact point
\fill[poporange] (6,0) circle (0.09);
\node[below right, popdark] at (6,0) {\small Impact};
% wreck velocity
\draw[->, very thick, poporange] (6,0) -- (7.3,1.3) node[right] {\small $\vec{v}$ (wreck)};
\draw[dashed, popmuted] (6,0) -- (7.3,0);
\draw[popdark] (6.7,0) arc (0:48:0.7) node[midway, right, xshift=1pt] {\small $\theta$};
\end{tikzpicture}
\legende{The taxi moves East, the pickup moves North; after locking together, the wreck moves with velocity $\vec{v}$ at angle $\theta$ North of East.}
\end{popfigure}

You are the team of \textbf{scientific investigators} attached to the case. Your report will be forwarded to the Divisional Delegation of Transport and used in a road-safety campaign.

\courssub{Tasks}

\textbf{Part A --- Velocity of the Wreck (Conservation of Momentum in 2D)}\par
\begin{enumerate}
\item Write down the momentum vector of each vehicle \textbf{just before} the collision, giving the $x$-component (East) and $y$-component (North) separately, in $\mathrm{kg\,m\,s^{-1}}$.
\item State the principle of conservation of linear momentum and explain briefly why it applies here \emph{component by component}, even though the collision is violent.
\item Hence find the components $v_x$ and $v_y$ of the velocity of the $3000\ \mathrm{kg}$ wreck just after impact.
\item Deduce the \textbf{magnitude} $v$ of this velocity and the \textbf{angle} $\theta$ it makes with the East direction. Give $\theta$ to one decimal place.
\end{enumerate}

\textbf{Part B --- Energy Analysis (Nature of the Collision)}\par
\begin{enumerate}
\setcounter{enumi}{4}
\item Calculate the total kinetic energy of the two vehicles \textbf{before} the collision.
\item Calculate the kinetic energy of the wreck \textbf{after} the collision.
\item Determine the kinetic energy \textbf{lost} during the impact. Into what forms has this energy been transformed? State, with justification, whether the collision is elastic, inelastic or perfectly inelastic.
\end{enumerate}

\textbf{Part C --- Impulse and the Passenger (Force--Time Reasoning)}\par
\begin{enumerate}
\setcounter{enumi}{7}
\item Compute the change of momentum $\Delta p$ of the $70\ \mathrm{kg}$ passenger, initially moving at $20\ \mathrm{m\,s^{-1}}$, who is brought to rest during the crash.
\item Using $F_{\text{av}} = \dfrac{\Delta p}{\Delta t}$, estimate the average force exerted on the passenger in \textbf{Scenario A} ($\Delta t_A = 0.02\ \mathrm{s}$) and in \textbf{Scenario B} ($\Delta t_B = 0.15\ \mathrm{s}$).
\item Express each force as a multiple of the passenger's weight ($W = mg$). Comment on the comparison.
\item Sketch, on the same axes, two force--time graphs (one per scenario) showing the same area under each curve. What does this area represent?
\end{enumerate}

\textbf{Part D --- Communication}\par
\begin{enumerate}
\setcounter{enumi}{11}
\item In at most ten lines, write a \textbf{safety recommendation} for the road-safety campaign, addressed to drivers and passengers in the South-West Region. Your text must use at least three of these words correctly: \emph{momentum, impulse, stopping time, conservation, kinetic energy}.
\end{enumerate}

Each group then presents Parts A--C on the board; the class challenges the assumptions and checks the units.

\courssub{Model Answers and Examiner's Commentary}

\textbf{Part A}\par
\textbf{1.} Momenta before impact (East $=x$, North $=y$):
\begin{align*}
\text{Taxi:}\quad & p_{1x} = m_1 u_1 = 1200 \times 20 = 24\,000\ \mathrm{kg\,m\,s^{-1}}, \qquad p_{1y} = 0,\\
\text{Pickup:}\quad & p_{2x} = 0, \qquad p_{2y} = m_2 u_2 = 1800 \times 15 = 27\,000\ \mathrm{kg\,m\,s^{-1}}.
\end{align*}

\textbf{2.} \emph{Principle:} If the resultant external force on a system is zero (or negligible during the event), the total linear momentum of the system remains constant. During the very short impact, the mutual collision forces are enormous compared with friction from the road, so external impulse is negligible. Since momentum is a \textbf{vector}, the principle holds separately along $x$ and along $y$.

\textbf{3.} Let $\vec{v} = (v_x, v_y)$ be the velocity of the combined wreck of mass $M = 1200 + 1800 = 3000\ \mathrm{kg}$:
\begin{align*}
x\text{-axis:}\quad & 24\,000 = 3000\,v_x \;\Rightarrow\; v_x = 8\ \mathrm{m\,s^{-1}},\\
y\text{-axis:}\quad & 27\,000 = 3000\,v_y \;\Rightarrow\; v_y = 9\ \mathrm{m\,s^{-1}}.
\end{align*}

\textbf{4.} Magnitude and direction:
\[
v = \sqrt{v_x^2 + v_y^2} = \sqrt{8^2 + 9^2} = \sqrt{145} \approx 12.0\ \mathrm{m\,s^{-1}},
\qquad
\tan\theta = \frac{v_y}{v_x} = \frac{9}{8} \;\Rightarrow\; \theta \approx 48.4^\circ \text{ North of East}.
\]

\textbf{Part B}\par
\textbf{5.} Kinetic energy before:
\[
E_{k,\text{before}} = \tfrac{1}{2}(1200)(20)^2 + \tfrac{1}{2}(1800)(15)^2 = 240\,000 + 202\,500 = 442\,500\ \mathrm{J}.
\]

\textbf{6.} Kinetic energy after:
\[
E_{k,\text{after}} = \tfrac{1}{2}(3000)\left(\sqrt{145}\right)^2 = \tfrac{1}{2}(3000)(145) = 217\,500\ \mathrm{J}.
\]

\textbf{7.} Energy lost:
\[
\Delta E_k = 442\,500 - 217\,500 = 225\,000\ \mathrm{J} = 225\ \mathrm{kJ}.
\]
This energy is transformed into \textbf{deformation} of the bodywork (permanent bending of metal), \textbf{sound}, and \textbf{heat}. Since kinetic energy is not conserved \emph{and} the two vehicles move together after impact, the collision is \textbf{perfectly inelastic}. Note that about $51\%$ of the initial kinetic energy is dissipated --- a perfectly inelastic collision is the most dissipative type consistent with momentum conservation.

\textbf{Part C}\par
\textbf{8.} Change of momentum of the passenger (brought from $20\ \mathrm{m\,s^{-1}}$ to rest):
\[
\Delta p = m\,\Delta u = 70 \times 20 = 1400\ \mathrm{kg\,m\,s^{-1}}.
\]

\textbf{9.} Average forces:
\[
F_A = \frac{1400}{0.02} = 70\,000\ \mathrm{N}, \qquad
F_B = \frac{1400}{0.15} \approx 9\,333\ \mathrm{N}.
\]

\textbf{10.} The passenger's weight is $W = mg = 70 \times 9.81 \approx 687\ \mathrm{N}$. Hence
\[
\frac{F_A}{W} \approx 102, \qquad \frac{F_B}{W} \approx 13.6.
\]
Without restraint, the passenger experiences a force about \textbf{one hundred times his own weight}, which is almost certainly fatal. With seat belt and airbag, the same change of momentum is spread over a stopping time $7.5$ times longer, so the average force is divided by $7.5$: the impulse $F\,\Delta t = \Delta p$ is fixed by the momentum change, only its distribution in time changes.

\textbf{11.} Force--time sketch:

\begin{popfigure}
\begin{tikzpicture}[scale=1.0, >=stealth]
\draw[->] (0,0) -- (5.2,0) node[below left] {\small $t$ (s)};
\draw[->] (0,0) -- (0,3.4) node[below left] {\small $F$ (N)};
% Scenario A: tall narrow pulse (rectangular for average force)
\draw[very thick, poporange] (0,0) -- (0.15,0) -- (0.15,3.0) -- (0.35,3.0) -- (0.35,0);
\node[poporange, align=center] at (0.6,3.2) {\small A: $\Delta t = 0.02$ s};
% Scenario B: low wide pulse
\draw[very thick, popblue] (0,0) -- (0.3,0) -- (0.3,0.5) -- (2.6,0.5) -- (2.6,0);
\node[popblue, align=center] at (2.9,0.7) {\small B: $\Delta t = 0.15$ s};
% Shaded areas (impulse)
\fill[poppink, opacity=0.35] (0.15,0) rectangle (0.35,3.0);
\fill[popblueL, opacity=0.5] (0.3,0) rectangle (2.6,0.5);
\node[align=center, popmuted] at (3.9,2.4) {\small Equal areas:\\ \small same $\Delta p = 1400\ \mathrm{N\,s}$};
\end{tikzpicture}
\legende{Both force--time pulses enclose the same area $\int F\,dt = \Delta p = 1400\ \mathrm{kg\,m\,s^{-1}}$; the airbag lowers the peak force by lengthening the stopping time.}
\end{popfigure}

The area under each curve is the \textbf{impulse}, equal to the change of momentum of the passenger: identical in both scenarios.

\textbf{Part D (Sample Recommendation)}\par
\emph{``Every moving vehicle carries momentum, and in a crash that momentum must be removed by an impulse. Physics cannot reduce the impulse --- your change of momentum is fixed by your speed --- but it can stretch the stopping time. Seat belts and airbags increase the stopping time of your body, dividing the average force by a large factor and saving lives. Drivers: reduce speed, because both momentum and kinetic energy grow with speed. Passengers: fasten your seat belt on every journey, even short ones. Let conservation of momentum work for you, not against you.''}

\begin{aretenir}[What This Activity Consolidates]
\begin{itemize}
\item Momentum is a \textbf{vector}: in two dimensions, conserve it \textbf{component by component}.
\item Locking together $\Rightarrow$ \textbf{perfectly inelastic}: momentum conserved, kinetic energy not.
\item $F_{\text{av}}\,\Delta t = \Delta p$: for a fixed $\Delta p$, a \textbf{longer stopping time} means a \textbf{smaller average force} --- the principle behind every safety device.
\end{itemize}
\end{aretenir}

\newpage

\coursec{Methods and worked examples}

\courssub{Method 1: Computing Momentum and Impulse as Vectors}

\begin{definition}[Momentum and impulse]
The \cle{momentum} of a body of mass $m$ moving with velocity $\vec{v}$ is the vector
\[ \vec{p}=m\vec{v}, \qquad \text{unit: } \mathrm{kg\,m\,s^{-1}}\ (\text{equivalently N\,s}). \]
The \cle{impulse} of a force $\vec{F}$ acting from time $t_1$ to $t_2$ is
\[ \vec{J}=\vec{F}\,\Delta t\ \ (\text{constant force}), \qquad \vec{J}=\int_{t_1}^{t_2}\vec{F}\,dt\ \ (\text{variable force}). \]
\end{definition}

\begin{propriete}[Impulse--momentum theorem]
The impulse of the resultant force acting on a body equals its change of momentum:
\[ \vec{J}=\Delta\vec{p}=m\vec{v}-m\vec{u}. \]
This follows directly from Newton's second law $\vec{F}=\dfrac{d\vec{p}}{dt}$ integrated over time. (The derivation from $\vec{F}=m\vec{a}$ for constant mass is examinable.)
\end{propriete}

\begin{methode}[Momentum and impulse of a moving body]
\begin{enumerate}
\item Identify the mass $m$ (in kg) and write the velocity as a \cle{vector} $\vec{v}$ (choose unit vectors $\vec{i}$, $\vec{j}$ or a sign convention in 1D).
\item Compute the momentum $\vec{p}=m\vec{v}$. Its unit is $\mathrm{kg\,m\,s^{-1}}$ (equivalently N\,s).
\item For a change of velocity, compute $\Delta\vec{p}=m\vec{v}-m\vec{u}$ \emph{vectorially}: subtract components, never magnitudes blindly.
\item The impulse of a force is $\vec{J}=\vec{F}\,\Delta t$ (constant force) or $\vec{J}=\displaystyle\int_{t_1}^{t_2}\vec{F}\,dt$; the impulse--momentum theorem gives $\vec{J}=\Delta\vec{p}$.
\item If the force varies, the impulse is the \cle{area under the force--time graph}.
\end{enumerate}
\textbf{Reflex:} momentum and impulse are vectors — always state a positive direction and keep signs consistent.
\end{methode}

\begin{exoresolu}[Impulse on a football]
A football of mass $0.45$ kg is moving horizontally with velocity $(12\vec{i}+5\vec{j})\ \mathrm{m\,s^{-1}}$ when a player kicks it. Immediately after the kick its velocity is $(-8\vec{i}+14\vec{j})\ \mathrm{m\,s^{-1}}$.
\begin{enumerate}
\item[(a)] Find the momentum of the ball before and after the kick.
\item[(b)] Find the impulse imparted to the ball by the kick, and state its magnitude.
\item[(c)] The contact between boot and ball lasts $0.08$ s. Find the average force exerted on the ball.
\end{enumerate}
\tcblower
\textbf{Solution.} Take $\vec{i}$, $\vec{j}$ as fixed horizontal unit vectors.

\textbf{(a)} Momentum is $\vec{p}=m\vec{v}$ with $m=0.45$ kg.
\[ \vec{p}_{\text{before}}=0.45(12\vec{i}+5\vec{j})=(5.4\vec{i}+2.25\vec{j})\ \mathrm{kg\,m\,s^{-1}}. \]
\[ \vec{p}_{\text{after}}=0.45(-8\vec{i}+14\vec{j})=(-3.6\vec{i}+6.3\vec{j})\ \mathrm{kg\,m\,s^{-1}}. \]

\textbf{(b)} By the impulse--momentum theorem,
\begin{align*}
\vec{J}&=\Delta\vec{p}=\vec{p}_{\text{after}}-\vec{p}_{\text{before}}\\
&=(-3.6-5.4)\vec{i}+(6.3-2.25)\vec{j}\\
&=(-9\vec{i}+4.05\vec{j})\ \mathrm{N\,s}.
\end{align*}
Its magnitude is
\[ |\vec{J}|=\sqrt{9^{2}+4.05^{2}}=\sqrt{81+16.4025}=\sqrt{97.4025}\approx 9.87\ \mathrm{N\,s}. \]

\textbf{(c)} The average force satisfies $\vec{J}=\vec{F}_{\text{av}}\Delta t$, so
\[ \vec{F}_{\text{av}}=\frac{\vec{J}}{\Delta t}=\frac{-9\vec{i}+4.05\vec{j}}{0.08}=(-112.5\vec{i}+50.6\vec{j})\ \mathrm{N}. \]
Magnitude: $|\vec{F}_{\text{av}}|=\dfrac{9.87}{0.08}\approx 123\ \mathrm{N}$.

\textbf{Check:} the impulse points roughly opposite to the initial $\vec{i}$-motion (the ball is reversed in the $\vec{i}$-direction) and in the positive $\vec{j}$-direction — consistent with the observed change of velocity.
\end{exoresolu}

\courssub{Method 2: Applying Conservation of Linear Momentum}

\begin{propriete}[Principle of conservation of linear momentum]
If the resultant \emph{external} force on a system is zero (or its impulse during the event is negligible), then the total momentum of the system is constant:
\[ \sum m_i\vec{u}_i=\sum m_i\vec{v}_i. \]
Internal forces between the bodies (collision or explosion forces) occur in equal and opposite pairs by Newton's third law, so they cancel in the total. (The justification from Newton's second and third laws is examinable.)
\end{propriete}

\begin{methode}[Conservation of momentum in 1D]
\begin{enumerate}
\item Define the \cle{system} (the colliding or separating bodies) and check that external impulses during the event are negligible (short collision, smooth surface).
\item Choose a positive direction and write the velocity of each body \emph{with its sign}.
\item Write total momentum before = total momentum after:
\[ m_1u_1+m_2u_2=m_1v_1+m_2v_2. \]
\item Solve for the unknown. A negative answer means motion opposite to the chosen positive direction.
\item For bodies that \emph{coalesce} (stick together), use a single combined mass $(m_1+m_2)$ after impact.
\end{enumerate}
\textbf{Reflex:} draw a "before" and an "after" diagram with arrows — most sign errors disappear.
\end{methode}

\begin{exoresolu}[Trucks coupling on a straight rail]
A railway truck $A$ of mass $6000$ kg moving at $4\ \mathrm{m\,s^{-1}}$ on a straight horizontal track collides with a stationary truck $B$ of mass $4000$ kg. On impact the two trucks couple together and move as one.
\begin{enumerate}
\item[(a)] Find their common speed just after the collision.
\item[(b)] Find the kinetic energy lost in the collision and comment on the nature of the collision.
\end{enumerate}
\tcblower
\textbf{Solution.} Take the direction of motion of $A$ as positive. The track is horizontal and the collision is brief, so momentum is conserved.

\textbf{(a)} Before: $u_A=4$, $u_B=0$. After: combined mass $6000+4000=10000$ kg moving at speed $v$.
\begin{align*}
m_Au_A+m_Bu_B&=(m_A+m_B)v\\
6000(4)+4000(0)&=10000\,v\\
24000&=10000\,v \quad\Longrightarrow\quad v=2.4\ \mathrm{m\,s^{-1}}.
\end{align*}
The trucks move together at $2.4\ \mathrm{m\,s^{-1}}$ in the original direction of $A$.

\textbf{(b)} Kinetic energy before:
\[ KE_{\text{before}}=\tfrac12(6000)(4^2)+0=48000\ \mathrm{J}. \]
Kinetic energy after:
\[ KE_{\text{after}}=\tfrac12(10000)(2.4^2)=\tfrac12(10000)(5.76)=28800\ \mathrm{J}. \]
Loss: $48000-28800=19200\ \mathrm{J}$.

Since kinetic energy is \emph{not} conserved, the collision is \cle{inelastic} (perfectly inelastic, since the bodies stick). The lost energy becomes heat, sound and permanent deformation of the coupling mechanism.
\end{exoresolu}

\courssub{Method 3: Elastic Collisions in One Dimension}

\begin{methode}[Solving a 1D elastic collision]
\begin{enumerate}
\item Write \emph{two} equations: conservation of momentum
\[ m_1u_1+m_2u_2=m_1v_1+m_2v_2 \]
and conservation of kinetic energy (or, quicker, Newton's experimental law with $e=1$):
\[ v_2-v_1=u_1-u_2 \quad(\text{speed of separation}=\text{speed of approach}). \]
\item Solve the simultaneous linear system for $v_1$ and $v_2$.
\item Useful memorised results for a target at rest ($u_2=0$):
\[ v_1=\frac{m_1-m_2}{m_1+m_2}u_1,\qquad v_2=\frac{2m_1}{m_1+m_2}u_1. \]
\item Interpret signs: $v_1<0$ means the incoming body rebounds.
\end{enumerate}
\textbf{Reflex:} for an elastic collision never use energy conservation alone — you need both equations.
\end{methode}

\begin{exoresolu}[Elastic collision between two spheres]
A sphere $P$ of mass $2$ kg moving at $6\ \mathrm{m\,s^{-1}}$ collides directly with a sphere $Q$ of mass $3$ kg moving at $1\ \mathrm{m\,s^{-1}}$ in the \emph{same} direction on a smooth horizontal table. The collision is perfectly elastic. Find the velocities of $P$ and $Q$ after the collision.
\tcblower
\textbf{Solution.} Positive direction: direction of motion of $P$. Given $u_P=6$, $u_Q=1$.

\textbf{Momentum:}
\[ 2(6)+3(1)=2v_P+3v_Q \quad\Longrightarrow\quad 2v_P+3v_Q=15. \tag{1} \]

\textbf{Elastic condition} (separation speed = approach speed):
\[ v_Q-v_P=u_P-u_Q=6-1=5. \tag{2} \]

From (2): $v_Q=v_P+5$. Substitute in (1):
\begin{align*}
2v_P+3(v_P+5)&=15\\
5v_P+15&=15\\
v_P&=0,\qquad v_Q=5\ \mathrm{m\,s^{-1}}.
\end{align*}

\textbf{Answer:} $P$ is brought to rest; $Q$ moves on at $5\ \mathrm{m\,s^{-1}}$ in the positive direction.

\textbf{Verification.} Momentum after: $2(0)+3(5)=15$ \checkmark. KE before: $\tfrac12(2)(36)+\tfrac12(3)(1)=36+1.5=37.5$ J. KE after: $0+\tfrac12(3)(25)=37.5$ J \checkmark. Both conserved, as required for an elastic collision.
\end{exoresolu}

\courssub{Method 4: Newton's Law of Restitution (General 1D Impact)}

\begin{methode}[Impact with coefficient of restitution $e$]
\begin{enumerate}
\item Write conservation of momentum: $m_1u_1+m_2u_2=m_1v_1+m_2v_2$.
\item Write \cle{Newton's experimental law}:
\[ v_2-v_1=-e\,(u_2-u_1)\quad\Longleftrightarrow\quad \text{speed of separation}=e\times\text{speed of approach}, \]
with $0\le e\le 1$ ($e=1$ elastic, $e=0$ perfectly inelastic).
\item Solve the pair of linear equations.
\item For impact with a fixed wall, $u_2=v_2=0$, giving simply $v_1=-e\,u_1$ (rebound speed $=e\times$ approach speed).
\item Energy lost $=\tfrac12 m_1u_1^2+\tfrac12 m_2u_2^2-\tfrac12 m_1v_1^2-\tfrac12 m_2v_2^2$.
\end{enumerate}
\begin{attention}[Sign trap in the restitution law]
The law relates \emph{speeds of separation and approach}: always write it as $v_2-v_1=e(u_1-u_2)$ with body 2 initially \emph{ahead} of body 1. Reversing the order flips signs and ruins the solution.
\end{attention}
\end{methode}

\begin{exoresolu}[Ball rebounding from a wall]
A ball of mass $0.5$ kg moves at $10\ \mathrm{m\,s^{-1}}$ towards a fixed vertical wall. The coefficient of restitution between ball and wall is $e=0.6$.
\begin{enumerate}
\item[(a)] Find the velocity of the ball immediately after impact.
\item[(b)] Find the impulse exerted by the wall on the ball.
\item[(c)] Find the kinetic energy lost in the impact.
\end{enumerate}
\tcblower
\textbf{Solution.} Take the initial direction of the ball as positive: $u=10\ \mathrm{m\,s^{-1}}$.

\textbf{(a)} For a fixed wall, rebound speed $=e\times$ approach speed, direction reversed:
\[ v=-eu=-0.6(10)=-6\ \mathrm{m\,s^{-1}}. \]
The ball rebounds at $6\ \mathrm{m\,s^{-1}}$ away from the wall.

\textbf{(b)} Impulse = change of momentum:
\[ J=m(v-u)=0.5(-6-10)=0.5(-16)=-8\ \mathrm{N\,s}. \]
The wall exerts an impulse of magnitude $8$ N\,s on the ball, directed away from the wall.

\textbf{(c)}
\[ \Delta KE=\tfrac12 m u^2-\tfrac12 m v^2=\tfrac12(0.5)(100-36)=0.25\times 64=16\ \mathrm{J}. \]
So $16$ J is dissipated (sound, heat, deformation) in the impact.

\textbf{Remark:} note that the impulse involves the \emph{sum} $u+|v|$ because the velocity changes sign — a very common source of error is to write $J=m(u-v)$ with unsigned speeds.
\end{exoresolu}

\courssub{Method 5: Collisions in Two Dimensions}

\begin{methode}[2D collision problems]
\begin{enumerate}
\item Draw the situation: line of centres at impact, directions of motion before and after.
\item Apply conservation of momentum \cle{as a vector equation}: write one equation per axis ($\vec{i}$ and $\vec{j}$ components separately).
\item For a smooth sphere striking a smooth fixed surface: the velocity component \emph{parallel} to the surface is unchanged; the component \emph{perpendicular} to the surface is multiplied by $-e$.
\item For oblique impact of two smooth spheres: momentum is conserved along the line of centres (with restitution $e$ along that line), and the velocity of each sphere \emph{perpendicular} to the line of centres is unchanged.
\item Use $\tan\theta=\dfrac{\text{perpendicular component}}{\text{parallel component}}$ for deflection angles.
\end{enumerate}
\textbf{Reflex:} resolve \emph{before} writing any equation; never apply restitution along a direction other than the line of centres.
\end{methode}

\begin{exoresolu}[Oblique impact with a smooth wall]
A smooth sphere of mass $0.4$ kg moves on a smooth horizontal floor and strikes a smooth vertical wall. Just before impact its velocity is $(4\vec{i}+3\vec{j})\ \mathrm{m\,s^{-1}}$, where $\vec{i}$ is perpendicular to the wall (pointing towards it) and $\vec{j}$ is parallel to the wall. The coefficient of restitution is $e=\tfrac12$.
\begin{enumerate}
\item[(a)] Find the velocity of the sphere after the impact.
\item[(b)] Find the angle through which the direction of motion is deflected.
\item[(c)] Find the impulse exerted by the wall on the sphere.
\end{enumerate}
\tcblower
\textbf{Solution.}

\textbf{(a)} Component parallel to the wall ($\vec{j}$) is unchanged: $v_j=3$.
Component perpendicular to the wall is reversed and multiplied by $e$: $v_i=-e\,u_i=-\tfrac12(4)=-2$.
\[ \vec{v}=(-2\vec{i}+3\vec{j})\ \mathrm{m\,s^{-1}}. \]

\textbf{(b)} Before impact the velocity makes angle $\alpha$ with the wall where $\tan\alpha=\dfrac{u_i}{u_j}=\dfrac{4}{3}$, so $\alpha\approx 53.1^{\circ}$. After impact the angle with the wall is $\beta$ with $\tan\beta=\dfrac{|v_i|}{v_j}=\dfrac{2}{3}$, so $\beta\approx 33.7^{\circ}$. The total deflection is
\[ \alpha+\beta\approx 53.1^{\circ}+33.7^{\circ}=86.8^{\circ}. \]
(The incoming and outgoing paths lie on opposite sides of the normal, so the deflection is the sum.)

\textbf{(c)} Only the $\vec{i}$-component of momentum changes:
\[ \vec{J}=m(\vec{v}-\vec{u})=0.4\big[(-2-4)\vec{i}+(3-3)\vec{j}\big]=-2.4\,\vec{i}\ \mathrm{N\,s}. \]
The impulse has magnitude $2.4$ N\,s and is directed along the normal, away from the wall — as expected, since a smooth wall can only push along its normal.
\end{exoresolu}

\courssub{Method 6: Explosions, Recoil and Rocket Motion}

\begin{methode}[Explosions and recoil]
\begin{enumerate}
\item An explosion is a collision in reverse: total momentum is conserved (often zero before, hence zero after).
\item If the system starts at rest, write $m_1\vec{v}_1+m_2\vec{v}_2=\vec{0}$, so the fragments move in \emph{opposite} directions with momenta of equal magnitude: $m_1v_1=m_2v_2$.
\item Kinetic energy \emph{increases}: the chemical (or elastic) energy released equals the gain in KE.
\item For a \cle{rocket}: thrust $=u_{\text{rel}}\left|\dfrac{dm}{dt}\right|$, where $u_{\text{rel}}$ is the exhaust speed relative to the rocket; apply $F=ma$ with the \emph{instantaneous} mass.
\end{enumerate}
\textbf{Reflex:} in recoil problems the lighter fragment always gets the larger speed.
\end{methode}

\begin{exoresolu}[Recoil of a gun and rocket thrust]
\textbf{Part A.} A rifle of mass $4$ kg fires a bullet of mass $20$ g horizontally. The bullet leaves the muzzle at $400\ \mathrm{m\,s^{-1}}$. Find the recoil speed of the rifle and the total kinetic energy generated by the explosion of the charge.

\textbf{Part B.} A rocket of instantaneous mass $1200$ kg ejects gas at $1800\ \mathrm{m\,s^{-1}}$ (relative to the rocket) at a rate of $8\ \mathrm{kg\,s^{-1}}$. Find the thrust and the initial acceleration of the rocket in deep space (no external force).
\tcblower
\textbf{Part A.} System: rifle + bullet, initially at rest, so total momentum before $=0$.

Conservation of momentum (taking the bullet's direction as positive):
\[ 0=m_b v_b+m_r v_r \quad\Longrightarrow\quad 0=(0.020)(400)+4\,v_r. \]
\[ v_r=-\frac{8}{4}=-2\ \mathrm{m\,s^{-1}}. \]
The rifle recoils at $2\ \mathrm{m\,s^{-1}}$ opposite to the bullet.

Kinetic energy generated:
\[ KE=\tfrac12(0.020)(400)^2+\tfrac12(4)(2)^2=1600+8=1608\ \mathrm{J}. \]
This energy comes from the chemical energy of the charge. Note the bullet carries $1600/1608\approx 99.5\%$ of the kinetic energy, although the momenta of rifle and bullet have equal magnitudes.

\textbf{Part B.} Thrust:
\[ F=u_{\text{rel}}\left|\frac{dm}{dt}\right|=1800\times 8=14400\ \mathrm{N}. \]
Newton's second law with instantaneous mass $m=1200$ kg:
\[ a=\frac{F}{m}=\frac{14400}{1200}=12\ \mathrm{m\,s^{-2}}. \]
As fuel burns, $m$ decreases while the thrust stays constant, so the acceleration \emph{increases} with time — this is why rocket acceleration is not uniform.
\end{exoresolu}

\courssub{Method 7: Force--Time Graphs and Safety Applications}

\begin{methode}[Reading and using $F$--$t$ graphs]
\begin{enumerate}
\item The impulse equals the \cle{area} between the $F$--$t$ curve and the time axis: $J=\displaystyle\int F\,dt$ (triangles, rectangles or trapezia for straight-line graphs; integration for curved graphs).
\item Set this area equal to the change of momentum $m\Delta v$ to find an unknown speed, mass or time.
\item The \emph{average} force is $\bar{F}=\dfrac{J}{\Delta t}=\dfrac{\text{area}}{\text{base}}$.
\item \textbf{Safety principle:} for a fixed change of momentum, increasing the impact time $\Delta t$ decreases the peak and average force. This explains crumple zones, airbags, seatbelts, crash helmets, bending the knees on landing, and follow-through in sport.
\end{enumerate}
\end{methode}

\begin{exoresolu}[Airbag force--time graph]
In a road accident test, a dummy's head of mass $5$ kg travelling at $15\ \mathrm{m\,s^{-1}}$ is brought to rest by an airbag. The force $F$ (in N) exerted on the head varies with time $t$ (in s) according to
\[ F(t)=450000\,t\,(0.1-t),\qquad 0\le t\le 0.1. \]
\begin{enumerate}
\item[(a)] Sketch the force--time graph and find the maximum force.
\item[(b)] Show that the impulse delivered by the airbag equals the change of momentum of the head.
\item[(c)] Without the airbag the head would be stopped by the dashboard in $0.01$ s with approximately constant force. Compare the forces in the two cases and comment.
\end{enumerate}
\tcblower
\textbf{Solution.}

\textbf{(a)} Expanding, $F(t)=45000t-450000t^{2}$: a downward parabola with $F(0)=F(0.1)=0$. The maximum occurs where
\[ F'(t)=45000-900000t=0 \quad\Longrightarrow\quad t=0.05\ \mathrm{s}. \]
\[ F_{\max}=450000(0.05)(0.1-0.05)=450000\times 0.0025=1125\ \mathrm{N}. \]

\begin{popfigure}
\begin{center}
\begin{tikzpicture}
\begin{axis}[width=8.5cm,height=5.5cm,axis lines=left,xlabel={$t$ (s)},ylabel={$F$ (N)},xmin=0,xmax=0.12,ymin=0,ymax=1300,xtick={0,0.05,0.1},ytick={0,1125},domain=0:0.1,samples=60]
\addplot[popblue,very thick] {450000*x*(0.1-x)};
\addplot[poporange,dashed] coordinates {(0.05,0) (0.05,1125)};
\node[popdark] at (axis cs:0.078,950) {$F_{\max}=1125$};
\end{axis}
\end{tikzpicture}
\end{center}
\legende{Force--time curve during airbag contact (parabolic model). The shaded region under the curve is the impulse.}
\end{popfigure}

\textbf{(b)} Impulse = area under the curve:
\begin{align*}
J&=\int_0^{0.1}\left(45000t-450000t^{2}\right)dt
=\Big[22500t^{2}-150000t^{3}\Big]_0^{0.1}\\
&=22500(0.01)-150000(0.001)=225-150=75\ \mathrm{N\,s}.
\end{align*}
Change of momentum of the head (brought from $15\ \mathrm{m\,s^{-1}}$ to rest):
\[ \Delta p=m\,\Delta v=5\times 15=75\ \mathrm{kg\,m\,s^{-1}}=75\ \mathrm{N\,s}. \]
Hence $J=\Delta p$: the impulse--momentum theorem is verified. (Quick check: the area of a parabola is $\tfrac23\times\text{base}\times\text{height}=\tfrac23(0.1)(1125)=75$ N\,s \checkmark.)

\textbf{(c)} Without the airbag, a constant force acting for $0.01$ s must deliver the same impulse $75$ N\,s:
\[ \bar{F}=\frac{J}{\Delta t}=\frac{75}{0.01}=7500\ \mathrm{N}. \]
With the airbag the peak force is only $1125$ N. Multiplying the stopping time by $10$ has divided the force by a factor of about $7$: this is exactly the principle of the airbag — \emph{same change of momentum, longer time, smaller force} — and it is why airbags, seatbelts and crumple zones save lives on Cameroonian roads as everywhere else.
\end{exoresolu}

\courssub{Method 8: Centre of Mass and Momentum of a System}

\begin{methode}[Centre of mass and total momentum]
\begin{enumerate}
\item For particles $m_i$ at positions $\vec{r}_i$: $\vec{r}_{CM}=\dfrac{\sum m_i\vec{r}_i}{\sum m_i}$; in 1D, $\bar{x}=\dfrac{\sum m_ix_i}{\sum m_i}$.
\item The total momentum of the system equals the momentum of a single particle carrying the total mass at the centre of mass:
\[ \vec{P}=M\vec{v}_{CM}. \]
\item If no external force acts, $\vec{v}_{CM}$ is \emph{constant} — even during collisions or explosions between the parts.
\item Use this to find the motion of one fragment when the other(s) and the initial velocity are known.
\end{enumerate}
\end{methode}

\begin{exoresolu}[Exploding shell and the centre of mass]
A shell of mass $6$ kg is fired and, at the top of its flight, is moving horizontally at $50\ \mathrm{m\,s^{-1}}$ when it explodes into two fragments of masses $4$ kg and $2$ kg. The $4$ kg fragment is brought to rest (momentarily) by the explosion. Air resistance is neglected.
\begin{enumerate}
\item[(a)] Find the velocity of the $2$ kg fragment just after the explosion.
\item[(b)] Verify that the velocity of the centre of mass is unchanged by the explosion.
\item[(c)] Find the kinetic energy released by the explosion.
\end{enumerate}
\tcblower
\textbf{Solution.} Take the direction of the shell's motion as positive. Before explosion: $M=6$ kg, $u=50\ \mathrm{m\,s^{-1}}$.

\textbf{(a)} Momentum is conserved (the explosion forces are internal; gravity acts vertically and has negligible impulse during the instant of explosion):
\begin{align*}
Mu&=m_1v_1+m_2v_2\\
6(50)&=4(0)+2v_2\\
300&=2v_2 \quad\Longrightarrow\quad v_2=150\ \mathrm{m\,s^{-1}}.
\end{align*}
The $2$ kg fragment shoots forward at $150\ \mathrm{m\,s^{-1}}$.

\textbf{(b)} Velocity of centre of mass after the explosion:
\[ v_{CM}=\frac{m_1v_1+m_2v_2}{m_1+m_2}=\frac{4(0)+2(150)}{6}=\frac{300}{6}=50\ \mathrm{m\,s^{-1}}. \]
This equals the shell's velocity before the explosion: the centre of mass continues exactly as if nothing had happened, and it would follow the original parabolic trajectory until a fragment hits the ground.

\textbf{(c)} Kinetic energy before:
\[ KE_{\text{before}}=\tfrac12(6)(50)^2=7500\ \mathrm{J}. \]
Kinetic energy after:
\[ KE_{\text{after}}=\tfrac12(4)(0)^2+\tfrac12(2)(150)^2=22500\ \mathrm{J}. \]
Energy released by the explosion:
\[ \Delta KE=22500-7500=15000\ \mathrm{J}, \]
supplied by the chemical energy of the explosive charge.

\textbf{Remark:} note the contrast with collisions — in an explosion kinetic energy \emph{increases}, while momentum is still conserved. Momentum conservation is the universal tool; energy bookkeeping tells you the \emph{type} of event.
\end{exoresolu}

\begin{aretenir}[Key points of the chapter]
\begin{itemize}
\item Momentum $\vec{p}=m\vec{v}$ and impulse $\vec{J}=\int\vec{F}\,dt$ are \cle{vectors}; $\vec{J}=\Delta\vec{p}$ always.
\item Total momentum is conserved whenever the external impulse is negligible — in collisions \emph{and} in explosions.
\item Kinetic energy is conserved only in \cle{elastic} collisions ($e=1$); it decreases in inelastic collisions ($e<1$) and increases in explosions.
\item Newton's experimental law: speed of separation $=e\times$ speed of approach, along the line of centres (or the normal to a wall).
\item In 2D, resolve: momentum is conserved component by component; perpendicular components to a smooth wall are unchanged.
\item Rocket thrust $=u_{\text{rel}}\left|\dfrac{dm}{dt}\right|$; the centre of mass of an isolated system moves with constant velocity.
\item Safety devices work by lengthening the impact time: same $\Delta\vec{p}$, smaller force.
\end{itemize}
\end{aretenir}

\newpage

\coursec{Exercises}

\courssub{I check my knowledge}

\begin{exercice}[Multiple choice: the nature of momentum]
Choose the correct answer and justify briefly.
\begin{enumerate}
\item The momentum of a body of mass $m$ moving with velocity $\vec{v}$ is:
\begin{enumerate}
\item a scalar equal to $mv$;
\item a vector $\vec{p}=m\vec{v}$, parallel to $\vec{v}$;
\item a vector $m\vec{v}$, perpendicular to $\vec{v}$;
\item a scalar equal to $\tfrac{1}{2}mv^{2}$.
\end{enumerate}
\item A $2\ \mathrm{kg}$ stone and a $1\ \mathrm{kg}$ stone have the same kinetic energy. Which has the greater momentum?
\begin{enumerate}
\item the $2\ \mathrm{kg}$ stone;
\item the $1\ \mathrm{kg}$ stone;
\item they are equal;
\item impossible to tell.
\end{enumerate}
\end{enumerate}
\end{exercice}

\begin{exercice}[True or false]
State whether each statement is true or false, and justify your answer.
\begin{enumerate}
\item If the total external force on a system is zero, the total momentum of the system is constant.
\item In a perfectly inelastic collision, kinetic energy is conserved.
\item The impulse of a force equals the area under the force--time graph.
\item During an explosion of a stationary body, the fragments always move in the same direction.
\end{enumerate}
\end{exercice}

\begin{exercice}[Definitions]
Define each of the following, giving the SI unit of each quantity:
\begin{enumerate}
\item the linear momentum of a particle;
\item the impulse of a force;
\item an elastic collision;
\item the centre of mass of a system of particles.
\end{enumerate}
\end{exercice}

\begin{exercice}[Direct calculations]
\begin{enumerate}
\item A lorry of mass $8000\ \mathrm{kg}$ travels along a straight road at $15\ \mathrm{m\,s^{-1}}$. Calculate its momentum.
\item A tennis ball of mass $0.058\ \mathrm{kg}$ moving at $30\ \mathrm{m\,s^{-1}}$ is stopped by a racket in $0.05\ \mathrm{s}$. Find the average force exerted on the ball.
\item A constant force of $40\ \mathrm{N}$ acts on a body for $0.25\ \mathrm{s}$. Find the impulse delivered to the body.
\end{enumerate}
\end{exercice}

\courssub{I practise}

\begin{exercice}[Kicking a football]
A $0.45\ \mathrm{kg}$ football moving at $18\ \mathrm{m\,s^{-1}}$ is kicked back along the same line at $24\ \mathrm{m\,s^{-1}}$; the foot is in contact with the ball for $0.08\ \mathrm{s}$.
\begin{enumerate}
\item Find the impulse given to the ball.
\item Find the average force exerted on the ball.
\end{enumerate}
\end{exercice}

\begin{exercice}[Man and canoe]
A $60\ \mathrm{kg}$ man jumps horizontally at $3\ \mathrm{m\,s^{-1}}$ from a stationary $90\ \mathrm{kg}$ canoe.
\begin{enumerate}
\item Find the recoil speed of the canoe.
\item State the physical principle you have used, and explain why it applies here.
\end{enumerate}
\end{exercice}

\begin{exercice}[Trolleys that stick together]
A $2\ \mathrm{kg}$ trolley moving at $5\ \mathrm{m\,s^{-1}}$ collides head-on with a $3\ \mathrm{kg}$ trolley moving at $1\ \mathrm{m\,s^{-1}}$ in the opposite direction; the trolleys stick together after impact.
\begin{enumerate}
\item Find their common velocity just after the collision.
\item Find the percentage of kinetic energy lost in the collision.
\end{enumerate}
\end{exercice}

\begin{exercice}[Force--time graph]
A force acting on a body varies with time as shown on a force--time graph consisting of a triangle of base $0.4\ \mathrm{s}$ and peak value $500\ \mathrm{N}$.
\begin{enumerate}
\item Find the impulse of the force.
\item Find the change in velocity of a $4\ \mathrm{kg}$ body, initially at rest, on which this force acts.
\end{enumerate}
\end{exercice}

\courssub{I prepare for the GCE}

\begin{exercice}[Collision at a crossroads \hfill (12 marks)]
A $1500\ \mathrm{kg}$ car travelling east at $25\ \mathrm{m\,s^{-1}}$ collides at a crossroads with a $2000\ \mathrm{kg}$ truck travelling north at $15\ \mathrm{m\,s^{-1}}$. The vehicles lock together on impact.
\begin{enumerate}
\item State the principle you will use and the conditions under which it is valid. \hfill (2)
\item Write down the components of the total momentum of the system before the collision along the east and north directions. \hfill (3)
\item Find the speed of the wreck immediately after the collision. \hfill (3)
\item Find the direction of motion of the wreck, giving the angle it makes with the east direction. \hfill (2)
\item Calculate the kinetic energy lost in the collision and comment on the nature of this collision. \hfill (2)
\end{enumerate}
\end{exercice}

\begin{exercice}[Ballistic pendulum and airbag safety \hfill (14 marks)]
\textbf{Part A.} A bullet of mass $20\ \mathrm{g}$ is fired horizontally at $400\ \mathrm{m\,s^{-1}}$ into a $2\ \mathrm{kg}$ wooden block hanging at rest, and becomes embedded in it.
\begin{enumerate}
\item Find the velocity of the block and bullet just after impact. \hfill (3)
\item Find the fraction of the initial kinetic energy lost in the impact, and state what becomes of this energy. \hfill (3)
\end{enumerate}
\textbf{Part B.} A passenger of mass $70\ \mathrm{kg}$ is brought to rest from $20\ \mathrm{m\,s^{-1}}$ during a crash.
\begin{enumerate}
\item[(c)] Compute the average force on the passenger if the stopping time is $0.01\ \mathrm{s}$ (no airbag) and if it is $0.12\ \mathrm{s}$ (with airbag). \hfill (4)
\item[(d)] Using the impulse--momentum theorem, explain clearly why an airbag reduces injuries, and name one other safety device based on the same principle. \hfill (4)
\end{enumerate}
\end{exercice}

\begin{exercice}[Explosion, rocket and centre of mass \hfill (14 marks)]
\textbf{Part A.} A stationary shell of mass $12\ \mathrm{kg}$ explodes into two fragments of masses $5\ \mathrm{kg}$ and $7\ \mathrm{kg}$. The $5\ \mathrm{kg}$ fragment moves off at $84\ \mathrm{m\,s^{-1}}$.
\begin{enumerate}
\item Find the velocity of the $7\ \mathrm{kg}$ fragment. \hfill (3)
\item Find the total kinetic energy released by the explosion, and explain its origin. \hfill (3)
\end{enumerate}
\textbf{Part B.} A rocket ejects exhaust gas at $2000\ \mathrm{m\,s^{-1}}$ relative to the rocket, at a rate of $150\ \mathrm{kg\,s^{-1}}$.
\begin{enumerate}
\item[(c)] Calculate the thrust on the rocket. \hfill (2)
\item[(d)] Using momentum considerations, explain why the rocket accelerates even in vacuum, where there is nothing to ``push against''. \hfill (3)
\end{enumerate}
\textbf{Part C.} Three particles of masses $1\ \mathrm{kg}$, $2\ \mathrm{kg}$ and $3\ \mathrm{kg}$ are placed at the points $(0,0)$, $(4,0)$ and $(0,3)$ metres respectively.
\begin{enumerate}
\item[(e)] Find the position of the centre of mass of the system. \hfill (2)
\item[(f)] Show that if the system is isolated, its centre of mass moves with constant velocity. \hfill (1)
\end{enumerate}
\end{exercice}

\newpage

\coursec{Solutions to the exercises}

\begin{corrige}[Exercise 1 — Multiple choice: the nature of momentum]
\begin{enumerate}
\item \textbf{Correct answer: (b) a vector $\vec{p}=m\vec{v}$, parallel to $\vec{v}$.} \\
\textbf{Justification:} Linear momentum is defined as the product of mass (a scalar) and velocity (a vector). Multiplying a vector by a positive scalar gives a vector with the \emph{same direction} as the original vector, so $\vec{p}$ is parallel to $\vec{v}$. Options (a) and (d) describe scalars (kinetic energy is $K=\frac{1}{2}mv^2$), and (c) is incorrect because the momentum vector is parallel, not perpendicular, to the velocity.

\item \textbf{Correct answer: (a) the $2\ \mathrm{kg}$ stone.} \\
\textbf{Justification:} Kinetic energy $K = \frac{1}{2}mv^2$ is the same for both stones. Let $m_1=2\ \mathrm{kg}$, $m_2=1\ \mathrm{kg}$. Then
\[ \frac{1}{2}m_1v_1^2 = \frac{1}{2}m_2v_2^2 \;\Rightarrow\; v_1 = v_2\sqrt{\frac{m_2}{m_1}} = \frac{v_2}{\sqrt{2}}. \]
Momentum is $p = mv$, so
\[ p_1 = m_1v_1 = 2\times\frac{v_2}{\sqrt{2}} = \sqrt{2}\,v_2, \qquad p_2 = m_2v_2 = v_2. \]
Since $\sqrt{2}>1$, the heavier stone has the greater momentum. (Equivalently, $p=\sqrt{2mK}$, so for equal $K$ the larger mass has the larger momentum.)
\end{enumerate}
\end{corrige}

\begin{corrige}[Exercise 2 — True or false]
\begin{enumerate}
\item \textbf{True.} If the net external force on a system is zero, the total momentum of the system is conserved (it remains constant in both magnitude and direction). This is the principle of conservation of linear momentum, which follows from $\vec{F}_{\mathrm{ext}} = \dfrac{d\vec{P}}{dt} = \vec{0}$.

\item \textbf{False.} In a perfectly inelastic collision, the colliding bodies stick together and move with a common velocity after impact. Kinetic energy is \emph{not} conserved; part of it is transformed into internal energy (heat, sound, permanent deformation). Only the total momentum is conserved.

\item \textbf{True.} Impulse $\vec{I} = \displaystyle\int_{t_1}^{t_2} \vec{F}\,dt$ is exactly the area under the force–time graph. For a constant force, $I = F\Delta t$ (area of a rectangle); for a varying force, the integral gives the area under the curve.

\item \textbf{False.} Before the explosion the body is stationary, so the total momentum is zero. By conservation of momentum, the vector sum of the momenta of all the fragments after the explosion must also be zero. Therefore the fragments cannot all move in the same direction: their momenta must cancel, so for two fragments they move in \emph{opposite} directions.
\end{enumerate}
\end{corrige}

\begin{corrige}[Exercise 3 — Definitions]
\begin{enumerate}
\item \textbf{Linear momentum of a particle:} The linear momentum $\vec{p}$ of a particle of mass $m$ moving with velocity $\vec{v}$ is defined by $\vec{p} = m\vec{v}$. It is a vector quantity, parallel to $\vec{v}$. SI unit: $\mathrm{kg\,m\,s^{-1}}$ (equivalently $\mathrm{N\,s}$).

\item \textbf{Impulse of a force:} The impulse $\vec{I}$ of a force $\vec{F}$ acting over the time interval $[t_1,t_2]$ is defined by $\vec{I} = \displaystyle\int_{t_1}^{t_2} \vec{F}\,dt$. For a constant force, $\vec{I} = \vec{F}\,\Delta t$. By the impulse–momentum theorem, the impulse equals the change in momentum of the body: $\vec{I} = \Delta\vec{p}$. SI unit: $\mathrm{N\,s}$ (equivalently $\mathrm{kg\,m\,s^{-1}}$).

\item \textbf{Elastic collision:} A collision in which both the total momentum and the total kinetic energy of the system are conserved. The bodies do not stick together and no kinetic energy is converted into internal energy.

\item \textbf{Centre of mass of a system of particles:} The centre of mass is the point whose position vector is
\[ \vec{R}_{\mathrm{cm}} = \frac{\sum_i m_i \vec{r}_i}{\sum_i m_i}, \]
where $m_i$ and $\vec{r}_i$ are the masses and position vectors of the particles. The whole system moves as if its total mass were concentrated at this point and all the external forces acted there; in particular $\vec{P} = M\vec{v}_{\mathrm{cm}}$. SI unit of $\vec{R}_{\mathrm{cm}}$: metre ($\mathrm{m}$).
\end{enumerate}
\end{corrige}

\begin{corrige}[Exercise 4 — Direct calculations]
\begin{enumerate}
\item Momentum $p = mv = 8000 \times 15 = 120\,000\ \mathrm{kg\,m\,s^{-1}}$ (i.e. $1.2\times10^5\ \mathrm{N\,s}$). Direction: along the road, in the direction of motion.

\item Change in momentum of the ball (taking the initial direction of motion as positive):
\[ \Delta p = m(v_f - v_i) = 0.058\,(0 - 30) = -1.74\ \mathrm{kg\,m\,s^{-1}}. \]
Magnitude of the impulse $= 1.74\ \mathrm{N\,s}$. Average force:
\[ F_{\mathrm{avg}} = \frac{|\Delta p|}{\Delta t} = \frac{1.74}{0.05} = 34.8\ \mathrm{N}, \]
directed opposite to the initial motion of the ball.

\item Impulse $I = F\,\Delta t = 40 \times 0.25 = 10\ \mathrm{N\,s}$, in the direction of the force.
\end{enumerate}
\end{corrige}

\begin{corrige}[Exercise 5 — Kicking a football]
Take the initial direction of motion of the ball as positive.
\begin{enumerate}
\item Initial momentum: $p_i = 0.45 \times 18 = 8.1\ \mathrm{kg\,m\,s^{-1}}$. \\
Final momentum: $p_f = 0.45 \times (-24) = -10.8\ \mathrm{kg\,m\,s^{-1}}$ (the ball moves back, hence the negative sign). \\
Impulse: $I = \Delta p = p_f - p_i = -10.8 - 8.1 = -18.9\ \mathrm{N\,s}$. \\
Magnitude $= 18.9\ \mathrm{N\,s}$; direction opposite to the initial motion of the ball.

\item Average force:
\[ F_{\mathrm{avg}} = \frac{I}{\Delta t} = \frac{-18.9}{0.08} = -236.25\ \mathrm{N}. \]
Magnitude $\approx 236\ \mathrm{N}$; the negative sign shows that the force exerted by the foot on the ball acts opposite to the initial motion, i.e. in the direction of the final velocity.
\end{enumerate}
\end{corrige}

\begin{corrige}[Exercise 6 — Man and canoe]
\begin{enumerate}
\item Take the direction of the man's jump as positive. Initially the man and the canoe are at rest, so the total momentum is $0$. After the jump, conservation of momentum gives
\[ m_{\mathrm{man}}v_{\mathrm{man}} + m_{\mathrm{canoe}}v_{\mathrm{canoe}} = 0
\;\Rightarrow\; 60 \times 3 + 90\,v_{\mathrm{canoe}} = 0
\;\Rightarrow\; v_{\mathrm{canoe}} = -\frac{180}{90} = -2\ \mathrm{m\,s^{-1}}. \]
The canoe recoils at $2\ \mathrm{m\,s^{-1}}$ in the direction opposite to the man's jump.

\item \textbf{Principle used:} conservation of linear momentum. It applies because the man–canoe system is isolated in the horizontal direction: during the very short time of the jump, the external horizontal forces (water resistance) are negligible. The vertical forces (weights and buoyancy) balance each other and do not affect the horizontal momentum.
\end{enumerate}
\end{corrige}

\begin{corrige}[Exercise 7 — Trolleys that stick together]
Take the direction of motion of the $2\ \mathrm{kg}$ trolley as positive.
\begin{enumerate}
\item Initial momentum:
\[ p_i = (2 \times 5) + \bigl(3 \times (-1)\bigr) = 10 - 3 = 7\ \mathrm{kg\,m\,s^{-1}}. \]
After the collision the trolleys stick together: combined mass $= 2+3 = 5\ \mathrm{kg}$, common velocity $v$. Conservation of momentum:
\[ 5v = 7 \;\Rightarrow\; v = 1.4\ \mathrm{m\,s^{-1}}, \]
in the positive direction (the direction initially taken by the $2\ \mathrm{kg}$ trolley).

\item Initial kinetic energy:
\[ K_i = \frac{1}{2}(2)(5^2) + \frac{1}{2}(3)(1^2) = 25 + 1.5 = 26.5\ \mathrm{J}. \]
Final kinetic energy:
\[ K_f = \frac{1}{2}(5)(1.4^2) = 2.5 \times 1.96 = 4.9\ \mathrm{J}. \]
Energy lost: $\Delta K = 26.5 - 4.9 = 21.6\ \mathrm{J}$. \\
Percentage lost:
\[ \frac{21.6}{26.5} \times 100\% \approx 81.5\%. \]
This large loss is typical of a perfectly inelastic collision: the lost kinetic energy becomes heat, sound and deformation of the coupling.
\end{enumerate}
\end{corrige}

\begin{corrige}[Exercise 8 — Force–time graph]
The graph is a triangle of base $0.4\ \mathrm{s}$ and height $500\ \mathrm{N}$.
\begin{enumerate}
\item The impulse equals the area under the force–time graph:
\[ I = \text{area of triangle} = \frac{1}{2} \times \text{base} \times \text{height} = \frac{1}{2} \times 0.4 \times 500 = 100\ \mathrm{N\,s}. \]

\item By the impulse–momentum theorem, $I = \Delta p = m\,\Delta v$. For a $4\ \mathrm{kg}$ body initially at rest:
\[ \Delta v = \frac{I}{m} = \frac{100}{4} = 25\ \mathrm{m\,s^{-1}}. \]
The body reaches a velocity of $25\ \mathrm{m\,s^{-1}}$ in the direction of the force.
\end{enumerate}
\end{corrige}

\begin{corrige}[Exercise 9 — Collision at a crossroads]
\textbf{Mark scheme and examiner expectations:}
\begin{itemize}
\item \textbf{(a) 2 marks:} State the principle (conservation of linear momentum) and the condition (external forces negligible during the very short collision, so their impulse is approximately zero).
\item \textbf{(b) 3 marks:} Correct components: east $p_x = 1500 \times 25 = 37\,500\ \mathrm{kg\,m\,s^{-1}}$; north $p_y = 2000 \times 15 = 30\,000\ \mathrm{kg\,m\,s^{-1}}$.
\item \textbf{(c) 3 marks:} Combined mass $= 3500\ \mathrm{kg}$; speed $v = \dfrac{\sqrt{37\,500^2+30\,000^2}}{3500} \approx 13.7\ \mathrm{m\,s^{-1}}$.
\item \textbf{(d) 2 marks:} Angle $\theta = \arctan(30\,000/37\,500) \approx 38.7^\circ$ north of east.
\item \textbf{(e) 2 marks:} KE lost $\approx 3.64\times10^5\ \mathrm{J}$ (about $53\%$ of the initial KE). Comment: perfectly inelastic collision (vehicles lock together), kinetic energy not conserved.
\end{itemize}

\textbf{Detailed solution:}
\begin{enumerate}
\item \textbf{Principle:} conservation of linear momentum, applied component by component. It is valid because during the very short collision time the external forces (friction, air resistance) are negligible compared with the large internal forces between the vehicles, so the impulse of the external forces is approximately zero.

\item \textbf{Momentum components before the collision:} \\
East ($x$-direction): $p_x = m_{\mathrm{car}} v_{\mathrm{car}} = 1500 \times 25 = 37\,500\ \mathrm{kg\,m\,s^{-1}}$. \\
North ($y$-direction): $p_y = m_{\mathrm{truck}} v_{\mathrm{truck}} = 2000 \times 15 = 30\,000\ \mathrm{kg\,m\,s^{-1}}$.

\item \textbf{Speed after the collision:} combined mass $M = 1500+2000 = 3500\ \mathrm{kg}$. Momentum is conserved in each direction:
\[ M v_x = p_x \;\Rightarrow\; v_x = \frac{37\,500}{3500} = 10.71\ \mathrm{m\,s^{-1}}, \qquad
M v_y = p_y \;\Rightarrow\; v_y = \frac{30\,000}{3500} = 8.57\ \mathrm{m\,s^{-1}}. \]
Hence
\[ v = \sqrt{v_x^2+v_y^2} = \sqrt{10.71^2+8.57^2} = \sqrt{188.3} \approx 13.7\ \mathrm{m\,s^{-1}}. \]

\item \textbf{Direction:}
\[ \theta = \arctan\!\left(\frac{v_y}{v_x}\right) = \arctan\!\left(\frac{8.57}{10.71}\right) = \arctan(0.8) \approx 38.7^\circ \quad \text{north of east}. \]

\item \textbf{Kinetic energy lost:}
\[ K_i = \frac{1}{2}(1500)(25^2) + \frac{1}{2}(2000)(15^2) = 468\,750 + 225\,000 = 693\,750\ \mathrm{J}, \]
\[ K_f = \frac{1}{2}(3500)(13.7^2) \approx 328\,500\ \mathrm{J}, \]
\[ \Delta K = K_i - K_f \approx 693\,750 - 328\,500 \approx 3.65\times10^5\ \mathrm{J} \quad (\text{about } 53\% \text{ of } K_i). \]
\textbf{Comment:} the collision is perfectly inelastic (the vehicles lock together), so kinetic energy is not conserved; the lost energy is transformed into heat, sound and the work of deforming the vehicles.
\end{enumerate}
\end{corrige}

\begin{corrige}[Exercise 10 — Ballistic pendulum and airbag safety]
\textbf{Mark scheme and examiner expectations:}
\begin{itemize}
\item \textbf{Part A (a) 3 marks:} Conservation of momentum: $0.02 \times 400 = 2.02\,v \Rightarrow v \approx 3.96\ \mathrm{m\,s^{-1}}$.
\item \textbf{Part A (b) 3 marks:} $K_i = 1600\ \mathrm{J}$, $K_f \approx 15.8\ \mathrm{J}$; fraction lost $\approx 0.99$. Energy becomes heat, sound and deformation.
\item \textbf{Part B (c) 4 marks:} Without airbag: $F = \dfrac{70 \times 20}{0.01} = 140\,000\ \mathrm{N}$. With airbag: $F = \dfrac{1400}{0.12} \approx 11\,700\ \mathrm{N}$.
\item \textbf{Part B (d) 4 marks:} Impulse–momentum theorem: $\Delta p$ fixed; increasing $\Delta t$ reduces the average force. Other device: seat belt, crumple zone or crash helmet.
\end{itemize}

\textbf{Detailed solution:}

\textbf{Part A.}
\begin{enumerate}
\item Let $m_b = 0.020\ \mathrm{kg}$, $u_b = 400\ \mathrm{m\,s^{-1}}$, $M = 2.00\ \mathrm{kg}$ (block initially at rest). After the bullet embeds itself, the combined mass is $M+m_b = 2.02\ \mathrm{kg}$ moving with velocity $v$. Conservation of momentum:
\[ m_b u_b = (M+m_b)v \;\Rightarrow\; v = \frac{0.02 \times 400}{2.02} = \frac{8}{2.02} \approx 3.96\ \mathrm{m\,s^{-1}}. \]

\item Initial kinetic energy:
\[ K_i = \frac{1}{2}m_b u_b^2 = \frac{1}{2} \times 0.02 \times 400^2 = 1600\ \mathrm{J}. \]
Final kinetic energy:
\[ K_f = \frac{1}{2}(2.02)(3.96^2) \approx 15.8\ \mathrm{J}. \]
Fraction lost:
\[ \frac{K_i - K_f}{K_i} = \frac{1600-15.8}{1600} \approx 0.990 \quad (\text{i.e. about } 99\%). \]
The lost kinetic energy is converted into heat, sound and the work done in deforming the block and the bullet (internal energy).
\end{enumerate}

\textbf{Part B.}
\begin{enumerate}
\item[(c)] Passenger mass $m = 70\ \mathrm{kg}$, initial speed $u = 20\ \mathrm{m\,s^{-1}}$, final speed $0$. The magnitude of the change of momentum is
\[ |\Delta p| = m\,|\Delta v| = 70 \times 20 = 1400\ \mathrm{kg\,m\,s^{-1}}. \]
Without airbag ($\Delta t = 0.01\ \mathrm{s}$):
\[ F_{\mathrm{avg}} = \frac{|\Delta p|}{\Delta t} = \frac{1400}{0.01} = 140\,000\ \mathrm{N}. \]
With airbag ($\Delta t = 0.12\ \mathrm{s}$):
\[ F_{\mathrm{avg}} = \frac{1400}{0.12} \approx 11\,700\ \mathrm{N}. \]

\item[(d)] The impulse–momentum theorem states $F_{\mathrm{avg}}\,\Delta t = \Delta p$. For a given crash, the change of momentum $\Delta p$ of the passenger is fixed by the initial speed. By increasing the stopping time $\Delta t$ (from $0.01\ \mathrm{s}$ to $0.12\ \mathrm{s}$, a factor of $12$), the airbag divides the average force on the passenger by $12$. This much smaller force reduces the pressure on the body and therefore the risk of serious injury. Another safety device based on the same principle is the \textbf{crumple zone} of a car (also seat belts and crash helmets), which lengthens the time over which the momentum of the occupants changes.
\end{enumerate}
\end{corrige}

\begin{corrige}[Exercise 11 — Explosion, rocket and centre of mass]
\textbf{Mark scheme and examiner expectations:}
\begin{itemize}
\item \textbf{Part A (a) 3 marks:} Conservation of momentum: $0 = 5 \times 84 + 7v \Rightarrow v = -60\ \mathrm{m\,s^{-1}}$.
\item \textbf{Part A (b) 3 marks:} KE released $= \frac{1}{2}(5)(84^2) + \frac{1}{2}(7)(60^2) = 30\,240\ \mathrm{J}$. Origin: chemical potential energy of the explosive.
\item \textbf{Part B (c) 2 marks:} Thrust $= v_{\mathrm{rel}} \dfrac{dm}{dt} = 2000 \times 150 = 300\,000\ \mathrm{N}$.
\item \textbf{Part B (d) 3 marks:} The rocket pushes the exhaust backward; by Newton's third law the exhaust pushes the rocket forward. No external medium is needed; momentum of the rocket + exhaust system is conserved.
\item \textbf{Part C (e) 2 marks:} $x_{\mathrm{cm}} = \dfrac{8}{6} \approx 1.33\ \mathrm{m}$; $y_{\mathrm{cm}} = \dfrac{9}{6} = 1.5\ \mathrm{m}$.
\item \textbf{Part C (f) 1 mark:} $\vec{F}_{\mathrm{ext}}=\vec{0} \Rightarrow \dfrac{d\vec{P}}{dt}=\vec{0} \Rightarrow \vec{P}=\text{constant}$; since $\vec{P}=M\vec{v}_{\mathrm{cm}}$, $\vec{v}_{\mathrm{cm}}$ is constant.
\end{itemize}

\textbf{Detailed solution:}

\textbf{Part A.}
\begin{enumerate}
\item The shell is initially at rest, so the total momentum before the explosion is $0$. Take the direction of the $5\ \mathrm{kg}$ fragment as positive and let $v$ be the velocity of the $7\ \mathrm{kg}$ fragment. Conservation of momentum:
\[ 5 \times 84 + 7v = 0 \;\Rightarrow\; v = -\frac{420}{7} = -60\ \mathrm{m\,s^{-1}}. \]
The $7\ \mathrm{kg}$ fragment moves at $60\ \mathrm{m\,s^{-1}}$ in the direction opposite to the $5\ \mathrm{kg}$ fragment.

\item Kinetic energy after the explosion:
\[ K = \frac{1}{2}(5)(84^2) + \frac{1}{2}(7)(60^2) = \frac{1}{2}(35\,280 + 25\,200) = 30\,240\ \mathrm{J}. \]
Since the initial kinetic energy was zero, the energy released by the explosion is $30\,240\ \mathrm{J}$. It comes from the chemical potential energy stored in the explosive, converted into kinetic energy of the fragments.
\end{enumerate}

\textbf{Part B.}
\begin{enumerate}
\item[(c)] The thrust is the rate of change of momentum of the exhaust gas:
\[ F = v_{\mathrm{rel}} \times \frac{dm}{dt} = 2000 \times 150 = 300\,000\ \mathrm{N} = 3.0 \times 10^5\ \mathrm{N}. \]

\item[(d)] The rocket accelerates by ejecting mass (exhaust gas) backwards at high speed relative to itself. The rocket and the exhaust gas form an isolated system, so their total momentum is conserved: the forward momentum gained by the rocket is equal in magnitude to the backward momentum carried away by the exhaust. Equivalently, the rocket exerts a backward force on the gas, and by Newton's third law the gas exerts a forward reaction force (the thrust) on the rocket. No external medium (such as air) is required, which is why rockets work in the vacuum of space.
\end{enumerate}

\textbf{Part C.}
\begin{enumerate}
\item[(e)] Total mass $M = 1+2+3 = 6\ \mathrm{kg}$. The coordinates of the centre of mass are
\[ x_{\mathrm{cm}} = \frac{1(0) + 2(4) + 3(0)}{6} = \frac{8}{6} = \frac{4}{3} \approx 1.33\ \mathrm{m}, \]
\[ y_{\mathrm{cm}} = \frac{1(0) + 2(0) + 3(3)}{6} = \frac{9}{6} = 1.5\ \mathrm{m}. \]
The centre of mass is at the point $\left(\frac{4}{3},\; 1.5\right)\ \mathrm{m}$.

\item[(f)] For an isolated system the net external force is zero: $\vec{F}_{\mathrm{ext}} = \vec{0}$. Newton's second law for a system of particles gives $\vec{F}_{\mathrm{ext}} = \dfrac{d\vec{P}}{dt}$, where $\vec{P}$ is the total momentum. Hence $\dfrac{d\vec{P}}{dt} = \vec{0}$, so $\vec{P}$ is constant. But $\vec{P} = M\vec{v}_{\mathrm{cm}}$ with constant total mass $M$, therefore $\vec{v}_{\mathrm{cm}}$ is constant: the centre of mass moves with constant velocity (or remains at rest if it was initially at rest).
\end{enumerate}
\end{corrige}

\newpage

\coursec{Summary sheet}

\begin{aretenir}[Bilan du chapitre]
\begin{popfigure}
\begin{center}
\begin{tikzpicture}[
  mindmap,
  level 1/.style={sibling angle=90, level distance=3.5cm},
  level 2/.style={sibling angle=45, level distance=2.5cm},
  every node/.style={concept, font=\small, execute at begin node=\hskip0pt},
  root concept/.append style={font=\bfseries\large, concept color=popblue},
  concept color=popblue!70
]
\node[root concept] {Chapitre}
  child[concept color=popgreen] {node[concept, align=center]{Notions\\clés}
    child {node[concept] {Définitions}}
    child {node[concept] {Théorèmes}}
    child {node[concept] {Formules}}}
  child[concept color=poporange] {node[concept] {Méthodes}
    child {node[concept] {Résolution}}
    child {node[concept] {Démonstration}}
    child {node[concept] {Application}}}
  child[concept color=poppurple] {node[concept] {Pièges}
    child {node[concept] {Erreurs fréquentes}}
    child {node[concept] {Cas particuliers}}
    child {node[concept] {Conditions}}}
  child[concept color=poppink] {node[concept] {Exercices}
    child {node[concept] {Types classiques}}
    child {node[concept] {Sujets d'examen}}
    child {node[concept] {Approfondissement}}};
\end{tikzpicture}
\end{center}
\legende{Carte mentale récapitulative du chapitre}
\end{popfigure}
\end{aretenir}

\findecours

\end{document}
